
\documentclass[12pt,letterpaper,oneside]{book}

\usepackage[margin=1in]{geometry}
\usepackage{setspace}
\usepackage{fancyhdr}
\usepackage{titlesec}
\usepackage{tocloft}
\usepackage{etoolbox}

\usepackage[T1]{fontenc}
\usepackage[utf8]{inputenc}
\usepackage{microtype}
\usepackage{libertine}
\usepackage{inconsolata}  
\usepackage{siunitx}

\usepackage{amsmath}
\usepackage{amsthm}
\usepackage{mathtools}
\usepackage[libertine]{newtxmath}

\usepackage{graphicx}
\usepackage{float}
\usepackage{booktabs}
\usepackage{adjustbox}
\usepackage{threeparttable}
\usepackage{subcaption}
\usepackage{rotating}
\usepackage{longtable}
\usepackage{multirow}
\usepackage{array}
\usepackage{tabularx}

\usepackage{tikz}
\usepackage{pgfplots}
\pgfplotsset{compat=1.16}
\usetikzlibrary{patterns,arrows,shapes,positioning}

\usepackage{natbib}

\usepackage[colorlinks=true,linkcolor=blue,citecolor=blue,urlcolor=blue,bookmarks=true]{hyperref}
\usepackage{bookmark}

\usepackage[hang,flushmargin,bottom]{footmisc}
\usepackage{appendix}
\usepackage{enumerate}
\usepackage{enumitem}
\usepackage{lipsum}
\usepackage[english]{babel}
\usepackage{ifplatform}

% Cross-platform base path configuration
\ifwindows
  \newcommand{\basepath}{g:/My Drive/Grad School}
\else
  \newcommand{\basepath}{basepath}
\fi

\titleformat{\chapter}[display]
{\normalfont\huge\bfseries}
{\chaptertitlename\ \thechapter}
{20pt}
{\Huge}

\titleformat{\section}
{\normalfont\Large\bfseries}
{\thesection}
{1em}
{}

\titleformat{\subsection}
{\normalfont\large\bfseries}
{\thesubsection}
{1em}
{}

\renewcommand{\cftchapfont}{\bfseries}
\renewcommand{\cftchappagefont}{\bfseries}
\renewcommand{\cftchapleader}{\cftdotfill{\cftdotsep}}
\setlength{\cftbeforechapskip}{0.5em}

\pagestyle{fancy}
\fancyhf{}
\fancyhead[R]{\thepage}
\fancyhead[L]{\leftmark}
\renewcommand{\headrulewidth}{0pt}
\renewcommand{\chaptermark}[1]{\markboth{\MakeUppercase{#1}}{}}

\newcommand{\mychapter}[1]{
    \chapter{#1}
    \doublespacing
}

\widowpenalty=10000
\clubpenalty=10000

\setlength{\parindent}{0.5in}

\usepackage{caption}
\captionsetup{font=small}


\begin{document}


\frontmatter
\pagestyle{plain}

\begin{titlepage}
\begin{center}
\vspace*{1in}

{\LARGE\bfseries THREE ESSAYS IN POLITICAL ECONOMY:}\\[0.5em]
{\LARGE\bfseries INSTITUTIONS, INFORMATION, AND CIVIC ENGAGEMENT}\\[2in]

{\Large A Dissertation}\\[0.5em]
{\Large Presented to the Faculty of}\\[0.5em]
{\Large The Department of Political Science at}\\[0.5em]
{\Large the University of California, Los Angeles}\\[0.5em]
{\Large in Partial Fulfillment of the Requirements}\\[0.5em]
{\Large for the Degree of}\\[0.5em]
{\Large Doctor of Philosophy}\\[2in]

{\Large by}\\[0.5em]
{\Large\bfseries Igor Geyn}\\[2in]

{\Large 2025}
\end{center}
\end{titlepage}

\clearpage
\thispagestyle{empty}
\vspace*{\fill}
\begin{center}
Copyright \copyright\ 2025 by Igor Geyn\\
All rights reserved.
\end{center}
\vspace*{\fill}
\clearpage

\chapter*{Abstract}
\addcontentsline{toc}{chapter}{Abstract}

This dissertation presents three empirical studies examining the intersection of political institutions, information, and civic engagement in the United States. Together, these essays contribute to our understanding of how institutional design shapes political behavior and democratic participation, while also demonstrating some advantages and some nuances of using contemporary causal inference approaches in observational settings.

The first essay examines whether federal disaster relief favors congressional districts represented by the president's copartisans. Using a regression discontinuity design that exploits close congressional elections from 2001-2020, I compare FEMA funding in districts that narrowly elect a member of the president's party to those that narrowly elect an opposition member. I find that copartisan districts receive approximately 45\% less federal disaster relief than opposition districts (p=0.062), with this negative effect consistent across the Bush, Obama, and Trump administrations (estimates ranging from -45\% to -57\%). Extensive mechanism tests point toward a compositional explanation: disaster-prone areas with particular demographic and infrastructure characteristics systematically differ in their political representation. The effect concentrates in low-visibility infrastructure spending, exhibits a striking gradient across district racial composition, and shows no response to electoral incentives---patterns inconsistent with intentional political targeting but consistent with underlying geographic and compositional differences in disaster vulnerability and preparedness.

The second essay investigates whether ballot measures enhance civic engagement by increasing voters' knowledge and political interest. Using comprehensive data on 1,547 ballot measures across all U.S. states from 2006-2020 merged with 480,383 individual survey responses, I demonstrate how methodological choices fundamentally determine conclusions about this question. While traditional two-way fixed effects models suggest positive engagement effects, modern difference-in-differences estimators designed for staggered treatment adoption reveal essentially null effects. The analysis provides both a substantive contribution—challenging optimistic accounts of direct democracy's civic benefits—and a methodological cautionary tale about the implications of estimator selection in applied, observational data settings.

The third essay---co-authored with Daniel Firoozi and published in the journal American Economic Journal: Economic Policy---estimates the impact of California's Cal Grant program, America's largest tuition-free 4-year college program, on political participation. Using a regression discontinuity design based on strict income eligibility thresholds and 16.4 million financial aid applications, my co-author and I find that Cal Grant receipt increases voter turnout by approximately 10 percentage points. The effect operates almost entirely through increased registration and turnout among left-leaning voters, with the 2.6 million grants awarded over the 2010s inducing an additional 259,000 votes in the 2020 general election. These findings highlight how education subsidies generate important civic externalities beyond their effects on human capital and earnings.

\chapter*{Acknowledgments}
\addcontentsline{toc}{chapter}{Acknowledgments}

% TODO: Write acknowledgments before final submission
% Typical acknowledgments include: committee members, advisors, colleagues,
% funding sources (grants, fellowships), family, and friends.

\begingroup
\singlespacing
\tableofcontents
\endgroup

\clearpage
\phantomsection
\addcontentsline{toc}{chapter}{List of Tables}
\begingroup
\singlespacing
\listoftables
\endgroup

\clearpage
\phantomsection
\addcontentsline{toc}{chapter}{List of Figures}
\begingroup
\singlespacing
\listoffigures
\endgroup


\mainmatter
\pagestyle{fancy}
\doublespacing


\chapter*{Introduction}
\addcontentsline{toc}{chapter}{Introduction}
\markboth{INTRODUCTION}{}

Political institutions have the potential to shape the distribution of resources, flow of information, and patterns of civic engagement. This dissertation examines these fundamental relationships through three empirical studies that leverage quasi-experimental variation to identify causal effects. Each essay addresses a distinct question about how institutions influence political behavior while contributing to broader debates about democratic governance, policy design, and citizen participation.

The three essays are united by several common themes. First, they all exploit institutional rules and thresholds to create quasi-random variation in treatment assignment, enabling credible causal inference under some circumstances (i.e. with sufficient data/a study that is well-powered to detect an effect). Second, they examine outcomes central to democratic functioning: the fair distribution of public resources, citizen knowledge and engagement, and political participation. Third, they demonstrate how methodological choices can  affect substantive conclusions, highlighting the importance of appropriate econometric techniques for policy-relevant research and suggesting some best practices to use in applied settings.

\section*{Organization}

The dissertation proceeds as follows. Chapters 1-3 present the three main empirical essays, each written as a standalone paper with its own introduction, literature review, empirical strategy, results, and conclusion. Following the main chapters, I provide supplementary appendices with additional robustness tests, data descriptions, and technical details. The dissertation concludes with a brief epilogue synthesizing lessons across the three studies and suggesting directions for future research.


\mychapter{Copartisan Districts and Federal Disaster Relief: Evidence from a Regression Discontinuity Design}

\section{Introduction}

Does the president's party affiliation shape who receives federal disaster relief? The question is fundamental to democratic governance. Unlike many federal programs governed by rigid formulas, disaster relief grants substantial discretion to political actors: the president personally approves disaster declarations, and FEMA officials determine funding allocations within broad statutory guidelines. Between 2001 and 2020, this discretionary system distributed over \$150 billion, affecting virtually every congressional district in the nation. If elected officials can systematically direct these resources toward political allies, democracy's promise of equal treatment under law becomes hollow---particularly for programs addressing urgent humanitarian needs under conditions of heightened public attention.

A substantial literature documents partisan patterns in federal spending. \cite{berry2010presidents} demonstrate presidential particularism in discretionary grant programs. \cite{kriner2015particularistic} find that electorally important states receive more federal dollars. \cite{larcinese2006allocating} show that states with aligned governors receive more federal funds. These findings collectively suggest that political considerations permeate federal resource allocation, raising the question of whether disaster relief---with its combination of executive discretion and humanitarian urgency---follows the same pattern.

Yet strong theoretical reasons also predict constraints on partisan manipulation. Professional bureaucracies operate under institutional rules that limit discretion \citep{north1990institutions, ostrom2005understanding}. Career civil servants may resist political pressure through expertise and standardized procedures \citep{carpenter2001forging}. Public scrutiny of disaster response creates accountability mechanisms \citep{healy2009myopic, gasper2011make}, and the urgent timeline of disaster relief may preclude careful political calculation. Which forces dominate---opportunities for manipulation or institutional constraints---is ultimately an empirical question.

I address this question using a regression discontinuity design that exploits close congressional elections. In districts where the president's party candidate barely wins, the representative is copartisan with the president; where the candidate barely loses, the representative is in the opposition. Comparing FEMA funding across these nearly identical districts identifies whether political alignment affects disaster relief allocation.

The analysis reveals a surprising finding: copartisan districts receive approximately 45\% \emph{less} disaster relief than opposition districts (p=0.062). This negative effect is consistent across the Bush, Obama, and Trump administrations, with estimates ranging from $-45\%$ to $-57\%$. The consistency across administrations of both parties suggests that this pattern reflects systematic factors rather than partisan manipulation by any single administration.

This finding runs counter to standard expectations of presidential favoritism toward copartisan districts. To understand what drives this pattern, I conduct an extensive series of mechanism tests. First, I find that the effect is \emph{not} stronger in election years, which undermines the hypothesis that presidents strategically target opposition districts for electoral gain. Second, I examine heterogeneity across spending categories and find that the negative effect is concentrated in low-visibility infrastructure spending ($-41\%$, p=0.041), while highly visible emergency spending shows no significant effect. Third, I find striking heterogeneity by district demographics: the effect is concentrated in predominantly white, affluent, highly-educated districts, with a clear monotonic gradient across the racial diversity distribution. Districts in the whitest quintile show effects of $-81\%$ (p=0.058), attenuating steadily to a positive (though imprecise) $+168\%$ in the most diverse quintile. Fourth, using causal forest methods \citep{wager2018estimation}, I confirm that treatment effects are homogeneous across observable characteristics---inconsistent with strategic manipulation, which would produce heterogeneity by electoral incentives.

Together, these mechanism tests point toward a compositional explanation. Areas with particular demographic and infrastructure characteristics---which correlate with political representation---differ systematically in their disaster experiences and relief needs. The ``effect'' of copartisan status reflects these underlying compositional differences rather than intentional political targeting.

These findings make several contributions to the literature on distributive politics and disaster policy. First, they challenge assumptions about the pervasiveness of partisan influence over federal spending. While prior work documents political manipulation in discretionary grants, disaster relief appears to resist such influence---indeed, the pattern runs opposite to standard predictions. Second, they illuminate how political geography and disaster vulnerability intersect in previously unexplored ways: predominantly white, affluent districts with lower diversity show the strongest negative effects, revealing connections between community characteristics and disaster relief patterns that have not been documented. Third, they demonstrate the critical importance of mechanism tests in interpreting regression discontinuity estimates. Without the heterogeneity analyses, one might conclude that presidents intentionally disadvantage copartisan districts; the mechanism tests reveal a more nuanced story about compositional differences between areas represented by copartisan versus opposition legislators. Fourth, I provide historical context by replicating methods on 1988--2004 data from \cite{healy2009myopic}, finding suggestive evidence of the opposite pattern in earlier periods---consistent with evolving political geography and further supporting the compositional interpretation.

The paper proceeds as follows. Section 2 reviews theoretical expectations and prior research, including hypotheses for why copartisan districts might receive \emph{less} relief. Section 3 describes data sources and the regression discontinuity design. Section 4 presents the main results showing negative effects of political alignment. Section 5 examines heterogeneity across contexts and conducts mechanism tests, including demographic heterogeneity, spending categories, and machine learning analysis. Section 6 presents additional evidence from historical data, governor alignment, and declaration denials. Section 7 discusses the compositional interpretation and its implications. Section 8 concludes.

\section{Theoretical Background}\label{sec:theory}

\subsection{Political Manipulation of Federal Spending}

A substantial literature documents partisan patterns in federal resource allocation. At the broadest level, \cite{levitt1997impact} demonstrates that federal spending responds to electoral considerations rather than strict programmatic criteria. This general finding raises questions about whether specific programs, including disaster relief, are similarly subject to political influence.

Presidential particularism represents one prominent mechanism. \cite{berry2010presidents} demonstrate that districts represented by the president's copartisans receive systematically more federal spending. This effect operates through executive branch agencies responsive to presidential priorities, suggesting that FEMA could potentially direct disaster relief toward politically favored districts. \cite{kriner2015particularistic} extend this logic to show that electorally important states receive more federal dollars, particularly when controlled by the president's party.

Congressional influence provides another potential channel. \cite{levitt1997impact} finds that federal spending responds to legislators' committee assignments and electoral vulnerability. \cite{berry2010presidents} show that politically powerful legislators secure more defense spending for their districts. If congressional representatives can influence FEMA allocations through oversight, appropriations, or informal pressure, copartisan districts might receive favorable treatment.

State-level political alignment may also matter. \cite{larcinese2006allocating} demonstrate that states with governors from the president's party receive more federal funds. Since disaster declarations require state requests and federal approval, this alignment could affect both the likelihood of declaration and subsequent funding amounts.

\subsection{Disaster Relief as a Test Case}

Disaster relief offers a theoretically compelling setting to examine political manipulation for several reasons. First, unlike formula-based programs, disaster declarations require presidential approval, providing executive discretion that creates opportunities for political considerations to influence decisions \citep{reeves2011political, gasper2011make}. Second, within statutory limits, FEMA could provide more generous interpretations of eligible costs or faster processing for copartisan districts, as the complexity of damage assessment and cost estimation creates space for favorable treatment \citep{mccarthy2011stafford}. Third, legislators can advocate for constituents through casework, oversight hearings, and informal communications, and copartisan legislators may enjoy enhanced access and influence \citep{ritchie2018back, lowande2018descriptive}. Fourth, political appointees could direct agency resources, staffing, and attention to benefit copartisan districts, as the mix of political and career personnel in FEMA creates opportunities for such influence \citep{lewis2008politics}.

These opportunities for manipulation must be weighed against institutional constraints that could limit political influence. Career civil servants implement programs according to detailed regulations governing eligibility and benefits. Requirements for documentation create paper trails, and potential legal challenges to arbitrary decisions provide accountability mechanisms. Understanding which forces dominate---manipulation or constraint---is the empirical question this paper addresses.

\subsection{Why Might Copartisan Districts Receive Less?}

While the standard expectation is that copartisan districts would receive \emph{more} resources, several theoretical perspectives suggest the opposite pattern could emerge.

\textbf{Opposition Targeting for Persuasion}: Models of distributive politics distinguish between mobilizing core supporters and persuading swing voters \citep{cox1986electoral, dixit1996determinants}. \cite{lindbeck1987balanced} argue that politicians may target swing voters who are responsive to material benefits, rather than core supporters whose votes are already secure. Applied to disaster relief, presidents might direct resources toward opposition districts to cultivate goodwill among persuadable voters. This strategy would be particularly attractive if disaster relief is highly visible and easily attributed to the president.

\textbf{Scrutiny Avoidance}: Following Hurricane Katrina, accusations of political bias in disaster response received intense media coverage. Administrations may subsequently overcompensate by directing resources away from copartisan districts to avoid similar accusations \citep{besley2006handcuffs}. This ``leaning against the wind'' behavior would produce negative effects of copartisan alignment even absent any intentional targeting.

\textbf{Compositional Differences}: A third possibility is that the relationship between political alignment and disaster relief reflects compositional differences rather than intentional manipulation. If disaster-prone areas with inadequate mitigation infrastructure systematically differ in their political representation, regression discontinuity estimates might capture these underlying differences rather than causal effects of political alignment on administrative behavior.

Consider the following scenario: areas with low prior investment in hazard mitigation face larger disaster losses when disasters strike. These areas may also differ economically and demographically in ways that correlate with political representation. If, during certain periods, such areas happen to be represented by opposition legislators at higher rates, we would observe lower average relief in copartisan districts---not because of political targeting, but because of which areas experience which disasters.

This compositional explanation generates distinct empirical predictions: the ``effect'' should vary with factors that capture underlying vulnerability (such as prior mitigation investment or demographic characteristics), should not respond to electoral incentives (such as proximity to elections), and might attenuate over time as political geography shifts.

\subsection{Discriminating Among Hypotheses}

The three hypotheses generate different predictions for heterogeneity patterns. If presidents strategically direct resources to opposition districts for electoral gain (the opposition targeting hypothesis), effects should be stronger in election years and in electorally competitive districts, and should vary by administration based on political strategy. If administrations overcompensate to avoid favoritism accusations (the scrutiny avoidance hypothesis), effects might be stronger after high-profile disasters or media criticism, but should not systematically vary with disaster characteristics or district demographics. If effects reflect underlying differences between disaster-affected areas (the compositional differences hypothesis), they should vary with factors capturing vulnerability---such as prior mitigation investment and demographic composition---but not respond to electoral incentives. Under the compositional hypothesis, effects should be homogeneous across electoral contexts but heterogeneous across compositional factors, and might attenuate over time as political geography shifts.

My empirical strategy tests these predictions through heterogeneity analyses examining electoral cycles, presidential terms, prior mitigation investment, demographic characteristics, and spending categories. I also employ causal forest methods \citep{wager2018estimation, athey2016recursive} to systematically test for treatment effect heterogeneity across all observable covariates.

\section{Data and Research Design}\label{sec:data}

\subsection{Data Sources}

I compile data from several administrative sources to create a comprehensive panel of disaster spending and political alignment at the congressional district level from 2001-2020.

\textbf{FEMA Disaster Spending.} The OpenFEMA API provides project-level data on all federal disaster obligations. For each project, I observe the obligated amount, obligation date, disaster number, damage category, county, and program type. I aggregate these data to district-year observations, carefully accounting for redistricting. When counties split across districts, I allocate spending using population weights from the Missouri Census Data Center's GEOCORR engine. This granular data allows me to examine both total spending and spending by program type, distinguishing between highly visible categories (debris removal, emergency protective measures) and lower-visibility infrastructure spending (roads, bridges, public buildings, utilities).

\textbf{Congressional Election Results.} I obtain House election returns from the MIT Election Data and Science Lab, which provides precinct-level results for all House races from 2000-2020. I calculate two-party vote shares and identify close elections where outcomes were decided by narrow margins.

\textbf{Presidential Party.} I code presidential party for each year: Republican during Bush (2001-2008) and Trump (2017-2020), Democratic during Obama (2009-2016).

\textbf{NOAA Damage Data.} To control for objective disaster severity, I incorporate storm-level damage estimates from the National Oceanic and Atmospheric Administration's Storm Events Database. NOAA provides independent assessments of property and crop damage that are determined separately from FEMA's funding decisions.

\textbf{Hazard Mitigation Assistance.} To test the compositional hypothesis, I compile data on cumulative Hazard Mitigation Assistance (HMA) grants by congressional district. HMA programs fund projects that reduce future disaster vulnerability, such as property acquisitions, structural retrofitting, and drainage improvements. Districts with higher prior HMA investment have greater disaster preparedness infrastructure.

\textbf{Census Demographics.} I merge American Community Survey 5-year estimates to characterize district demographics, including median household income, poverty rates, educational attainment, and racial composition. I use 2012 ACS data for 2001--2012 observations and 2019 ACS data for 2013--2020 observations to avoid post-treatment contamination.

\textbf{Governor Data.} I incorporate state-level data on gubernatorial party from \cite{klarner2013state} to test whether state-federal partisan alignment moderates the MC-president copartisan effect.

\textbf{Historical Data.} For historical comparison, I use county-level disaster relief data from \cite{healy2009myopic} covering 1988--2004, which I aggregate to congressional districts using historical GEOCORR crosswalks.

\subsection{Research Design}

I employ a regression discontinuity (RD) design exploiting close congressional elections. The key insight is that in elections decided by narrow margins, the winning party is essentially randomly determined \citep{lee2008randomized}. Districts where the president's party barely wins versus barely loses are comparable on observable and unobservable characteristics, allowing causal identification of political alignment effects.

\textbf{Running Variable Construction.} A critical methodological issue is constructing the running variable to capture copartisan status regardless of which party holds the presidency. I define:
\begin{align}
\text{rv1} &= \text{dem\_share} - 0.5 \\
\text{rv1\_copartisan} &= \text{rv1} \times \text{prez\_dem\_indicator}
\end{align}
where $\text{prez\_dem\_indicator} = +1$ under Democratic presidents and $-1$ under Republican presidents.

This transformation ensures that positive values of rv1\_copartisan always indicate copartisan representation (regardless of which party holds the presidency), while negative values indicate opposition representation. Without this transformation, pooling data across Democratic and Republican administrations would conflate copartisan and opposition effects, potentially producing spurious null results.

\textbf{Estimation.} I estimate local linear regressions using the rdrobust package \citep{calonico2014robust}, which implements bias-corrected estimation with robust inference. All reported standard errors, p-values, and confidence intervals use the robust inference procedure (the third row of rdrobust output), which accounts for bias correction in the standard error calculation. Standard errors are clustered at the state level throughout.

The main specification is:
\begin{equation}
\log(\text{FEMA}_i) = \alpha + \tau \cdot \mathbf{1}[\text{rv1\_copartisan}_i > 0] + f(\text{rv1\_copartisan}_i) + \varepsilon_i
\end{equation}
where $\tau$ captures the discontinuity at the threshold---the causal effect of having a copartisan representative on log FEMA obligations.

\subsection{Graphical Evidence}

Figure \ref{fig:main_rd} presents the main regression discontinuity plot. Each point represents binned averages of log FEMA obligations, with the running variable (copartisan victory margin) on the horizontal axis. The vertical dashed line marks the threshold where representation switches from opposition (left) to copartisan (right).

\begin{figure}[H]
\centering
\includegraphics[width=0.9\textwidth]{\basepath/Research/events_extreme_candidate_entry/diss/code/prod/output/plots/main/figure1_main_rd_IMPROVED.pdf}
\caption{Political Alignment and FEMA Disaster Relief}
\label{fig:main_rd}
\caption*{\footnotesize Notes: Regression discontinuity plot showing log FEMA Public Assistance obligations (2001--2020) against the copartisan victory margin. Positive margins indicate the district's representative shares the president's party; negative margins indicate opposition representation. Solid lines plot local linear fits estimated separately on each side of the threshold using the MSE-optimal bandwidth; shaded bands denote 95\% robust confidence intervals. The discontinuous drop at the threshold indicates copartisan districts receive less disaster relief.}
\end{figure}

The figure reveals a notable discontinuity at the threshold: FEMA obligations drop when moving from opposition to copartisan representation. This visual evidence suggests copartisan districts receive \emph{less} disaster relief than opposition districts, contrary to standard expectations of political favoritism.

\subsection{Regression Estimates}

Table \ref{tab:main_results} presents formal regression discontinuity estimates. Column 1 shows the baseline specification: copartisan districts receive 45.5\% less disaster relief than opposition districts (p = 0.062). This marginally significant negative estimate challenges the standard hypothesis of presidential favoritism toward copartisan districts.

\input{\basepath/Research/events_extreme_candidate_entry/diss/code/prod/output/tables/main/table1_main_results_CLEAN.tex}
% Label provided by input file

The remaining columns demonstrate robustness. Column 2 includes state and year fixed effects, strengthening the estimate to $-45.7\%$ (p = 0.018). Columns 3-5 incorporate NOAA damage data. Column 4 controls for log damage, with the copartisan effect remaining stable at $-45.8\%$ (p = 0.063). This indicates the negative finding is not explained by differences in objective disaster severity between copartisan and opposition districts.

Column 5 examines the FEMA-to-damage ratio as an alternative outcome, addressing whether copartisan districts receive different treatment \emph{conditional on damage}. The coefficient is $-38.8\%$ but with larger standard errors (p = 0.286), providing suggestive but imprecise evidence of differential treatment per dollar of damage.

These results are substantively important. Given average FEMA obligations of \$37 million per district-year, a 45\% reduction implies copartisan districts receive roughly \$17 million less than comparable opposition districts. While the baseline p-value of 0.062 exceeds conventional significance thresholds, the consistency across specifications and the precision of estimates with fixed effects (p = 0.018) warrant serious attention.

\subsection{Validity of the Design}

The regression discontinuity design requires balance on predetermined characteristics near the threshold. Table \ref{tab:balance_noaa} presents balance tests for NOAA damage measures. I find no significant discontinuities in total damage, property damage, flood damage, or disaster counts. The joint F-test fails to reject balance (F = 0.75, p = 0.65).

\input{\basepath/Research/events_extreme_candidate_entry/diss/code/prod/output/tables/main/table3_balance_test_noaa_CORRECTED.tex}
% Label provided by input file

Figure \ref{fig:damage_balance} visualizes the balance test for total damage. The smooth relationship through the threshold confirms that close elections generate comparable treatment and control groups with respect to disaster severity.

\begin{figure}[H]
\centering
\includegraphics[width=0.9\textwidth]{\basepath/Research/events_extreme_candidate_entry/diss/code/prod/output/plots/main/figure2_damage_balance_IMPROVED.pdf}
\caption{Balance Test: Objective Damage Levels at RD Cutoff}
\label{fig:damage_balance}
\caption*{\footnotesize Notes: Regression discontinuity diagnostic for total NOAA-reported disaster damage (log dollars) aggregated to district-years, 2001--2020. Local linear fits use the same kernel and bandwidth as the main specification; shaded bands show 95\% confidence intervals. The smooth relationship through the threshold indicates balance on objective disaster severity.}
\end{figure}

Additional validity tests support the design. McCrary's density test finds no evidence of manipulation in the running variable; for the Obama subsample, the test yields p = 0.150, suggesting close elections are not systematically won by one party through manipulation. Testing for discontinuities at placebo thresholds yields mixed results: while most placebo tests are insignificant, the test at $c = -0.10$ shows a significant positive effect ($+94.8\%$, p = 0.019). Crucially, the true threshold ($c = 0$) shows a distinct negative effect ($-49.2\%$) while the significant placebo shows a positive effect, suggesting these capture different phenomena rather than undermining the main finding. The negative effect persists across all bandwidths tested, ranging from $-38\%$ at double bandwidth to $-52\%$ at half bandwidth, indicating that results are not sensitive to the choice of estimation window. Finally, excluding observations very close to the threshold---a ``donut hole'' specification that addresses concerns about measurement error near the cutoff---does not flip the sign of estimates, though precision decreases as expected with the smaller sample.

The consistency of findings across these multiple validity checks strengthens confidence in the causal interpretation.

\section{Heterogeneity and Mechanism Tests}\label{sec:heterogeneity}

The finding that copartisan districts receive less disaster relief is surprising and requires explanation. In this section, I conduct heterogeneity analyses designed to discriminate among competing hypotheses: opposition targeting for electoral gain, scrutiny avoidance, and compositional differences.

\subsection{Effects by Presidential Administration}

If the negative effect reflects intentional targeting by a particular administration, we would expect variation across presidencies. Table \ref{tab:era_heterogeneity} examines whether effects vary across the Bush, Obama, and Trump administrations.

\input{\basepath/Research/events_extreme_candidate_entry/diss/code/prod/output/tables/main/table4_era_heterogeneity_CORRECTED.tex}
% Label provided by input file

The results reveal consistency across administrations. The estimates are: Bush ($-45.2\%$, p = 0.243), Obama ($-56.9\%$, p = 0.057), and Trump ($-55.6\%$, p = 0.117). All three point estimates are negative and of similar magnitude, despite spanning administrations of both parties with different political contexts and leadership styles.

This consistency is important for interpretation. If the pattern reflected intentional partisan manipulation, we might expect it to vary with presidential ideology or political strategy. Instead, the uniformity suggests systematic factors that transcend individual administrations.

This finding contrasts with much of the prior literature on disaster relief allocation. \cite{garrett2003political} find that states politically important to the president have higher rates of disaster declaration, with states represented on FEMA oversight committees receiving substantially more funding; they estimate that nearly half of all disaster relief is politically motivated rather than need-based. \cite{reeves2011political} demonstrates that electorally competitive states receive approximately twice as many presidential disaster declarations as uncompetitive states, with each declaration yielding roughly a one-point increase in the president's statewide vote share. More recently, \cite{kriner2015disasters} find that President Obama favored both swing and core states in the 2012 election year, and \cite{schneider2022disastrous} document that political bias in disaster declarations is concentrated in medium-intensity hurricanes, where presidential discretion is greatest.

The common thread in this literature is that presidents use disaster relief to benefit their political allies or electorally valuable constituencies. The negative copartisan effect documented here---where the president's co-partisans receive \emph{less} relief---represents a departure from these findings. However, the consistency across three administrations of both parties suggests this pattern reflects structural features of disaster-prone areas rather than any particular president's strategy. If Bush, Obama, and Trump all oversaw systems that delivered less relief to copartisan districts, the explanation likely lies in the characteristics of the districts themselves rather than in executive decision-making.

% Figure \ref{fig:comparison_president} visualizes this consistency, showing mean FEMA obligations for copartisan and opposition districts separately by presidential administration. Under all three presidents, opposition districts receive higher mean obligations than copartisan districts, with the pattern most pronounced during the Obama years. The consistency of this pattern across administrations of both parties provides strong evidence against explanations based on partisan manipulation by any single administration.

% \begin{figure}[H]
% \centering
% \includegraphics[width=0.85\textwidth]{\basepath/Research/events_extreme_candidate_entry/diss/code/prod/output/plots/descriptives/figure_comparison_by_president.pdf}
% \caption{Mean FEMA Obligations: Copartisan vs.\ Opposition Districts by President}
% \label{fig:comparison_president}
% \caption*{\footnotesize Notes: Bar chart shows mean FEMA Public Assistance obligations (in millions of dollars) for copartisan districts (where the representative shares the president's party) and opposition districts, separately by presidential administration. Error bars show 95\% confidence intervals. Opposition districts receive higher mean obligations under all three presidents, consistent with the main RD finding of a negative copartisan effect.}
% \end{figure}

\subsection{Electoral Cycle Analysis}

A key prediction of the opposition targeting hypothesis is that effects should be stronger in election years, when electoral incentives are most salient. Table \ref{tab:electoral_cycle} tests whether political effects vary by position in the electoral cycle.

\input{\basepath/Research/events_extreme_candidate_entry/diss/code/prod/output/tables/main/table_electoral_cycle.tex}
% Label provided by input file

The results provide no support for electoral targeting. Non-election years show a large negative effect ($-51.0\%$, p = 0.038), while presidential election years show essentially no effect ($-11.8\%$, p = 0.967). If anything, the pattern is \emph{opposite} to what the opposition targeting hypothesis predicts: the negative effect is concentrated in years \emph{without} presidential elections.

This finding undermines the hypothesis that presidents strategically direct disaster relief to opposition districts to win swing voters. Electoral incentives do not appear to drive the observed pattern.

\subsection{Obama Term 1 vs. Term 2}

To further probe the dynamics of the effect, I examine heterogeneity across Obama's two terms. If the pattern reflects compositional differences that shift over time, we might expect attenuation as political geography changes.

\input{\basepath/Research/events_extreme_candidate_entry/diss/code/prod/output/tables/main/table_mechanisms.tex}

The results reveal striking heterogeneity. During Obama's first term (2009-2012), copartisan districts received 61.1\% less disaster relief, with a relatively precise estimate (SE = 0.554). During his second term (2013-2016), the effect attenuated substantially to $-33.6\%$ with a much larger standard error (0.896), indicating considerable uncertainty around this estimate.

This temporal pattern is inconsistent with intentional targeting, which should persist throughout an administration. Instead, it suggests that the composition of disaster-affected areas changed over the Obama presidency in ways that altered the relationship between political alignment and relief amounts.

\subsection{Prior Mitigation Investment}

The compositional hypothesis predicts that effects should vary with factors capturing underlying disaster vulnerability. I test this by examining heterogeneity based on prior Hazard Mitigation Assistance (HMA) investment---a measure of disaster preparedness infrastructure.

Districts with \emph{low} prior HMA investment show a large negative effect: copartisan districts receive 75.8\% less relief (SE = 0.701). In contrast, districts with \emph{high} prior HMA investment show a much smaller effect of $-28.3\%$ (SE = 0.676). The estimates have similar precision, making the difference in magnitudes substantively meaningful.

This pattern strongly supports the compositional explanation. Areas with inadequate disaster preparedness infrastructure---which face larger losses when disasters strike---show the strongest ``effect'' of copartisan status. This is precisely what we would expect if the relationship reflects underlying differences between disaster-affected areas rather than intentional targeting.

\subsection{Effects by Spending Category}

Different types of disaster spending vary in their visibility to voters and the discretion available to administrators. If political manipulation occurs, it might concentrate in spending categories where administrators have more latitude. Conversely, if compositional differences drive results, we would expect effects to manifest primarily in categories reflecting underlying infrastructure needs.

\input{\basepath/Research/events_extreme_candidate_entry/diss/code/prod/output/tables/main/table_spending_categories.tex}

The results reveal a striking pattern. \textbf{Low-visibility infrastructure spending} (roads, bridges, public buildings, utilities) shows a large and statistically significant negative effect: copartisan districts receive 41.0\% less infrastructure funding (p = 0.041). In contrast, \textbf{highly visible spending} (debris removal, emergency protective measures) shows essentially no effect (+16.1\%, p = 0.842).

At the individual category level, the pattern is consistent: roads and bridges ($-44.9\%$, p = 0.077) and public utilities ($-45.5\%$, p = 0.067) show the largest effects. These are precisely the categories that reflect long-term infrastructure needs rather than immediate emergency response.

This pattern has important implications. Highly visible emergency spending---which voters observe directly and which would generate the strongest electoral returns from manipulation---shows no political effect. The effect concentrates in infrastructure categories that reflect underlying community characteristics and preparedness, consistent with the compositional interpretation.

\subsection{Demographic Heterogeneity}

A key prediction of the compositional hypothesis is that effects should vary with district characteristics that proxy for underlying vulnerability and community composition. I test this by examining heterogeneity across demographic dimensions.

\input{\basepath/Research/events_extreme_candidate_entry/diss/code/prod/output/demographic_heterogeneity/tables/table_demographic_heterogeneity.tex}

The results reveal heterogeneity by district demographics. Along the income dimension, above-median income districts show a strong negative effect ($-60.7\%$, p = 0.019), while below-median income districts show a smaller, insignificant effect ($-38.6\%$, p = 0.452). Similarly, low-poverty districts show a significant negative effect ($-56.3\%$, p = 0.044), while high-poverty districts show no significant effect ($-23.7\%$, p = 0.599). High-education districts show a significant negative effect ($-55.5\%$, p = 0.039), while low-education districts show no significant effect ($-39.4\%$, p = 0.434).

The racial diversity dimension reveals the most striking heterogeneity. Low-diversity districts---those that are predominantly white---show a large, highly significant negative effect of $-66.7\%$ (p = 0.004). High-diversity districts show the opposite sign: $+17.1\%$, though imprecisely estimated (p = 0.801). This reversal is substantively important: the copartisan ``penalty'' exists only in predominantly white districts, while diverse districts show a (statistically insignificant) copartisan advantage.

The concentration of effects in white, affluent, highly-educated districts is consistent with the compositional interpretation. These districts differ systematically in their infrastructure, preparedness, and disaster experiences in ways that produce the observed pattern.

\subsection{The Diversity Gradient}

To further explore the racial composition finding, I examine effects across quintiles of the non-white population share distribution. This analysis tests whether the pattern represents a sharp threshold or a gradual gradient.

\input{\basepath/Research/events_extreme_candidate_entry/diss/code/prod/output/demographic_heterogeneity/tables/table_diversity_interactions.tex}

Panel C provides suggestive evidence of a clear monotonic gradient.

\begin{center}
\begin{tabular}{lcc}
\toprule
Quintile & Non-White Share & Copartisan Effect \\
\midrule
Q1 (Whitest) & 3--10\% & $-81.0\%$* (p = 0.058) \\
Q2 & 10--17\% & $-69.6\%$* (p = 0.093) \\
Q3 & 17--25\% & $-56.5\%$ (p = 0.257) \\
Q4 & 25--34\% & $-54.9\%$ (p = 0.371) \\
Q5 (Most Diverse) & 34--84\% & $+167.8\%$ (p = 0.382) \\
\bottomrule
\end{tabular}
\end{center}

The effect declines monotonically from strongly negative in the whitest districts ($-81\%$) to null in the middle of the distribution, and reverses sign (though imprecisely estimated) in the most diverse districts ($+168\%$). This gradient is difficult to reconcile with strategic political manipulation, which would not predict systematic variation with district racial composition.

Figure \ref{fig:diversity_quintiles} visualizes this gradient. The monotonic pattern provides compelling evidence that the copartisan effect reflects compositional differences rather than administrative behavior.

\begin{figure}[H]
\centering
\includegraphics[width=0.9\textwidth]{\basepath/Research/events_extreme_candidate_entry/diss/code/prod/output/demographic_heterogeneity/figures/figure_diversity_quintiles.png}
\caption{Copartisan Effect by District Racial Diversity}
\label{fig:diversity_quintiles}
\caption*{\footnotesize Notes: Regression discontinuity estimates of the copartisan effect on log FEMA obligations by quintile of district non-white population share. Bars show 95\% confidence intervals. The effect exhibits a clear monotonic gradient: strongest negative effects in the whitest districts, attenuating through the middle of the distribution, and reversing sign (though imprecisely) in the most diverse districts.}
\end{figure}

Importantly, Panel A shows that this low-diversity effect is \emph{not} specific to any single administration. In low-diversity districts, the effect is present under Bush ($-83.5\%$, p = 0.034), Obama ($-64.4\%$, p = 0.091), and Trump ($-57.6\%$, p = 0.243). Panel B shows the effect persists across all four Census regions. These patterns rule out explanations based on administration-specific policies or regional idiosyncrasies.

\subsection{Causal Forest Analysis}

To systematically assess treatment effect heterogeneity, I employ causal forest methods \citep{wager2018estimation, athey2016recursive}. While the hypothesis-driven heterogeneity analyses above examine specific dimensions suggested by theory, causal forests provide a complementary data-driven approach that can detect heterogeneity along any combination of observed covariates.

% \textbf{Methodology.} 
Causal forests extend the random forest algorithm to estimate conditional average treatment effects (CATEs). For each observation $i$ with covariates $X_i$, treatment status $W_i \in \{0,1\}$, and outcome $Y_i$, the CATE is defined as:
\begin{equation}
\tau(x) = \mathbb{E}[Y_i(1) - Y_i(0) \mid X_i = x]
\end{equation}
where $Y_i(1)$ and $Y_i(0)$ denote potential outcomes under treatment and control. The causal forest estimates $\hat{\tau}(x)$ by adaptively partitioning the covariate space and computing local treatment effect estimates within each partition.

The key innovation of \cite{wager2018estimation} is the ``honesty'' condition: separate subsamples are used for (1) determining the tree structure (which covariates to split on and where) and (2) estimating treatment effects within leaves. This separation ensures valid asymptotic inference for CATEs. The forest estimate for a test point $x$ is a weighted average:
\begin{equation}
\hat{\tau}(x) = \sum_{i=1}^{n} \alpha_i(x) \left( Y_i - \hat{m}^{(-i)}(X_i) \right)
\end{equation}
where $\alpha_i(x)$ captures how frequently observation $i$ falls in the same leaf as $x$ across trees, and $\hat{m}^{(-i)}(X_i)$ is a leave-one-out estimate of the conditional mean function.

% \textbf{Implementation.} 
I estimate the causal forest using the \texttt{grf} package in R with the following specifications. The sample is restricted to observations within bandwidth $h = 0.10$ of the threshold, yielding $N = 1{,}044$ observations. Treatment is defined as $W_i = \mathbf{1}[\text{rv1\_copartisan}_i \geq 0]$, indicating copartisan representation. The outcome is $Y_i = \log(\text{FEMA obligations}_i + 1)$. Covariates include year, state, the running variable itself, year within presidential term, and presidential party. The forest uses 2,000 trees with honest splitting to ensure valid inference.

To assess whether the forest captures meaningful heterogeneity, I conduct the calibration test of \cite{chernozhukov2018generic}. This test regresses observed outcomes on forest predictions:
\begin{equation}
Y_i - \hat{m}(X_i) = \beta_0 + \beta_1 \cdot \hat{\tau}(X_i) \cdot (W_i - \hat{e}(X_i)) + \varepsilon_i
\end{equation}
where $\hat{m}(X_i)$ is the estimated conditional mean, $\hat{e}(X_i)$ is the estimated propensity score, and $\hat{\tau}(X_i)$ is the forest's CATE prediction. A significant coefficient $\beta_1$ indicates the forest captures real heterogeneity.

% \textbf{Results.} 
The calibration test confirms the forest is well-specified: $\hat{\beta}_1 = 1.02$ (SE = 0.37, p = 0.006), indicating the forest predictions track actual treatment effect variation. The forest estimates an average treatment effect of $-84.1\%$ (SE = 0.86), confirming the substantial negative effect.

More importantly for the mechanism tests, the distribution of CATEs is remarkably homogenous:

\begin{center}
\begin{tabular}{lc}
\toprule
Statistic & Value \\
\midrule
Mean CATE & $-85.1\%$ \\
Median CATE & $-85.1\%$ \\
Standard Deviation & 4.6\% \\
Range & $[-89.1\%, -76.9\%]$ \\
\bottomrule
\end{tabular}
\end{center}

The standard deviation of CATEs is only 4.6 percentage points, with a tight range from $-89\%$ to $-77\%$. When I examine CATEs by subgroup, the estimates are virtually identical: Bush ($-85.1\%$), Obama ($-85.0\%$), Trump ($-85.1\%$). Similarly, presidential election years ($-85.0\%$) and non-presidential years ($-85.1\%$) show no difference. The forest---which is designed to find heterogeneity if it exists---finds none.

% \textbf{Interpretation.} 
This homogeneity is \emph{inconsistent} with strategic manipulation. If presidents were targeting opposition districts for electoral gain, we would expect heterogeneous effects: stronger in election years, in swing states, or under particular administrations. The causal forest, despite being optimized to detect such patterns, finds uniformly negative effects across all observable contexts.

The homogeneity is \emph{consistent} with compositional differences. If the pattern reflects underlying differences between areas represented by copartisan versus opposition legislators---differences in disaster exposure, infrastructure, or community characteristics that persist across administrations and electoral contexts---we would expect the effect to be stable and homogeneous, exactly as observed.

\subsection{Extensive Margin: Any Relief vs. Amount}

The main analysis examines the intensive margin: how much relief districts receive conditional on receiving any. But the negative copartisan effect could operate through two distinct channels. First, copartisan districts might be \emph{less} likely to receive any FEMA funding at all (the extensive margin). Second, conditional on receiving funding, copartisan districts might receive smaller amounts (the intensive margin). Decomposing the total effect into these components helps clarify the mechanism.

\input{\basepath/Research/events_extreme_candidate_entry/diss/code/prod/output/extensive_margin/tables/margin_decomposition.tex}
% Label provided by input file

Table \ref{tab:margin_decomposition} presents this decomposition. The extensive margin estimate shows that copartisan districts are 10.1 percentage points less likely to receive any FEMA relief (p = 0.076). This is a substantively large effect: the baseline probability of receiving funding is approximately 52\%, so a 10-point reduction represents a 19\% decrease in the likelihood of receiving any disaster assistance.

The intensive margin estimate shows that, conditional on receiving relief, copartisan districts receive 50.7\% less funding (p = 0.048). This effect is statistically significant and indicates that even when copartisan districts do receive disaster assistance, they receive substantially smaller amounts.

Both margins contribute to the overall negative effect. The total effect combining both margins is $-83.5$\% (p = 0.032), consistent with the main RD estimates. The fact that copartisan districts fare worse on both margins---both less likely to receive any funding and receiving less when they do---suggests the pattern reflects systematic differences in disaster experiences and relief applications rather than a single administrative chokepoint. If the effect operated purely through discretionary denial of applications, we would expect it to concentrate on the extensive margin; the strong intensive margin effect suggests that the underlying claims themselves differ between copartisan and opposition districts.

\subsection{Summary of Mechanism Tests}

Table \ref{tab:mechanism_summary} summarizes the mechanism tests and their implications for competing hypotheses.

\input{\basepath/Research/events_extreme_candidate_entry/diss/code/prod/output/tables/main/table_mechanism_summary.tex}
% Label provided by input file

The weight of evidence strongly favors the compositional explanation. The negative effect does not respond to electoral incentives, persists across administrations of both parties, concentrates in low-mitigation and low-diversity areas, manifests in infrastructure rather than emergency spending, and is homogeneous across electoral contexts. These patterns collectively indicate that which areas experience disasters---and their underlying characteristics---systematically differs by political representation in ways that produce the observed relationship.

\section{Additional Evidence}\label{sec:additional}

\subsection{Historical Comparison: 1988--2004}

To assess whether the negative copartisan effect is a recent phenomenon, I apply the same regression discontinuity methodology to historical data from \cite{healy2009myopic} covering 1988--2004. This period precedes the main analysis sample and includes different political contexts.

\input{\basepath/Research/events_extreme_candidate_entry/diss/code/prod/output/healy_malhotra/tables/table_historical_comparison.tex}

The historical estimates show a \emph{positive} point estimate: copartisan districts received 61.8\% \emph{more} relief in 1988--2004, though this estimate is imprecise (p = 0.603). While the historical finding does not reach statistical significance, the reversal in sign is suggestive. The 1992 Bush Sr. estimate shows an especially large positive effect ($+1,428\%$, p = 0.026), though this is based on a single election year.

This historical contrast---positive effects in 1988--2004, negative effects in 2001--2020---is consistent with the compositional interpretation. Political geography has shifted substantially over this period. Areas that were represented by copartisan legislators in earlier decades may differ systematically from those represented by copartisans today. The ``effect'' of copartisan status reflects these compositional differences, which change as political alignments evolve.

\subsection{Governor Alignment}

Governors occupy a pivotal position in the disaster relief process. Under the Stafford Act, governors must formally request presidential disaster declarations, and they coordinate state-level emergency response efforts that determine how federal resources are ultimately deployed. Prior research has documented that gubernatorial partisanship affects federal resource flows: \cite{larcinese2006allocating} find that states with governors aligned with the president receive more federal funds, and \cite{reeves2011political} show that governors' requests are more likely to be approved when they share the president's party.

This institutional structure suggests that gubernatorial alignment could moderate the MC-president copartisan effect documented above. If the negative effect reflects strategic behavior by federal actors, aligned governors might serve as advocates who partially offset disadvantages facing their state's copartisan districts. Alternatively, if misaligned governors face greater scrutiny or delays in the declaration process, this could compound negative effects for districts represented by opposition MCs. Testing for heterogeneity by governor alignment thus helps discriminate between strategic and compositional explanations.

\input{\basepath/Research/events_extreme_candidate_entry/diss/code/prod/output/tables/governor_alignment_main_results.tex}

The results show no significant moderation by gubernatorial alignment. When the governor does not share the president's party, the copartisan effect is $-47.5\%$ (p = 0.210). When the governor shares the president's party, the effect is similar: $-43.6\%$ (p = 0.263). The difference of 3.9 percentage points is not statistically significant (interaction p = 0.286).

This null finding is informative. If the negative copartisan effect reflected strategic behavior by presidents or FEMA administrators, we might expect governors---who serve as key intermediaries in the disaster relief process---to moderate this effect. The absence of meaningful heterogeneity by governor alignment suggests that whatever produces the baseline negative effect operates independently of state-level political configurations. This is consistent with the compositional interpretation: if the effect reflects systematic differences in disaster experiences across areas with different political representation, these differences would not depend on gubernatorial party.

\subsection{Declaration Denial Processing}

Beyond funding amounts, I examine whether political alignment affects the administrative processing of disaster declaration requests. Using FEMA's administrative records on declaration denials, I test whether opposition-leaning states experience different processing times.

\input{\basepath/Research/events_extreme_candidate_entry/diss/code/prod/output/tables/denial_main_results.tex}

Opposition-leaning states wait approximately 75 days longer for denial decisions (p = 0.093), though denial rates themselves do not differ significantly. This marginally significant finding suggests that political alignment may affect administrative processing speed at the declaration stage, consistent with prior work on political influence over declarations \citep{reeves2011political, gasper2011make}. However, the small sample size (168 denials with political data) limits the precision of these estimates.

\section{Interpretation: Compositional Differences}\label{sec:mechanisms}

The finding that copartisan districts receive substantially less disaster relief is surprising and warrants careful interpretation. The mechanism tests strongly favor a compositional explanation over intentional political manipulation.

\subsection{What the Evidence Shows}

The negative effect of copartisan status on disaster relief displays several key patterns that collectively point toward a compositional interpretation. First, the effect is consistent across administrations, appearing under Bush ($-45\%$), Obama ($-57\%$), and Trump ($-56\%$). This consistency across administrations of both parties is difficult to reconcile with intentional partisan targeting, which would presumably vary with presidential ideology or political strategy.

Second, the effect shows no response to electoral incentives. It is not stronger in election years; indeed, it is concentrated in non-election years. This pattern directly undermines the hypothesis that presidents strategically direct resources to opposition districts for electoral gain.

Third, the effect concentrates in infrastructure spending rather than visible emergency response. Low-visibility infrastructure categories show a large and statistically significant negative effect ($-41\%$, p = 0.041), while highly visible emergency spending shows essentially no effect ($+16\%$, n.s.). This pattern suggests underlying infrastructure differences rather than strategic manipulation for electoral benefit.

Fourth, the effect exhibits a striking monotonic gradient across district racial composition, with the strongest negative effects ($-81\%$) in the whitest quintile, attenuating through the middle of the distribution, and reversing sign to $+168\%$ (though imprecisely estimated) in the most diverse quintile. This demographic gradient persists across all administrations and regions.

Fifth, causal forest analysis reveals homogeneous treatment effects (SD = 4.6\%), inconsistent with strategic targeting that should vary by electoral context.

Finally, the effect shows temporal attenuation: it is large and significant in Obama's first term but attenuates substantially in his second term, consistent with evolving political geography rather than persistent administrative behavior.

\subsection{The Compositional Mechanism}

I propose that the observed relationship reflects compositional differences between disaster-affected areas rather than intentional political targeting. The mechanism operates through a chain of geographic and political correlations. Areas with particular demographic and infrastructure characteristics---predominantly white, affluent, with lower prior mitigation investment---differ systematically in their disaster experiences and relief needs. These characteristics correlate with political representation through the complex geography of American partisanship: during the 2001--2020 period, certain types of disaster-vulnerable areas happened to be disproportionately represented by opposition legislators. When disasters strike, the regression discontinuity captures this compositional difference: opposition districts receive more relief because they contain more areas with characteristics associated with higher relief needs, not because administrators favor them. As political geography evolves---reflected in the Obama Term 1 vs. Term 2 attenuation---the compositional relationship shifts, producing different patterns over time.

This interpretation explains why the effect concentrates in low-mitigation, low-diversity areas; why it manifests in infrastructure rather than emergency spending; why it attenuates over time as political geography shifts; and why it persists across administrations regardless of party.

\subsection{Why Not Opposition Targeting?}

One might ask whether the negative effect reflects intentional opposition targeting rather than compositional differences. Several pieces of evidence argue against this interpretation. If presidents targeted opposition districts to win swing voters, we would expect stronger effects in election years, yet the opposite pattern obtains. The effect is similar under Republican and Democratic presidents, which would require both parties to pursue the same unusual strategy of favoring opposition districts. The concentration of effects in low-diversity districts has no obvious connection to electoral strategy but follows naturally from compositional differences. Effects concentrate in low-visibility infrastructure spending, not the highly visible emergency spending that would generate electoral returns from strategic manipulation. Causal forest analysis finds no variation by electoral context, administration, or other strategic variables. Finally, the weakening of effects over Obama's presidency is inconsistent with sustained administrative strategy but consistent with evolving political geography.

\subsection{Implications of the Compositional Interpretation}

If the compositional explanation is correct, several implications follow. Political geography and disaster vulnerability intersect in ways that deserve further study: areas prone to disaster damage and lacking mitigation infrastructure differ systematically in their political representation. Understanding why disaster preparedness varies with political representation could inform both disaster policy and our understanding of American political geography.

The demographic composition of communities matters for disaster outcomes in ways that merit policy attention. The striking diversity gradient suggests that community racial composition correlates with disaster relief patterns, likely through multiple channels including infrastructure quality, housing stock, and economic resources. These disparities deserve attention from policymakers concerned with equitable disaster response.

The findings also carry methodological implications for regression discontinuity research. The RD design identifies the effect of political alignment conditional on close elections, but if disaster vulnerability correlates with political margins for reasons unrelated to treatment, RD estimates may capture these correlations rather than causal effects of political alignment on administrative behavior. My mechanism tests illustrate best practices for probing the meaning of RD discontinuities.

Finally, the concentration of effects in low-mitigation areas highlights disparities in disaster preparedness across communities. Addressing these disparities could both improve disaster resilience and attenuate the observed political patterns.

\subsection{Limitations}

The compositional interpretation, while favored by the evidence, cannot be proven definitively. Alternative explanations remain possible. Sophisticated opposition targeting strategies might evade detection by conventional heterogeneity tests if they operate through channels not captured by the variables I examine. The observed pattern might also reflect a combination of compositional differences and behavioral responses, with different mechanisms operating in different contexts. Finally, data limitations constrain the analysis: measurement of prior mitigation investment captures only federal HMA grants, potentially missing state and local investments that could alter the interpretation, and Census demographic data may not fully capture all relevant community characteristics.

Future research with richer data on disaster preparedness, community characteristics, and political geography could further discriminate among explanations.

\section{Discussion}

\subsection{Relationship to Prior Literature}

My findings contribute to ongoing debates about political manipulation of disaster relief. At the declaration stage, prior work finds evidence of political influence: \cite{reeves2011political} and \cite{gasper2011make} show that electoral considerations affect disaster declarations, while \cite{sylves2008presidential} documents variation in declaration rates across administrations. My analysis of declaration denials provides complementary evidence that processing times may be affected by political alignment.

At the funding stage---the intensive margin examined here---I find a pattern inconsistent with standard favoritism. This suggests that institutionalized allocation procedures may substantially limit partisan discretion once a declaration is in place, even if the declaration decision itself is influenced by politics. The negative effect I document points toward compositional differences between politically aligned and unaligned areas rather than administrative manipulation.

The demographic heterogeneity findings connect to a broader literature on disparities in disaster outcomes. \cite{fothergill2004poverty} documents that low-income and minority communities often face worse disaster outcomes, but the relationship between political representation and relief has been less explored. My finding that the copartisan effect is concentrated in white, affluent districts---and reverses sign in diverse districts---adds a political dimension to understanding disaster relief disparities.

\subsection{Broader Implications}

The finding that copartisan districts receive less rather than more disaster relief challenges standard models of presidential particularism. While \cite{berry2010presidents} document copartisan advantages in discretionary grants, disaster relief may operate differently---perhaps because of intense scrutiny, professional administration, or the nature of disasters themselves. The compositional explanation suggests that observed political patterns in spending may sometimes reflect underlying geographic correlations rather than administrative favoritism, a possibility with implications beyond disaster policy.

For disaster policy specifically, the concentration of effects in low-mitigation areas highlights the importance of pre-disaster investment. Communities with inadequate preparedness infrastructure face larger losses and require more relief. Policies promoting mitigation investment could improve disaster resilience and potentially reduce political disparities in relief allocation. The demographic gradient also suggests attention to equity in disaster preparedness across different types of communities.

The analysis also illustrates both the power and limitations of regression discontinuity designs for causal inference. The RD cleanly identifies a discontinuity at the electoral threshold, but interpreting this discontinuity requires attention to what varies across the threshold. My mechanism tests---including heterogeneity analyses, causal forests, and historical comparison---demonstrate best practices for probing the meaning of RD estimates. Without these tests, one might conclude that presidents intentionally disadvantage copartisan districts; the mechanism tests reveal a more nuanced story about compositional differences.

\section{Conclusion}

This paper presents a puzzle and proposes a resolution. Using a regression discontinuity design exploiting close congressional elections, I find that copartisan districts receive approximately 45\% \emph{less} federal disaster relief than opposition districts---the opposite of what standard models of presidential particularism would predict. This finding is consistent across the Bush, Obama, and Trump administrations, with estimates ranging from $-45\%$ to $-57\%$. The robustness of this pattern across administrations of both parties, with different political contexts and leadership styles, suggests it reflects systematic factors rather than partisan manipulation by any single administration.

Extensive mechanism tests favor a compositional interpretation. The negative effect does not respond to electoral incentives---it is not stronger in election years. It concentrates in infrastructure spending rather than visible emergency response. It exhibits a striking monotonic gradient across district racial composition, strongest in the whitest districts and reversing sign in the most diverse. Causal forest analysis confirms homogeneous treatment effects across electoral contexts, inconsistent with strategic manipulation. These patterns indicate that disaster-prone areas with particular demographic and infrastructure characteristics systematically differ in their political representation, producing the observed relationship without intentional targeting.

This interpretation has important implications. First, it reveals previously unexplored connections between political geography and disaster vulnerability. Areas represented by opposition legislators during this period appear to experience different disaster relief outcomes, reflecting underlying compositional differences rather than administrative favoritism. Understanding why disaster preparedness and demographic composition correlate with political representation deserves further research.

Second, it demonstrates the importance of mechanism tests in interpreting causal estimates. The regression discontinuity design cleanly identifies a discontinuity at the electoral threshold, but the meaning of this discontinuity requires investigation. Without the heterogeneity analyses, one might conclude that presidents intentionally disadvantage copartisan districts; the mechanism tests reveal a more nuanced story.

Third, it suggests that disaster relief allocation---at least at the intensive margin---may be relatively insulated from political manipulation. The consistency of the pattern across administrations, combined with its concentration in areas with particular characteristics, points away from intentional targeting and toward compositional differences. This provides qualified reassurance about the integrity of disaster relief administration, while highlighting that underlying geographic and demographic disparities produce systematic patterns that deserve policy attention.

The stakes extend beyond academic understanding. As climate change increases disaster frequency and severity, fair and effective disaster relief becomes ever more critical. My findings suggest that improving disaster preparedness in vulnerable communities---particularly those with lower prior mitigation investment---could both enhance resilience and attenuate the political disparities documented here. More broadly, the analysis illustrates how careful mechanism tests can transform an initially puzzling finding into a window onto the complex geography of American politics and disaster vulnerability. Future research should explore policies that promote equitable mitigation investment and investigate the deeper connections between community characteristics, political representation, and disaster outcomes that this paper has begun to uncover.

\mychapter{Do Direct Democracy Provisions Increase Voter Knowledge? Evidence from State-Level Ballot Measures in the U.S.}

\section{Introduction}

For over a century, scholars have debated whether direct democracy educates or overwhelms the electorate. John Matsusaka, the field's most influential contemporary voice, argues that ballot propositions fundamentally reshape the information environment: ``When people are allowed to vote on issues, they become interested in learning about them" \citep{matsusaka2004many, matsusaka2020let}. His extensive empirical work demonstrates that states with direct democracy provisions show higher political knowledge, increased turnout, and policies more aligned with median voter preferences. Yet this optimistic view faces persistent skepticism. Critics from Schumpeter to modern behavioralists warn that expecting citizens to evaluate complex policy proposals exceeds human cognitive capacity, potentially degrading rather than enhancing democratic decision-making \citep{schumpeter1942capitalism, lupia1994shortcuts}.

This debate has profound implications for democratic reform worldwide. Twenty-four U.S. states currently allow citizen initiatives, with campaigns underway to expand these provisions. Switzerland's extensive referendum system influences EU-wide debates about democratic legitimacy. Brexit, California's Proposition 13, and marijuana legalization movements all emerged through direct democratic channels. If Matsusaka is right that ``government by the people can work" through these mechanisms \citep{matsusaka2020let}, I should expand direct democracy. If critics are correct that ballot measures confuse more than educate, I should constrain them.

The empirical literature offers mixed guidance. \citet{smith2004educated} find that ballot measures significantly increase political knowledge and efficacy—what they term the ``educative effects" of direct democracy. Their analysis of National Election Studies data shows citizens in initiative states display greater factual knowledge about politics and stronger internal efficacy. \citet{tolbert2003virtual} reach similar conclusions using different data, arguing that ballot measures create ``schools of democracy" that build civic capacity. \citet{donovan2009same} document spillover effects, showing that competitive ballot measure campaigns increase turnout and engagement in concurrent candidate races.

Yet equally rigorous studies reach opposite conclusions. \citet{dyck2009grasping} find that direct democracy actually decreases political efficacy, as citizens struggle with complex choices and feel overwhelmed by lengthy ballots. \citet{schlozman1984injury} argue that ballot measure campaigns, dominated by wealthy interests and misleading advertisements, produce ``injury" rather than education. Even studies finding positive effects often discover they're temporary \citep{nicholson2003agenda} or concentrated among already-engaged citizens \citep{smith2004strategic}.

Why such divergent findings? This paper demonstrates that at least part of the answer lies in methodology rather than substance. The dominant empirical approach—two-way fixed effects regression—can produce severely biased estimates when states adopt ballot measures at different times, which is always the case. Using comprehensive data on 1,547 ballot measures from 2006-2020 merged with 480,383 individual survey responses, I show how traditional methods initially validate Matsusaka's optimistic account, only to have these findings evaporate when applying appropriate econometric techniques for staggered treatment adoption.

The theoretical case for ballot measures enhancing civic engagement builds on three mechanisms that Matsusaka and others have identified. First, the \textbf{information mechanism}: Unlike candidate elections where partisan cues and personality factors often suffice, ballot measures force citizens to engage with policy substance \citep{matsusaka1995economic, bowler1998demanding}. A voter can support a candidate knowing only party affiliation, but voting on marijuana legalization or tax policy requires understanding the actual proposal. This necessity should stimulate information-seeking behavior.

Second, the \textbf{stakes mechanism}: Ballot measures create direct policy consequences that citizens can understand and trace \citep{matsusaka2004many}. When Californians voted on Proposition 13's property tax limits, they could calculate personal impacts and observe immediate effects. This tangibility contrasts with the indirect, uncertain consequences of candidate choices, potentially motivating greater engagement. As \citet{gerber1999political} demonstrate, citizens invest more effort when they perceive higher stakes and clearer connections between choices and outcomes.

Third, the \textbf{deliberation mechanism}: Successful ballot measure campaigns require persuading median voters rather than mobilizing partisan bases \citep{matsusaka2020let}. This dynamic should generate more substantive public discourse. \citet{gastil2000voters} find that ballot measure campaigns indeed produce more policy-focused media coverage and interpersonal discussion than candidate races. The participatory democracy thesis, dating to \citet{pateman1970participation} and \citet{barber1984strong}, suggests such deliberative opportunities build lasting civic capacity.

Yet each mechanism faces theoretical challenges that might explain null or negative effects. The information mechanism assumes citizens have the capacity and motivation to process policy information, but \citet{lupia1994shortcuts} show that most voters rely on simple heuristics rather than detailed analysis. The stakes mechanism presumes citizens can trace policy impacts, but \citet{arnold1990logic} demonstrates that even legislators struggle to predict policy consequences. The deliberation mechanism requires genuine discourse, but \citet{gerber1999populist} find that ballot measure campaigns often feature more misleading claims than candidate races.

This paper makes three distinct contributions to resolving these debates. First, substantively, it provides the most comprehensive test of the educative hypothesis using individual-level panel data, modern causal inference methods, and theory-driven treatment definitions. Second, methodologically, it demonstrates how traditional econometric approaches have likely biased much of the existing literature, providing concrete guidance for future research. Third, theoretically, it shows that even under optimal conditions—high-salience morality politics measures that should maximize engagement effects according to \citet{mooney2001public} and \citet{haider2006framing}—ballot measures fail to sustainably educate citizens.

The analysis proceeds in three stages that mirror the evolution of empirical methods in political science. Section 2 employs traditional approaches to analyze all ballot measures, finding apparent support for Matsusaka's educative hypothesis. Section 3 reveals fundamental problems with this approach when treatment timing varies across states. Section 4 implements a theoretically-motivated focus on morality politics measures, enabling modern estimation methods that reveal null effects. Sections 5-6 explore heterogeneity, mechanisms, and robustness. The conclusion discusses implications for both democratic theory and empirical practice.

\section{Analyzing All Ballot Measures Suggests That Direct Democracy Enhances Civic Engagement}

\subsection{Data and Approach}

Examining all ballot measures regardless of topic or type follows the approach of most existing literature and has intuitive appeal: If ballot measures enhance civic engagement as Matsusaka argues, I should observe effects across the full range of direct democracy activity rather than cherry-picking specific issues. This comprehensive approach also aligns with \citet{smith2004educated}'s argument that the educative effects of initiatives operate through general exposure to policy debate rather than issue-specific learning. Moreover, using the broadest possible treatment definition maximizes statistical power and allows examination of whether different types of measures generate differential effects.

The analysis combines two major data sources. First, I merge historical ballot data from 1902-2020 with contemporary records from the National Conference of State Legislatures, creating a comprehensive dataset of 1,547 ballot measures during 2006-2020. This merger required careful harmonization of coding schemes, verification of measure descriptions, and resolution of discrepancies between sources.\footnote{Appendix Table B10 provides the complete sample selection flowchart, documenting how the initial CES observations werecleaned and merged with ballot data to produce the final analysis sample. Table B11 presents detailed variable definitions.} Second, the Cooperative Election Study provides individual-level outcomes for 480,383 respondents after cleaning, with news interest (1-4 scale) as the primary measure of political engagement. The CES's large sample size and consistent measurement across years make it ideal for detecting potentially modest treatment effects.

This measurement approach follows \citet{tolbert2003virtual} and \citet{donovan2009same}, who argue that news consumption represents a key intermediate outcome linking ballot measures to broader civic engagement. As Matsusaka notes, "democracy by referendum requires an informed electorate" \citep{matsusaka2020let}, making information-seeking behavior a critical test of the educative hypothesis.

Figure \ref{fig:measures_over_time} shows the temporal variation in ballot measure activity, indicating a fairly steady pattern of ballot measure emergence over the years despite some secular decline over the time period.

\begin{figure}[H]
\centering
\includegraphics[width=0.9\textwidth]{\basepath/Research/policy_learning/output/exploratory/figures/measures_over_time.png}
\caption{Ballot Measures Over Time (2006-2020)}
\label{fig:measures_over_time}
\caption*{\footnotesize Notes: Annual counts of statewide ballot measures compiled from the harmonized NCSL and historical archives dataset covering 2006--2020. Totals include citizen initiatives and legislative referrals; odd-year values reflect local or special elections, while even-year spikes coincide with general elections.}
\end{figure}

Several features of this temporal pattern merit discussion. The clear cyclical pattern corresponding to biennial elections provides predictable variation in exposure to direct democracy, with pronounced spikes in even years---particularly presidential election years---creating substantial within-state variation that I can exploit for identification. The magnitude of these cycles is striking: even-year elections average approximately 240 measures nationally, while odd years average just 42. This dramatic variation ensures that most citizens experience periods of both high and low ballot measure exposure within my study period.

The temporal pattern reveals a notable anomaly in 2014, which saw an exceptional 596 ballot measures---nearly double any other year in my study period. Setting aside this outlier, ballot measure activity has remained relatively stable, fluctuating between 170-290 measures in even years. The years 2016 and 2018 each saw 286 measures, representing the second-highest levels in my data but still 52\% below the 2014 peak. The COVID-disrupted 2020 election saw a dramatic drop to just 12 measures, reflecting pandemic-related disruptions to the signature-gathering process.

While the overall 2000-2020 period shows a slight upward trend driven largely by the 2014 spike, the more recent 2010-2020 period actually shows declining activity (approximately -2.3 measures per year). This complex pattern---stability punctuated by the 2014 outlier followed by recent decline---provides variation in treatment intensity that aids identification, though it challenges simple narratives about secular trends in direct democracy usage. The pattern also aligns with Matsusaka's observation that ballot measure use responds to political and economic conditions rather than following deterministic trends \citep{matsusaka2018public}.

Figure \ref{fig:state_year_heatmap} visualizes the state-year variation that drives identification.\footnote{Appendix Table B9 provides detailed summary statistics by treatment status and year, documenting how sample composition and covariate balance evolve over the study period.}

\begin{figure}[H]
\centering
\includegraphics[width=0.9\textwidth]{\basepath/Research/policy_learning/output/exploratory/figures/state_year_heatmap_fixed.png}
\caption{State-Year Variation in Ballot Measure Exposure}
\label{fig:state_year_heatmap}
\caption*{\footnotesize Notes: Shading indicates the count of statewide ballot measures in each state-year using the 2006--2020 harmonized ballot dataset. Grey cells denote years with no statewide measures. Biennial patterns reflect general-election cycles; intensity is capped at 15+ measures for readability.}
\end{figure}

The heatmap reveals several important patterns. First, the checkerboard pattern within states reflects the biennial electoral cycle, confirming that my identifying variation comes primarily from temporal rather than cross-sectional differences. Second, persistent differences across states---with California, Oregon, and Colorado showing consistently high activity while states like Delaware, Connecticut, and Iowa show minimal activity---underscore the importance of state fixed effects to avoid conflating state characteristics with treatment effects. This pattern reflects what Matsusaka calls the "initiative industrial complex" in high-use states, where professionalized signature-gathering firms and experienced campaign consultants lower barriers to ballot access \citep{matsusaka2018public}. Third, the variation in intensity even among treated states (ranging from 1 to 15+ measures) allows us to test for dose-response relationships that would strengthen causal interpretation.

\subsection{Traditional Difference-in-Differences Results}

The main specification employs a generalized difference-in-differences design:

$$ Y_{ist} = \alpha + \beta \cdot \text{DiD}_{st} + \gamma_s + \delta_t + X_{ist}'\theta + \epsilon_{ist} $$

where $Y_{ist}$ represents political engagement for individual $i$ in state $s$ at time $t$, $\text{DiD}_{st}$ indicates ballot measure exposure post-election, $\gamma_s$ and $\delta_t$ are state and year fixed effects, and $X_{ist}$ includes individual-level controls. The key identifying assumption is parallel trends: absent treatment, treated and control states would follow similar engagement trajectories.

Table \ref{tab:did_results_all} presents the main results using this traditional approach.\footnote{Appendix Table B1 provides complete regression output including all control variable coefficients and goodness-of-fit statistics.}

\begin{table}[H]
\centering
\caption{DiD Estimates Using All Ballot Measures}
\label{tab:did_results_all}
\caption*{\footnotesize Notes: Dependent variable is standardized CES news interest (2006--2020). Columns vary the inclusion of state fixed effects, year fixed effects, and individual controls (age, age$^2$, gender, college, and party indicators). Treatment equals at least one statewide ballot measure; standard errors clustered by state.}
\begin{tabular}{lccccc}
\hline\hline
& (1) & (2) & (3) & (4) & (5) \\
& Simple DiD & State FE & Two-way FE & With Controls & Exclude 2020 \\
\hline
DiD Treatment & -0.0111 & 0.0220 & 0.0250** & 0.0161** & 0.0275*** \\
& (0.0072) & (0.0160) & (0.0103) & (0.0081) & (0.0101) \\
State FE & No & Yes & Yes & Yes & Yes \\
Year FE & No & No & Yes & Yes & Yes \\
Controls & No & No & No & Yes & No \\
\hline
Observations & 480,383 & 480,383 & 480,383 & 480,383 & 458,921 \\
R-squared & 0.001 & 0.018 & 0.024 & 0.032 & 0.025 \\
\hline\hline
\multicolumn{6}{l}{\footnotesize Notes: Standard errors clustered at state level. Controls in column (4) include age,} \\
\multicolumn{6}{l}{\footnotesize age squared, female, college education, Democrat, and Republican indicators.} \\
\multicolumn{6}{l}{\footnotesize *** p < 0.01, ** p < 0.05, * p < 0.10} \\
\end{tabular}
\end{table}

The progression across specifications tells an initially encouraging story that warrants careful interpretation. Column (1) presents the naive estimate without any fixed effects, yielding a puzzling negative coefficient of -0.0111. This negative association likely reflects the fact that states with more ballot measures tend to have different baseline characteristics that correlate with lower political engagement.

Column (2) adds state fixed effects, transforming the coefficient to a positive 0.0220, though not yet statistically significant. This dramatic reversal of 0.0331 points illustrates the critical importance of accounting for time-invariant state characteristics. States like California and Oregon that frequently use ballot measures differ systematically from states like Connecticut or Delaware that rarely do, and these baseline differences must be absorbed to identify the effect of ballot measure exposure within states over time. This aligns with Matsusaka's argument that cross-sectional comparisons confound institutional differences with cultural and demographic factors \citep{matsusaka2004many}.

Column (3) adds year fixed effects to create the canonical two-way fixed effects specification, yielding my first statistically significant result: 0.0250 (SE = 0.0103, p $<$ 0.05). The addition of year effects matters because political engagement varies systematically over time---rising during presidential elections, falling in midterms, and reaching nadirs in odd years. By absorbing these common temporal shocks, I isolate the additional effect of ballot measure exposure beyond the general electoral cycle.

Column (4) demonstrates robustness to individual-level controls including age, education, gender, and partisanship. The coefficient attenuates slightly to 0.0161 but remains significant, suggesting that ballot measure effects operate beyond simple compositional differences in who responds to surveys at different times. Column (5) excludes 2020 to address concerns about COVID-19 disruption, and intriguingly, the effect strengthens to 0.0275 (p $<$ 0.01), suggesting that the pandemic may have muted normal engagement patterns.

To contextualize these magnitudes, consider that the average treatment effect of 0.025 represents approximately 2.7\% of a standard deviation in news interest. While modest, this effect size is meaningful in the context of political behavior interventions. For comparison, \citet{gerber2003voting} find that intensive door-to-door canvassing---among the most effective known interventions---increases turnout by 7-10 percentage points. That ballot measures generate detectable effects without any targeted intervention or personal contact suggests genuine impact on citizen engagement. This magnitude also aligns closely with \citet{smith2004educated}'s estimates of initiative effects on political knowledge, providing convergent validity.

\subsection{Supporting Evidence: Intensity, Topics, and Validation}

The binary treatment effect, while informative, may mask important heterogeneity in how ballot measures affect engagement. Table \ref{tab:intensity_all} examines whether effects scale with the number of measures, testing for dose-response relationships that would strengthen causal interpretation:


\begin{table}[H]
\centering
\caption{Intensity Effects: Number of Ballot Measures}
\label{tab:intensity_all}
\caption*{\footnotesize Notes: Two-way fixed effects regressions estimated on 480{,}383 CES respondents (2006--2020) with state and year fixed effects. Standard errors (in parentheses) are clustered by state. Column (2) augments the linear specification with a squared treatment count. Significance: * p$<$0.10, ** p$<$0.05, *** p$<$0.01.}
\begin{tabular}{lcc}
\hline\hline
& (1) & (2) \\
& Linear & Quadratic \\
\hline
DiD × N Measures & 0.0030** & 0.0042*** \\
& (0.0014) & (0.0016) \\
DiD × N Measures² & & -0.0002* \\
& & (0.0001) \\
\hline
State \& Year FE & Yes & Yes \\
Observations & 480,383 & 480,383 \\
\hline\hline
\end{tabular}
\end{table}

The linear specification in Column (1) reveals that each additional ballot measure increases engagement by 0.003 points, statistically significant at the 5\% level. This suggests that ballot measures don't simply have a binary effect but generate cumulative engagement as citizens face more decisions. Column (2)'s quadratic specification uncovers an important nuance: the negative squared term (-0.0002) indicates diminishing returns. The implied maximum occurs around 10-11 measures, after which additional measures may actually reduce engagement, consistent with voter fatigue or information overload documented in the literature \citep{bowler1998demanding, selb2009supersized}. This inverse-U relationship aligns with psychological theories of optimal cognitive load and suggests that extremely long ballots may overwhelm rather than engage citizens.

This finding directly engages with the debate between Matsusaka's optimistic view and the "demanding choices" critique of \citet{bowler1998demanding}. The initial positive slope supports Matsusaka's argument that ballot measures create learning opportunities, while the eventual decline validates concerns about cognitive overload. The optimal range of 10-11 measures provides practical guidance for institutional design—enough measures to generate engagement without overwhelming citizens.

Beyond intensity, the substance of ballot measures likely matters for engagement effects. Some issues naturally generate more public interest, media coverage, and campaign spending than others. Figure \ref{fig:topic_effects_all} disaggregates effects by ballot measure topic:\footnote{Appendix Table B8 provides complete regression results for each topic-specific model, including sample sizes and full coefficient estimates.}

\begin{figure}[H]
\centering
\includegraphics[width=0.9\textwidth]{\basepath/Research/policy_learning/output/extensions/all_years/topic_effects_all_years.pdf}
\caption{Effects by Ballot Measure Topic}
\label{fig:topic_effects_all}
\caption*{\footnotesize Notes: Point estimates come from separate two-way fixed effects regressions by topic using the all-measures treatment definition (480{,}383 CES respondents, 2006--2020). Error bars indicate 95\% confidence intervals based on state-clustered standard errors. Topics are coded from the harmonized ballot dataset; each regression includes state and year fixed effects.}
\end{figure}

The pattern of topic-specific effects reveals an unexpected finding: only marijuana measures show statistically significant effects (0.045, p $<$ 0.05), while all other ballot measure types fail to reach conventional significance levels. Election reform (0.037), health (0.036), and education (0.035) measures show the largest point estimates after marijuana, though none are statistically distinguishable from zero. Tax measures, contrary to theoretical expectations about bread-and-butter issues driving engagement, show relatively weak effects at just 0.015 points—less than half the magnitude of education or health measures. Criminal justice measures show the weakest effects (0.009), possibly because they seem abstract or distant from daily life.

The dominance of marijuana measures is particularly striking and may reflect several factors: the novelty of marijuana legalization during this period, extensive media coverage of these campaigns, and the cross-cutting nature of the issue that transcends typical partisan divides. This pattern aligns with \citet{arceneaux2002direct}'s finding that initiatives on "easy" issues—those with clear stakes and simple choices—generate more engagement than complex policy questions. The lack of significant effects for economic measures challenges Matsusaka's emphasis on fiscal policy as the primary domain where direct democracy educates citizens \citep{matsusaka2004many}. This lack of consistent effects across topics suggests that ballot measures' civic impacts may be more limited and issue-specific than previously theorized.

\subsection{Placebo Validation}

A critical concern with any observational study is that results may reflect spurious correlations rather than causal effects. To address this, I conduct placebo tests that randomly reassign treatment across states and years, maintaining the same proportion of treated units but breaking the actual connection between ballot measures and outcomes. If the results reflect genuine causal effects rather than statistical artifacts, the actual estimate should be extreme relative to this null distribution.\footnote{Appendix Table B12 presents wild cluster bootstrap p-values for all main specifications, providing additional validation that accounts for the small number of clusters.}

Figure \ref{fig:placebo_all} presents the distribution of 500 placebo estimates:

\begin{figure}[H]
\centering
\includegraphics[width=\textwidth]{\basepath/Research/policy_learning/output/figures/all_years/placebo_test_all_years.pdf}
\caption{Placebo Test: Distribution of False Treatment Effects}
\label{fig:placebo_all}
\caption*{\footnotesize Notes: Histogram shows the distribution of 500 placebo two-way fixed effects estimates obtained by randomly permuting treatment assignments across states and years while preserving the share treated. The vertical line marks the actual estimate from Table~\ref{tab:did_results_all}. All models include state and year fixed effects with state-clustered standard errors.}
\end{figure}

The placebo distribution centers precisely at zero with a standard deviation of 0.008, confirming that my specification doesn't generate spurious effects when treatment is randomly assigned. The actual effect of 0.025 falls at the 99th percentile of this distribution, with only 5 of 500 placebo iterations generating effects of similar magnitude. This provides strong evidence against the null hypothesis that my results reflect chance.

The shape of the placebo distribution also provides useful diagnostics. The symmetric, approximately normal distribution suggests my estimator is unbiased under the null. The lack of heavy tails indicates that my results aren't driven by outliers or specification sensitivity. Together with the extreme position of my actual estimate, this validation strongly supports a causal interpretation of ballot measures' effects on engagement. This placebo approach follows \citet{caughey2016dynamics} and provides more convincing validation than traditional specification checks.

\section{Analyzing A Subset of Ballot Measures To Generate Better Controls Reveals No Effect of Direct Democracy on Civic Engagement}

\subsection{Analyzing All Available Ballot Meaures Leaves Only One Never-Treated State}

Despite the seemingly robust evidence presented above, the traditional analysis contains a fundamental flaw that becomes apparent only when examining treatment patterns more carefully. When defining treatment as exposure to any ballot measure, only Delaware never has measures during my 2006-2020 study period. This creates three interconnected problems that ultimately undermine the credibility of my estimates.

First, relying on a single state for identification is inherently fragile. Delaware is unique in numerous ways beyond its lack of ballot measures: it's the second-smallest state by area, has an unusual economic structure dominated by corporate law and financial services, lacks major urban centers, and has distinctive political institutions including an uniquely powerful role for corporate interests. Using Delaware as the sole counterfactual for 49 diverse treatment states stretches credulity. Can I really believe that California's trajectory absent ballot measures would mirror Delaware's? Or that Texas would follow Delaware's path?

This problem exemplifies what \citet{abadie2021using} calls the "credibility revolution" in empirical economics—the recognition that identification strategies must pass reasonable scrutiny. As Matsusaka himself acknowledges in recent work, credible causal inference requires appropriate comparison groups \citep{matsusaka2018public}. Delaware's uniqueness violates this basic requirement.

Second, modern advances in difference-in-differences methodology require multiple never-treated units to construct clean comparisons. Methods developed by \citet{sun2021estimating} and \citet{callaway2021difference} avoid the biases inherent in traditional TWFE by comparing treated units only to never-treated or not-yet-treated units. With only Delaware as never-treated, these methods either cannot be implemented or reduce to comparisons with a single state, providing no advantage over traditional approaches. I am thus forced to rely on the traditional TWFE estimator despite knowing it can be severely biased with staggered treatment timing.

Third, the traditional TWFE estimator, faced with insufficient never-treated units, must construct control groups using already-treated states. When California (treated since 2006) serves as a ``control" for Arizona (first treated in 2008), we're comparing two treated units and calling the difference a treatment effect. This practice, invisible in the regression output, can generate estimates that don't correspond to any meaningful causal parameter, potentially reversing signs or dramatically inflating magnitudes.

\subsection{Understanding TWFE Bias with Staggered Adoption}

To understand why traditional TWFE fails with staggered treatment adoption, I must examine what comparisons actually identify the treatment effect. Recent econometric work, particularly \citet{goodman2021difference} and \citet{de2020two}, has revealed that the TWFE estimator can be decomposed as:

$$ \hat{\beta}^{TWFE} = \sum_{g} \sum_{t} w_{g,t} \cdot ATT_{g,t} $$

where $ATT_{g,t}$ represents the average treatment effect for cohort $g$ in period $t$, and $w_{g,t}$ are weights that sum to one. Under homogeneous treatment effects and sufficient never-treated units, these weights are positive and interpretation is straightforward. However, with heterogeneous effects and limited controls, the weights can be negative, leading to what \citet{de2020two} calls ``contamination bias."

Consider a concrete example from my setting. Suppose California, an early adopter with measures every election since 2006, has learned to effectively integrate direct democracy, generating engagement effects of 0.05 points. Meanwhile, suppose Georgia, which first has measures in 2014, experiences no effect because its political culture doesn't support meaningful ballot measure campaigns. When TWFE uses post-2014 California as a control for newly-treated Georgia, it subtracts California's positive effect from Georgia's zero effect, generating a negative ``treatment effect" that represents neither state's actual experience.

This problem directly relates to Matsusaka's concept of "initiative states" versus "non-initiative states" \citep{matsusaka2004many}. Initiative states have developed what he calls the infrastructure of direct democracy—experienced campaign consultants, established signature-gathering firms, and citizens accustomed to ballot measures. Non-initiative states lack this infrastructure. TWFE's problematic comparisons essentially treat established initiative states as controls for newly-experimenting states, conflating infrastructure effects with treatment effects.

Even worse, if treatment effects vary over time---perhaps declining as novelty wears off or increasing as institutions adapt---these dynamic patterns further contaminate estimates. A state in its fifth year of treatment serving as control for a newly-treated state introduces systematic bias that cannot be signed ex ante. The resulting TWFE estimate becomes an uninterpretable weighted average of heterogeneous effects, contamination bias, and potentially negative weights.

With only Delaware as never-treated, my TWFE estimate relies heavily on these problematic comparisons. The Bacon decomposition, which I'll present when focusing on morality politics measures where it can be properly implemented, reveals that over half my identifying variation comes from such contaminated comparisons. This explains why modern estimators, when they can be implemented, often find dramatically different results from traditional TWFE.

\subsection{Extension to Objective Political Knowledge}

The results presented thus far rely on self-reported news interest, which captures subjective engagement but may not reflect actual political learning. Citizens might report following news more closely during ballot measure campaigns without genuinely understanding politics better—a form of illusory engagement. To address this concern and following recent work emphasizing the importance of distinguishing between perceived and actual political sophistication \citep{delli1996americans, prior2007post, lupia2016uninformed}, I extend the analysis to examine objective political knowledge.

The Cooperative Election Study includes a battery of factual questions about Congress that provide a more stringent test of the educative hypothesis. Respondents are asked which party controls the House and Senate, whether Republicans or Democrats need more seats for majority control, and similar factual questions that tap into basic political knowledge. I construct a Congress Knowledge index ranging from 0 to 1, representing the proportion of questions answered correctly. With a mean of 0.561 (SD = 0.387), this measure shows substantial variation and captures citizens' actual understanding of political institutions rather than their self-perceived engagement.

To ensure clean identification and comparability across outcomes, I examine both measures using the morality politics specification, which provides 19 never-treated states rather than the single control state (Delaware) available when using all ballot measures.

Importantly, the correlation between news interest and Congress knowledge is modest (r = 0.208), suggesting these measures capture distinct dimensions of civic engagement. This weak relationship aligns with prior work showing that subjective political interest often diverges from objective political knowledge \citep{prior2007post}. Citizens may feel more engaged with politics—and report following news more closely—without actually learning factual information about political institutions.

Table \ref{tab:knowledge_comparison} presents parallel specifications to my main results, using Congress Knowledge as the outcome variable. Note that these results use the morality politics specification to ensure comparability across outcomes, as this specification provides better identification with 19 never-treated states rather than just Delaware.

\begin{table}[H]
\centering
\caption{Ballot Measure Effects: News Interest vs. Congress Knowledge (Morality Politics Specification)}
\label{tab:knowledge_comparison}
\caption*{\footnotesize Notes: Two-way fixed-effects regressions estimated on CES 2006--2020 respondents using the morality-politics treatment definition. Columns compare subjective news interest (1--4) and objective Congress knowledge (0--1). Controls, fixed effects, and state-clustered standard errors align with the row labels.}
\begin{tabular}{lccc}
\hline\hline
& \multicolumn{2}{c}{Outcome Variable} \\
\cline{2-3}
& News Interest & Congress Knowledge \\
& (1-4 scale) & (0-1 scale) \\
\hline
Ballot Measure & -0.005 & 0.002 \\
& (0.008) & (0.006) \\
\hline
Individual Controls & No & No \\
State FE & Yes & Yes \\
Year FE & Yes & Yes \\
Observations & 305,086 & 310,860 \\
R² & 0.170 & 0.044 \\
\hline
With Controls & -0.004 & -0.001 \\
& (0.008) & (0.006) \\
\hline
Intensity (per measure) & -0.0003 & 0.0001 \\
& (0.0005) & (0.0002) \\
\hline\hline
\end{tabular}
\caption*{\footnotesize Notes: Two-way fixed-effects regressions estimated on CES 2006--2020 respondents using the morality-politics treatment definition (marijuana, gambling, abortion, marriage). Estimates shown are parallel across outcomes; controls include age, age$^2$, gender, education, and party indicators where indicated. Standard errors are clustered by state.}
\end{table}

The results are consistently null across all specifications for both outcomes when using the morality politics definition. For Congress Knowledge, the baseline effect is 0.002 (SE = 0.006), neither statistically nor substantively significant. This represents less than 0.5\% of the knowledge scale and approximately 0.005 standard deviations—effectively zero. News interest, which showed positive effects of 0.025 (p < 0.05) in Section 2's all-measures specification, now shows a null effect of -0.005 (SE = 0.008) when focusing on morality politics measures with proper identification. Adding individual controls or using continuous treatment intensity yields similarly null results for both outcomes. Even my most flexible specifications cannot detect meaningful gains in either knowledge or interest from exposure to morality politics ballot measures. This stark reversal for news interest—from significantly positive to null—suggests the earlier findings may have been artifacts of having only Delaware as a control state rather than genuine causal effects.

Figure \ref{fig:comparison_effects} visualizes the stark contrast between my two outcome measures:

\begin{figure}[H]
\centering
\includegraphics[width=0.9\textwidth]{\basepath/Research/policy_learning/output/figures/knowledge_vs_interest_comparison.png}
\caption{Ballot Measure Effects: Self-Reported Interest vs. Objective Knowledge}
\label{fig:comparison_effects}
\caption*{\footnotesize Notes: Points show two-way fixed-effects estimates with 95\% state-clustered confidence intervals using the morality-politics specification (marijuana, gambling, abortion, marriage) on CES 2006--2020 respondents. Outcomes are residualized on state and year fixed effects prior to plotting for visual comparability.}
\end{figure}


The divergence becomes even clearer when examining heterogeneous effects by issue area. While fiscal and morality politics measures generated significant increases in news interest in the all-measures specification presented in Section 2 (0.034 and 0.029 respectively), when I focus specifically on morality politics measures with proper identification, neither news interest nor Congress knowledge shows detectable gains. For Congress knowledge specifically, the effects are consistently null across all morality politics subtypes: marijuana (0.004, SE = 0.008), gambling (0.001, SE = 0.007), abortion (0.006, SE = 0.009), and marriage (-0.002, SE = 0.008). Similarly, news interest effects that appeared positive in Section 2 disappear when examined with the cleaner morality politics specification. This pattern holds across all issue types: even measures that theoretically should be most engaging—those touching on deeply held moral values—fail to increase either subjective engagement or objective political knowledge when properly identified. The consistency of null findings across both outcomes and all issue subtypes suggests ballot measures don't meaningfully educate citizens about politics, and that earlier positive findings likely reflected identification problems rather than genuine causal effects.

To further validate these null findings, I conduct placebo tests for the knowledge outcome. When randomly reassigning treatment status 1,000 times, the actual estimate falls at the 45th percentile of the null distribution—well within what I would expect from chance alone. This contrasts sharply with news interest, where the actual estimate lies at the 99th percentile. The placebo test results reinforce the interpretation that any effects of ballot measures are limited to subjective engagement rather than objective learning.

Several explanations could account for this divergence between subjective and objective measures. First, ballot measure campaigns might increase citizens' confidence in their political knowledge without actual learning—what \citet{kuklinski2000misinformation} term "confident ignorance." Campaign messaging often emphasizes mobilization over education, encouraging citizens to participate rather than deeply understand issues.

Second, the knowledge questions focus on congressional institutions while ballot measures concern state-level policies. However, if ballot measures truly served as "schools of democracy" as proponents claim \citep{tolbert2003virtual}, I should expect spillover effects to general political knowledge. The absence of such spillovers suggests limited educative impact.

Third, citizens might acquire narrow, issue-specific knowledge about ballot measures without broader political learning. my Congress knowledge measure cannot capture whether citizens learned about specific tax proposals or marijuana regulations. Yet from a democratic theory perspective, this narrow learning would still fall short of the comprehensive civic education that Matsusaka and others hypothesize \citep{matsusaka2020let}.

The precision of my estimates allows us to rule out even modest knowledge effects. The 95\% confidence interval excludes knowledge gains larger than 0.014 points (approximately 0.036 standard deviations). For context, this upper bound is smaller than one-fifth the knowledge gap between college and high school graduates. If ballot measures have any effect on political knowledge, it is trivially small.

These findings cast doubt on the educative value of direct democracy, at least as measured through general political knowledge. While citizens may report greater interest in news during ballot measure campaigns—possibly reflecting increased exposure to political advertising and campaign communications—this doesn't translate into improved understanding of political institutions. This pattern aligns with skeptical accounts suggesting that direct democracy overwhelms rather than educates citizens \citep{schumpeter1942capitalism, dyck2009grasping}.

The implications extend beyond ballot measures. If high-salience campaigns on concrete policy questions don't increase political knowledge, I should be skeptical that other institutional reforms can substantially enhance citizen competence. The constraints on political learning appear more binding than reformers acknowledge, suggesting limits to how much institutional design can improve democratic citizenship.

\subsection{The Parallel Trends Problem}

Even setting aside the fundamental problem of inadequate control units, the traditional analysis faces another critical challenge: the parallel trends assumption may be violated. This assumption---that treatment and control states would follow similar trajectories absent treatment---underpins all difference-in-differences designs. When violated, I cannot separate treatment effects from pre-existing differential trends.

\begin{figure}[H]
\centering
\includegraphics[width=\textwidth]{\basepath/Research/policy_learning/output/figures/all_years/parallel_trends_by_cohort.pdf}
\caption{Pre-Treatment Trends by Adoption Cohort}
\label{fig:parallel_trends_cohort}
\caption*{\footnotesize Notes: Cohorts are defined by the first year states adopt morality politics measures. Lines plot cohort-average residualized news interest (demeaned by state and year) from the CES sample, using the 2006--2020 panel. Shaded bands denote periods when the cohort is treated. The diverging trends before treatment, particularly for late adopters (2014-2016), violate the parallel trends assumption (aggregate test: p = 0.030).}
\end{figure}

\begin{figure}[H]
\centering
\includegraphics[width=\textwidth]{\basepath/Research/policy_learning/output/figures/all_years/parallel_trends_event_study.pdf}
\caption{Event Study: Pre-Treatment Trends in Relative Time}
\label{fig:parallel_trends_event}
\caption*{\footnotesize Notes: Shows average news interest relative to treatment year across all treated states. Event-time coefficients computed from a two-way fixed effects model with state and year fixed effects using the morality-politics treatment definition. The solid line reports cohort-weighted averages; the dashed line shows the mean for the 19 never-treated states. The upward trend before treatment (t $<$ 0) suggests selection into treatment rather than causal effects. Standard errors are clustered by state.}
\end{figure}

While individual pre-treatment coefficients vary by cohort, the aggregate pattern reveals concerning differential trends. States that eventually adopt morality politics measures show diverging political engagement trajectories before treatment, particularly among late adopters (2014-2016). This staggered adoption pattern—with different states receiving treatment in different years—creates complex identification challenges that traditional methods cannot address.

Formal testing confirms this concern. Interacting treatment status with a linear time trend in the pre-period yields a coefficient of -0.01339 with a p-value of 0.030, rejecting the parallel trends assumption at conventional significance levels. The violation is particularly pronounced for late-adopting states (p $<$ 0.001), suggesting heterogeneous pre-treatment dynamics across adoption cohorts. This violation fundamentally undermines causal interpretation: if treatment states werealready on different trajectories, I cannot attribute post-treatment differences to ballot measures themselves.

The parallel trends violation likely reflects reverse causality or omitted variables, a problem Matsusaka acknowledges in his discussion of endogenous institutions \citep{matsusaka2018public}. States experiencing increasing political activism may be more likely to generate ballot measures, creating the appearance of treatment effects when causality runs in the opposite direction. The cohort-specific patterns—with late adopters showing the strongest violations—suggest that the political dynamics driving ballot measure adoption evolved over my study period. Alternatively, unobserved factors like demographic changes, economic conditions, or political polarization might drive both ballot measure adoption and political engagement, generating spurious correlations my fixed effects cannot absorb.

\section{A Theory-Driven Solution: Focusing on Morality Politics}

\subsection{Theoretical Motivation}

I adopt a theoretically-motivated refinement: focusing specifically on morality politics ballot measures. This approach is grounded in substantial theoretical and empirical literature suggesting these measures should be most likely to generate civic engagement effects.

\citet{mooney2001public} defines morality politics as issues where ``at least one advocacy coalition...portrays the issue as one of morality or sin and uses moral arguments in its policy advocacy." In the American context, four issues consistently fit this definition: marijuana legalization, gambling expansion, abortion restrictions, and same-sex marriage. These issues share distinctive characteristics that should theoretically maximize civic engagement effects.

First, morality politics measures generate exceptional media coverage that penetrates beyond typical political audiences. \citet{anderson2013medical} documents how marijuana legalization campaigns receive sustained national attention, with coverage appearing not just in political sections but in lifestyle, health, and culture reporting. This broad coverage creates information-rich environments where even typically disengaged citizens encounter policy arguments. Unlike technical measures about bond authorizations or administrative procedures that receive minimal coverage, morality politics saturates media environments. This aligns with Matsusaka's argument that effective civic education requires sufficient information penetration \citep{matsusaka2020let}.

Second, moral framing makes these issues personally relevant regardless of material stake. A citizen might ignore a complex tax measure affecting only high earners, but everyone can access moral intuitions about marijuana or marriage. This personal relevance should motivate information-seeking and political discussion. \citet{smith2004ballot} finds that morality politics measures generate significantly more interpersonal discussion than other ballot measures, with citizens reporting conversations at work, church, and family gatherings. This social dimension may amplify individual-level engagement effects, creating what \citet{campbell2019affect} calls "affective mobilization."

Third, morality politics scrambles traditional partisan alignments, forcing citizens to think independently rather than simply following party cues. Libertarian Republicans might support marijuana legalization while socially conservative Democrats oppose it. Catholic Democrats might oppose abortion rights while secular Republicans support them. These cross-pressures prevent citizens from relying on simple heuristics, potentially stimulating genuine deliberation. \citet{biggers2011policy} shows that campaign advertising on morality politics measures is less partisan and more issue-focused than candidate advertising, providing substantive information rather than just party appeals. This aligns with \citet{lupia2015uninformed}'s argument that removing partisan cues can enhance rather than diminish citizen competence.

Fourth, morality politics campaigns attract substantial resources from national organizations, ensuring professional campaigns with clear messaging. National organizations like the Marijuana Policy Project, Americans for Responsible Gaming, Planned Parenthood, and National Organization for Marriage provide funding, expertise, and messaging that generate sustained campaigns. This infrastructure ensures that morality politics measures receive the kind of professional campaign treatment that might activate civic engagement. As \citet{stratmann2006effectiveness} demonstrates, campaign spending on initiatives significantly affects both outcomes and citizen knowledge.

\subsection{Empirical Advantages}

Focusing on morality politics provides empirical advantages that address the methodological problems identified earlier. Table \ref{tab:treatment_comparison} compares different treatment definitions and their implications for identification:\footnote{Appendix Table B13 provides the distribution of treatment timing across states, showing which states adopted morality politics measures in which years.}

\begin{table}[H]
\centering
\caption{Comparison of Treatment Definitions}
\label{tab:treatment_comparison}
\caption*{\footnotesize Notes: Two-way fixed effects estimates for alternative treatment definitions using the 2006--2020 CES sample. All specifications include state and year fixed effects with state-clustered standard errors (in parentheses). ``Never Treated States'' counts the units that never receive the indicated treatment. Significance: ** p$<$0.05.}
\begin{tabular}{lrll}
  \toprule
Treatment & Never Treated States & TWFE Estimate & Std. Error \\
  \midrule
Any Ballot Measure (Broad) & 1 & 0.0250** & (0.0103) \\
  Morality Politics (Focused) & 19 & 0.0288 & (0.0188) \\
  Marijuana Only & 23 & 0.0448** & (0.0184) \\
  Morality excl. Abortion & 19 & 0.0234 & (0.0184) \\
   \bottomrule
\end{tabular}
\end{table}

The improvement in identification is notable. Moving from any ballot measure to morality politics measures increases never-treated states from 1 to 19. These 19 states span diverse regions and political cultures: Vermont and Connecticut in the Northeast, Georgia and Tennessee in the South, Iowa and Indiana in the Midwest, and Wyoming in the Mountain West.\footnote{The complete list of never-treated states is: Connecticut, Delaware, Georgia, Hawaii, Idaho, Indiana, Iowa, Kansas, Kentucky, Minnesota, New Hampshire, Pennsylvania, South Carolina, Tennessee, Texas, Utah, Vermont, West Virginia, and Wyoming.} This geographic and political diversity provides a much more credible control group than Delaware alone.

The expanded control group enables implementation of modern difference-in-differences methods that avoid TWFE's problematic comparisons. With 19 never-treated states, I can implement Sun-Abraham and Callaway-Sant'Anna estimators that compare treated units only to never-treated units, avoiding the contamination bias that plagues traditional estimates. The diversity of control states also makes parallel trends more plausible---it's more believable that treatment states would follow similar trajectories to a diverse set of 19 states than to Delaware alone.

This refinement exemplifies what \citet{sekhon2009opiates} calls "theory-driven empiricism"—using substantive knowledge to improve identification rather than letting data availability drive research questions. By focusing on measures where theory predicts the strongest effects, I conduct a more decisive test of the educative hypothesis while simultaneously improving econometric identification.

Figure \ref{fig:morality_distribution} shows the distribution of morality politics measures across types and years:

\begin{figure}[H]
\centering
\begin{tikzpicture}
\begin{axis}[
    ybar,
    width=12cm,
    height=6cm,
    xlabel={Type of Morality Politics Measure},
    ylabel={Number of Measures (2006-2020)},
    ymin=0,
    symbolic x coords={Marijuana,Gambling,Abortion,Marriage},
    xtick=data,
    nodes near coords,
    nodes near coords align={vertical},
    bar width=0.7cm,
    ]
\addplot coordinates {(Marijuana,187) (Gambling,142) (Abortion,68) (Marriage,40)};
\end{axis}
\end{tikzpicture}
\caption{Distribution of Morality Politics Measures by Type}
\label{fig:morality_distribution}
\caption*{\footnotesize Notes: Counts from the harmonized ballot-measure dataset (2006--2020) restricted to morality politics topics: marijuana legalization, gambling expansion, abortion regulation, and same-sex marriage. Totals reflect distinct statewide ballot measures irrespective of success.}
\end{figure}

The distribution reveals that marijuana and gambling measures dominate, together accounting for 75\% of morality politics measures. This concentration has advantages and disadvantages. On one hand, it provides substantial variation for identifying marijuana and gambling effects specifically. On the other hand, it means my aggregate estimates are driven primarily by these two issue areas. The relative scarcity of abortion and marriage measures (together just 25\%) means I have less power to detect their specific effects and must be cautious about generalizing from marijuana and gambling to all morality politics.

This distribution reflects what \citet{haider2006framing} identifies as the "issue attention cycle" in morality politics—marijuana and gambling represent contemporary moral debates while marriage equality was largely resolved during my study period. The dominance of these "libertarian" morality issues versus traditional "conservative" morality issues may explain different patterns of engagement across citizen subgroups.

\section{Modern DiD Estimators Mostly Reveal Null Findings}

\subsection{Implementing Modern Difference-in-Differences}

Adding 19 never-treated states provides a credible control group and allows for the implementation of modern estimators specifically designed to handle the challenges of staggered treatment adoption. These methods, developed in response to growing recognition of TWFE's limitations, avoid problematic comparisons and allow for heterogeneous treatment effects across cohorts and time.

The \textbf{Sun-Abraham (2021)} estimator addresses TWFE's problems by estimating separate treatment effects for each adoption cohort relative to never-treated units, then aggregating these cohort-specific effects. Formally, it estimates:

$$ Y_{ist} = \alpha_{is} + \delta_t + \sum_{g \neq 0} \sum_{\tau} \beta_{g,\tau} \cdot \mathbf{1}[G_i = g] \cdot \mathbf{1}[t - G_i = \tau] + \epsilon_{ist} $$

where $G_i$ indicates unit $i$'s treatment cohort (with $G_i = 0$ for never-treated), and $\tau$ represents relative time to treatment. By interacting cohort indicators with relative time, Sun-Abraham allows each cohort to have its own dynamic treatment path while using only never-treated units for identification.\footnote{Appendix Table B2 presents the full set of relative-time coefficients from the Sun-Abraham estimator, showing how treatment effects evolve before and after adoption.} The approach is particularly valuable when treatment effects vary across adoption cohorts, as I might expect if early adopters differ systematically from late adopters—exactly what Matsusaka's framework predicts given different levels of direct democracy infrastructure \citep{matsusaka2018public}.

The \textbf{Callaway-Sant'Anna (2021)} estimator takes a different but complementary approach, using doubly-robust estimation that combines outcome regression and propensity score methods. For each cohort-time pair, it estimates:

$$ ATT(g,t) = E[Y_t(g) - Y_t(\infty) | G = g] $$

where $Y_t(g)$ represents potential outcomes when first treated at time $g$, and $Y_t(\infty)$ represents never-treated potential outcomes. Identification comes from comparing treated units only to never-treated or not-yet-treated units, completely avoiding already-treated controls.\footnote{Appendix Table B3 provides group-time average treatment effects before aggregation, revealing heterogeneity across treatment cohorts.} The doubly-robust approach provides consistency if either the outcome model or propensity score model is correctly specified, offering insurance against misspecification.

The \textbf{Bacon decomposition} provides diagnostic insight by decomposing the traditional TWFE estimate into its component comparisons. Following \citet{bacon2021difference}, it calculates how much I ight TWFE places on different types of comparisons:

\begin{itemize}
    \item \textbf{Clean comparisons}: Treated vs. never-treated units
    \item \textbf{Problematic comparisons}: Early vs. late treated, or within-cohort comparisons that may receive negative weights
\end{itemize}

This decomposition reveals whether TWFE's estimate is driven by clean identification or contaminated comparisons, directly testing whether the problems \citet{de2020two} identifies theoretically actually matter empirically.

\subsection{The Divergence: Traditional vs. Modern Estimates}

Table \ref{tab:main_comparison} presents the key comparison between traditional and modern estimators for the morality politics specification:

\begin{table}[H]
\centering
\caption{The Methodological Divergence: Traditional vs. Modern Estimators}
\label{tab:main_comparison}
\caption*{\footnotesize Notes: Estimates computed on the morality-politics sample (480{,}383 CES respondents, 2006--2020). Traditional rows use two-way fixed effects with state and year fixed effects; modern rows implement Sun--Abraham (2021) and Callaway--Sant'Anna (2021) estimators relying on the 19 never-treated states. Standard errors (in parentheses) are clustered by state.}
\begin{tabular}{lccc}
\hline\hline
Method & Estimate & Std. Error & Interpretation \\
\hline
\multicolumn{4}{l}{\textit{Traditional Methods:}} \\
Simple DiD & -0.0185* & (0.0100) & Negative without FE \\
Two-way FE & 0.0288 & (0.0188) & Positive, marginally sig. \\
With State Trends & 0.0269 & (0.0203) & Slightly attenuated \\
\hline
\multicolumn{4}{l}{\textit{Modern Methods:}} \\
Sun-Abraham (2021) & -0.0001 & --- & Essentially zero \\
Callaway-Sant'Anna (2021) & -0.0014 & (0.0133) & Null effect \\
\hline\hline
\end{tabular}
\end{table}

Traditional TWFE suggests morality politics measures increase engagement by 0.0288 points (approximately 3\% of a standard deviation), marginally significant with p $\approx$ 0.13. While not overwhelming evidence, this would support theoretical predictions about morality politics' engagement effects. Adding state-specific linear trends slightly attenuates the estimate to 0.0269, suggesting modest sensitivity to trend specifications.

The modern estimators tell a completely different story. Sun-Abraham's estimate of -0.0001 is not merely statistically insignificant but substantively negligible---one ten-thousandth of a point on a four-point scale. This isn't a matter of insufficient power; the estimate is essentially zero. Callaway-Sant'Anna similarly finds -0.0014 with a standard error of 0.0133, yielding a 95\% confidence interval of [-0.027, 0.025]. Even the upper bound suggests effects smaller than traditional TWFE's point estimate.

The modern estimators suggest that what TWFE interprets as treatment effects are actually artifacts of problematic comparisons and heterogeneous treatment effects. When I properly account for staggered adoption using methods designed for this purpose, the apparent effects vanish entirely. This stark divergence exemplifies what \citet{roth2022s} calls the "credibility revolution" in difference-in-differences—many published findings using TWFE may not survive scrutiny with appropriate methods.

\subsection{Diagnosing the Problem}

Three interconnected issues explain this dramatic divergence between traditional and modern estimates:

\textbf{1. Parallel Trends Violation}

The parallel trends assumption---crucial for any difference-in-differences design---requires that treatment and control units would follow parallel trajectories absent treatment. Figure \ref{fig:event_study_morality} presents an event study examining this assumption:

\begin{figure}[H]
\centering
\includegraphics[width=\textwidth]{\basepath/Research/policy_learning/output/figures/all_years/event_study_morality.pdf}
\caption{Event Study: Morality Politics Measures}
\label{fig:event_study_morality}
\caption*{\footnotesize Notes: Estimated using the Sun--Abraham cohort-specific event-study specification with 19 never-treated states as controls. Points plot average treatment effects at each relative period; vertical lines show 95\% confidence intervals with state-clustered standard errors. Period 0 corresponds to the election containing the first morality-politics measure.}
\end{figure}

The pre-treatment coefficients (periods -3 to -2) trend upward approaching treatment, suggesting that states receiving morality politics measures werealready experiencing increasing political engagement before treatment. While individual coefficients aren't statistically significant, formal testing strongly rejects parallel trends. Interacting treatment group with a linear time trend in the pre-period yields a coefficient of 0.0042 with p = 0.016, providing strong evidence of differential pre-trends.\footnote{Appendix Table B4 presents state-specific parallel trends tests, identifying Alaska and Washington as exhibiting the most severe violations.}

State-specific analysis reveals that Alaska and Washington exhibit the most severe violations, with pre-treatment trends significantly different from never-treated states. These states experienced substantial political changes during my study period unrelated to ballot measures---Alaska's political corruption scandals and subsequent reforms, Washington's tech boom and demographic transformation. These confounding trends contaminate estimates if not properly controlled.

The parallel trends violation fundamentally undermines causal interpretation. If treatment states werealready on different trajectories, I cannot attribute post-treatment differences to morality politics measures. Modern estimators partially address this through their comparison structure, but the violation suggests even their null estimates may be optimistic. This echoes \citet{kahn2020promise}'s warning that parallel trends violations are endemic in institutional studies where treatment reflects endogenous political processes.

\textbf{2. Problematic Comparisons in TWFE}

The Bacon decomposition reveals exactly what comparisons drive the TWFE estimate. Table \ref{tab:bacon_decomp} presents the breakdown:\footnote{Appendix Table B5 provides the complete Bacon decomposition with weights and average effects for each comparison type.}

\begin{table}[H]
\centering
\caption{Bacon Decomposition of TWFE Estimate}
\label{tab:bacon_decomp}
\caption*{\footnotesize Notes: Decomposition of the two-way fixed effects estimate for the morality-politics treatment (2006--2020 CES sample). Weights indicate the share of the TWFE coefficient contributed by each comparison type following Goodman-Bacon (2021). Classifications label whether comparisons rely on clean never-treated controls or potentially contaminated treated units.}
\begin{tabular}{lcc}
\hline\hline
Comparison Type & I ight & Classification \\
\hline
Never vs. Timing Groups & 42\% & Clean \\
Earlier vs. Later Treated & 38\% & Problematic \\
Within Cohort & 20\% & Problematic \\
\hline
Total Problematic & 58\% & \\
\hline\hline
\end{tabular}
\end{table}

This decomposition is illuminating and concerning. Only 42\% of TWFE's identifying variation comes from clean comparisons between treated and never-treated states---the comparisons I actually want. The remaining 58\% comes from problematic comparisons that can severely bias estimates.

The 38\% I ight on ``earlier vs. later treated" comparisons is particularly troubling. These comparisons use states like California and Colorado, which adopted morality politics measures early and frequently, as controls for states like Ohio or Florida that adopted later. But California in 2016, having had marijuana measures for a decade, is not a valid control for Ohio getting its first marijuana measure. California's post-treatment engagement reflects both its long history with direct democracy and contemporary treatment, contaminating the comparison. This directly relates to Matsusaka's distinction between states with developed versus nascent direct democracy infrastructure \citep{matsusaka2018public}.

The 20\% I ight on within-cohort comparisons can be even more problematic, as these may receive negative weights in the presence of treatment effect heterogeneity. \citet{de2020two} shows that such negative weights can reverse the sign of estimates, making harmful treatments appear beneficial or vice versa. With over half my identifying variation coming from contaminated comparisons, it's unsurprising that TWFE and modern estimators diverge dramatically.

\textbf{3. Heterogeneous Treatment Effects}

The assumption of homogeneous treatment effects, implicit in TWFE, appears severely violated. Table \ref{tab:timing_heterogeneity} examines how effects vary by adoption timing:

\begin{table}[H]
\centering
\caption{Heterogeneous Effects by Treatment Adoption Timing}
\label{tab:timing_heterogeneity}
\caption*{\footnotesize Notes: Two-way fixed effects estimates with state and year fixed effects using the morality-politics sample (480{,}383 CES respondents). Cohorts are split by first treatment year. Standard errors (in parentheses) are clustered by state. Opposite-signed estimates highlight the heterogeneous timing effects that bias pooled TWFE results. Significance: ** p$<$0.05.}
\begin{tabular}{lcc}
\hline\hline
Cohort & Effect & Std. Error \\
\hline
Early Adopters ($\leq$2010) & 0.0551** & (0.0241) \\
Late Adopters (>2010) & -0.0175 & (0.0192) \\
\hline
\multicolumn{3}{l}{\footnotesize Note: Opposite signs explain TWFE bias} \\
\hline\hline
\end{tabular}
\end{table}

This heterogeneity is striking and explains much of the TWFE bias. Early adopters---states that implemented morality politics measures before 2010---show positive effects of 0.0551, statistically significant and economically meaningful. These states, including California, Colorado, and Oregon, have long histories with direct democracy and developed infrastructure for ballot measure campaigns. Their citizens may be socialized to expect and engage with ballot measures, confirming Matsusaka's argument about the importance of institutional learning \citep{matsusaka2004many}.

Late adopters show negative effects of -0.0175, though not statistically significant. These states adopted morality politics measures during a different political era, often through very different processes. Many late-adopting states had measures placed on ballots by legislatures rather than citizen initiatives, potentially generating less organic engagement. The political context also differed, with late adoption occurring during heightened polarization that may have transformed ballot measures from deliberative opportunities into partisan battlegrounds, contradicting the deliberation mechanism central to the educative hypothesis.

TWFE pools these opposite-signed effects into a single parameter, producing a meaningless average. Worse, when using early adopters as controls for late adopters, TWFE subtracts positive effects from negative ones, amplifying apparent treatment effects through pure artifact. Modern estimators avoid this by maintaining separate cohort effects and using only clean comparisons, revealing that the pooled effect is essentially zero.

\subsection{Topic-Specific Analysis: Unpacking the Components}

To better understand what drives results, I examine effects separately by type of morality politics measure. Table \ref{tab:morality_topics} presents topic-specific estimates:\footnote{Appendix Table B8 provides complete regression results for each topic-specific model with full specifications.}

\begin{table}[H]
\centering
\caption{Effects by Morality Politics Topic}
\label{tab:morality_topics}
\caption*{\footnotesize Notes: Topic-specific two-way fixed effects estimates using the morality-politics treatment definition with state and year fixed effects. ``N Treated'' denotes the number of CES respondents exposed to each topic post-election. Standard errors (in parentheses) are clustered by state. Significance: * p$<$0.10, ** p$<$0.05, *** p$<$0.01.}
\begin{tabular}{lrrrl}
\hline\hline
Measure Type & Coefficient & Std. Error & N Treated & Assessment \\
\hline
Abortion & 0.2234*** & (0.0602) & 359 & Implausible outlier \\
Marijuana & 0.0448** & (0.0184) & 3,943 & Only robust effect \\
Gambling & 0.0207 & (0.0269) & 6,408 & Null \\
Marriage & 0.0183 & (0.0221) & 1,287 & Null \\
\hline\hline
\end{tabular}
\end{table}

These disaggregated results reveal serious problems with the aggregate analysis. The abortion coefficient of 0.2234 is implausibly large---suggesting that a single abortion measure increases news interest by nearly a quarter point, or 6\% of the entire scale range. This would make abortion measures approximately 10 times more powerful than intensive canvassing efforts. Such massive effects strain credulity and likely reflect unobserved confounders or specification problems rather than genuine causal effects.

The abortion result's implausibility becomes clearer when considering that it's identified from just 359 treated observations across a handful of states and years. These measures often coincided with major political battles over abortion rights, including legislative attempts to challenge Roe v. Wade. The apparent ``effect" likely captures broader political mobilization around abortion rights rather than the ballot measure itself. When modern estimators properly isolate the ballot measure effect from concurrent political activity, the effect disappears.

Marijuana measures show the only potentially robust effect at 0.0448 (p $<$ 0.05), identified from nearly 4,000 treated observations. This effect size, while modest, appears more plausible and is precisely estimated. Several factors might explain why marijuana measures uniquely generate engagement. First, marijuana legalization represented genuinely new policy terrain during my study period, potentially generating more learning than familiar issues. This aligns with \citet{nicholson2003agenda}'s argument that novel issues create more engagement than recurring topics. Second, marijuana campaigns attracted unusual coalitions crossing traditional partisan lines, potentially engaging citizens who typically follow party cues. Third, the concrete nature of marijuana policy---Will dispensaries open? Will people be arrested?---may be more engaging than abstract moral arguments.

However, even the marijuana effect proves fragile under scrutiny.\footnote{Appendix Table B7 presents leave-one-state-out results showing sensitivity to specific states, with coefficients ranging from 0.036 to 0.052.} Leave-one-state-out analysis reveals coefficients ranging from 0.036 to 0.052 depending on which state is excluded, suggesting that no single state drives results but also that results aren't robust to sample composition. Adding state-specific trends reduces the coefficient to 0.0301 and eliminates statistical significance. When examined with modern estimators focusing only on marijuana measures, the effect again vanishes, suggesting even this apparently robust finding reflects TWFE bias rather than genuine effects.

Gambling and marriage measures show null effects throughout, despite theoretical predictions that these moral issues should engage citizens. The null effects for marriage equality measures are particularly surprising given the intense political battles over this issue during my study period. This might reflect that marriage equality opinions werealready crystallized through national debate, leaving little room for ballot measures to generate additional engagement. This pattern supports \citet{carmines1989issue}'s theory that "easy" issues with clear moral dimensions may not require additional information-seeking.

\section{Extensions: Heterogeneity, Mechanisms, and Robustness}

\subsection{Heterogeneous Effects Across Citizens}

Even with null average effects, examining heterogeneity remains important for understanding whether specific subgroups might benefit from ballot measures. Political engagement interventions often show concentrated effects among particular populations, and identifying these patterns could inform targeted democratic reforms.\footnote{Appendix Table B6 presents complete interaction models examining heterogeneous effects across all demographic subgroups.} This analysis directly engages with \citet{matsusaka2005direct}'s argument that direct democracy may differentially benefit different groups.

Figure \ref{fig:heterogeneity_combined} examines effects across demographic groups:

\begin{figure}[H]
\centering
\includegraphics[width=\textwidth]{\basepath/Research/policy_learning/output/figures/all_years/heterogeneity_demographics_twfe.pdf}
\caption{Heterogeneous Treatment Effects by Demographics}
\label{fig:heterogeneity_combined}
\caption*{\footnotesize Notes: Two-way fixed-effects estimates by demographic subgroup using the all-ballot-measures specification and CES 2006--2020 sample. Horizontal lines denote 95\% state-clustered confidence intervals; filled markers indicate p<0.05. Modern staggered DiD estimators cannot be implemented because untreated comparison units are exhausted.}
\end{figure}

Several patterns emerge from this TWFE analysis, though these should be interpreted with extreme caution given that modern DiD methods cannot identify subgroup effects when properly accounting for staggered treatment. Of the eight demographic groups examined, only two show statistically significant effects.

First, education shows a modest difference, with non-college-educated citizens displaying a statistically significant effect of 0.031 (p<0.05) compared to an insignificant 0.012 for college-educated citizens. While this might suggest ballot measures engage less-educated citizens more than their college-educated counterparts, the overlapping confidence intervals and failure of modern methods to identify these differences caution against strong interpretation. This pattern contradicts theories suggesting complex policy decisions require educational resources to generate engagement, but aligns with \citet{lupia2015uninformed}'s argument that formal education may be less important for political competence than commonly assumed.

Second, partisan differences appear but lack statistical significance. Republicans show the largest point estimate (0.049) compared to Democrats (0.017) or Independents (0.025), but none reach conventional significance levels. The Republican effect might reflect that morality politics measures during my study period—particularly marijuana legalization—challenged conservative orthodoxy more than progressive positions. However, without statistical significance and given the estimation challenges with modern methods, these differences likely reflect TWFE bias rather than genuine partisan heterogeneity.

Third, age heterogeneity shows the clearest pattern, with middle-aged citizens (35-64) displaying the only highly significant effect at 0.061 (p<0.01). Younger citizens show essentially null effects (0.002), while older citizens show negative but insignificant effects (-0.024). The concentration among middle-aged citizens might reflect life-stage factors—raising children, paying mortgages, peak career engagement—that make policy outcomes feel most consequential. This aligns with \citet{wolfinger1980votes}'s life-cycle theory of political participation. However, this is the group most likely to be affected by TWFE's negative I ighting problems in staggered adoption settings.

Critically, Callaway-Sant'Anna estimation fails for all demographic subgroups due to insufficient identifying variation once negative I ighting is eliminated. This failure is not a technical limitation but reveals that the apparent heterogeneity shown here likely reflects TWFE's known biases—using already-treated units as controls, I ighting some comparisons negatively, and conflating treatment effect heterogeneity with treatment timing heterogeneity. With the overall Callaway-Sant'Anna estimate for the full sample at -0.0014, these TWFE patterns almost certainly represent methodological artifacts rather than true differential policy learning effects.

\subsection{Political Awareness and Information Processing}

A crucial theoretical question concerns how political sophistication moderates ballot measure effects. The information overload hypothesis suggests low-awareness citizens might be overwhelmed by complex policy choices, while the activation hypothesis \citep{smith2004ballot} suggests ballot measures might uniquely engage typically disconnected citizens. Following the morality politics literature \citep{mooney2001public}, certain high-salience issues---particularly those with moral dimensions---may be especially effective at mobilizing low-awareness citizens who can form opinions based on values rather than detailed policy knowledge.

Figure \ref{fig:awareness_test} examines this relationship:

\begin{figure}[H]
\centering
\includegraphics[width=\textwidth]{\basepath/Research/policy_learning/output/extensions/morality_politics/awareness_effects_morality.pdf}
\caption{Effects by Political Awareness Level}
\label{fig:awareness_test}
\caption*{\footnotesize Notes: Estimates obtained from two-way fixed effects models run separately by terciles of the CES political awareness index (morality-politics treatment, 2006--2020). Bars show point estimates; whiskers denote 95\% confidence intervals with state-clustered standard errors. Awareness terciles are defined using pre-election baseline responses.}
\end{figure}

The pattern challenges both theoretical perspectives. Low-awareness citizens show positive effects (0.087) that exceed those for medium (0.005) or high-awareness (-0.070) citizens. This contradicts information overload, as I 'd expect negative or null effects for citizens lacking political resources. But it also challenges simple activation stories, as the effect isn't monotonically decreasing with awareness.

One interpretation is that ballot measures provide accessible entry points for political engagement that don't require extensive political knowledge. A low-awareness citizen can form opinions about marijuana or gambling based on personal experience and values, while understanding candidate positions requires following political news. Ballot measures might thus serve as what \citet{zaller1992nature} calls "easy issues" that build efficacy and interest. This aligns with Matsusaka's argument that direct democracy can engage citizens who find representative politics alienating \citep{matsusaka2020let}.

Alternatively, high-awareness citizens might already be maximally engaged, creating ceiling effects that prevent additional activation. These citizens already follow political news closely; ballot measures simply redirect rather than increase their attention. Middle-awareness citizens might be in a ``valley" where they know enough to recognize ballot measures' complexity but lack confidence to engage deeply. This non-monotonic pattern challenges linear models of political sophistication.

Of course, this paper is not set up to adjudicate among these---or any other---competing explanations, and I leave this exercise for future research. The pattern does suggest, however, that if ballot measure effects exist, they may operate through different mechanisms than traditional political mobilization.

\subsection{Temporal Dynamics and Persistence}

If ballot measures generate learning and civic skills, effects should persist beyond immediate campaigns. Table \ref{tab:persistence} examines temporal dynamics:

\begin{table}[H]
\centering
\caption{Persistence of Effects}
\label{tab:persistence}
\caption*{\footnotesize Notes: Two-way fixed effects regression including contemporaneous and lagged indicators for morality-politics treatments, estimated on 480{,}383 CES respondents. All models include state and year fixed effects with state-clustered standard errors (in parentheses). The cumulative effect sums the period-specific point estimates. Significance: ** p$<$0.05.}
\begin{tabular}{lcc}
\hline\hline
Time Period & Effect & Std. Error \\
\hline
Contemporaneous (t) & 0.0464** & (0.0193) \\
One-year lag (t-1) & -0.0029 & (0.0068) \\
Two-year lag (t-2) & -0.0041 & (0.0186) \\
\hline
Cumulative effect & 0.0394 & (0.0277) \\
\hline\hline
\end{tabular}
\end{table}

The rapid decay of effects challenges civic republican theories of democratic participation. The contemporaneous effect of 0.0464 represents campaign-period mobilization, but this completely dissipates within one year. By the second year, point estimates turn negative (though insignificant), suggesting no lasting impact on political engagement.

This temporal pattern suggests ballot measures generate temporary attention rather than lasting transformation. Citizens engage during campaigns---attending to news, discussing issues with friends, perhaps researching positions---but return to baseline engagement once elections pass. The lack of persistence undermines arguments that direct democracy builds civic capacity or democratic skills that transfer across contexts. This finding contradicts \citet{smith2004educated}'s claim about lasting educative effects but aligns with \citet{nicholson2003agenda}'s evidence of temporary mobilization.

Several mechanisms might explain this rapid decay. First, the information acquired during campaigns might be issue-specific rather than transferable. Learning about marijuana policy doesn't necessarily help citizens understand education finance or healthcare reform, limiting spillover effects. Second, the motivations for engagement might be instrumental rather than intrinsic. Citizens engage to influence specific outcomes, not because ballot measures cultivate general political interest. Third, modern media environments might prevent sustained attention. With constant news cycles and social media, yesterday's ballot measure quickly becomes ancient history.

The lack of persistence also has methodological implications. If effects are purely contemporaneous, identification comes entirely from comparing election to non-election periods within states. This makes my estimates vulnerable to any unobserved factors varying with electoral cycles, potentially explaining why modern estimators find null effects when properly accounting for temporal patterns. This echoes \citet{erikson2002timeline}'s warning about conflating electoral cycles with policy effects.

\subsection{Comprehensive Robustness Analysis}

Given the dramatic divergence between traditional and modern estimators, extensive robustness checking is essential. Table \ref{tab:robustness_full} presents comprehensive tests:\footnote{Appendix Table B12 provides complete robustness results including wild cluster bootstrap p-values for all specifications.}

\begin{table}[H]
\centering
\caption{Comprehensive Robustness Checks}
\label{tab:robustness_full}
\caption*{\footnotesize Notes: Each row re-estimates the morality-politics specification on the 2006--2020 CES sample. Traditional variants retain state and year fixed effects with state-clustered standard errors (reported in parentheses unless otherwise noted). Modern estimators implement Sun--Abraham (2021), Callaway--Sant'Anna (2021), and De Chaisemartin--D'Haultfoeuille (2020) approaches using the 19 never-treated states. P-values for the wild cluster bootstrap use 999 Webb-weighted replications. Significance: * p$<$0.10, ** p$<$0.05, *** p$<$0.01.}
\small
\begin{tabular}{lcc}
\hline\hline
Specification & Coefficient & (SE) \\
\hline
\multicolumn{3}{l}{\textit{Traditional TWFE Variants:}} \\
Baseline TWFE & 0.0288 & (0.0188) \\
State-specific trends & 0.0269 & (0.0203) \\
Drop high-intensity states & 0.0166 & (0.0195) \\
Exclude 2020 (COVID) & 0.0292 & (0.0188) \\
Exclude states w/ bad pre-trends & 0.0298 & (0.0207) \\
Wild cluster bootstrap & 0.0288 & [p=0.082] \\
State-year clustering & 0.0288 & (0.0214) \\
Robust standard errors & 0.0288*** & (0.0095) \\
\hline
\multicolumn{3}{l}{\textit{Modern Estimators:}} \\
Sun-Abraham (2021) & -0.0001 & --- \\
Callaway-Sant'Anna (2021) & -0.0014 & (0.0133) \\
De Chaisemartin-D'Haultfoeuille & 0.0008 & (0.0156) \\
\hline
\multicolumn{3}{l}{\textit{Alternative Treatments:}} \\
Marijuana only & 0.0448** & (0.0184) \\
Morality excl. abortion & 0.0234 & (0.0184) \\
High-salience measures only & 0.0195 & (0.0211) \\
\hline\hline
\end{tabular}
\end{table}

The traditional TWFE variants show modest sensitivity but maintain positive point estimates throughout. Adding state-specific trends slightly attenuates effects, suggesting some confounding from differential trends. Dropping high-intensity states (California, Oregon, Colorado) reduces the coefficient by 42\%, indicating these states contribute disproportionately to positive findings. This sensitivity to specific states echoes Matsusaka's point about heterogeneous direct democracy cultures across states \citep{matsusaka2018public}. The wild cluster bootstrap, which addresses potential problems from few clusters, yields p = 0.082, suggesting marginal significance at best.

Excluding 2020 barely changes results, suggesting COVID-19 didn't drive findings. Excluding states with the worst pre-trends (Alaska, Washington) actually increases the coefficient slightly, indicating that parallel trends violations, while concerning, don't fully explain positive TWFE estimates. Alternative standard error calculations change precision but not point estimates, as expected.

The modern estimators consistently find null effects across different implementations. Beyond Sun-Abraham and Callaway-Sant'Anna, I implement \citet{de2020two}'s estimator, which also finds essentially zero effect (0.0008). The consistency across three distinct methodological approaches, each handling staggered adoption differently, strongly suggests that morality politics measures don't meaningfully affect political engagement. This robustness to estimator choice strengthens confidence in the null finding.

Alternative treatment definitions reveal that results depend critically on what measures I include. Marijuana-only shows the strongest effects, but as discussed earlier, these prove fragile under scrutiny. Excluding abortion measures---which showed implausibly large effects---reduces the aggregate estimate, confirming that outlier measures inflate results. Focusing only on high-salience measures (those with close votes or high campaign spending) yields null effects, contradicting the hypothesis that competitive campaigns drive engagement. This last finding is particularly damaging to the educative hypothesis, as high-salience campaigns should theoretically maximize information transmission and deliberation.

\section{Conclusion}

This paper examined whether ballot measures enhance civic engagement, revealing how methodological choices fundamentally shape conclusions in political science. Initial analysis using traditional two-way fixed effects found encouraging results: ballot measure exposure increases political engagement by 0.025 points (p < 0.05), with effects scaling by intensity and passing placebo tests. These findings would justify expanding direct democracy as a civic renewal tool, validating Matsusaka's optimistic account of government by referendum.

However, this narrative crumbled when confronting two methodological realities. First, defining treatment as any ballot measure left only Delaware as never-treated control, forcing problematic comparisons between early and late adopters. Second, traditional TWFE severely biases estimates when treatment effects are heterogeneous—as with most political institutions. The theoretically-motivated focus on morality politics measures (marijuana, gambling, abortion, marriage) expanded never-treated states from one to nineteen, enabling modern estimators that avoid TWFE's contaminated comparisons.

Modern estimators designed for staggered adoption found unambiguous null effects: Sun-Abraham yields -0.0001 and Callaway-Sant'Anna yields -0.0014. Combined with parallel trends violations (p = 0.016) and Bacon decomposition showing 58\% of TWFE's I ight from problematic comparisons, the evidence strongly suggests morality politics measures don't meaningfully affect engagement.

These findings make three contributions. For democratic theory, they challenge assumptions about citizen capacity—even high-salience moral issues with extensive campaigns fail to sustainably engage citizens, suggesting hard constraints on political attention that reforms cannot overcome. If citizens don't learn from direct votes on concrete policies, expecting them to engage meaningfully with complex governance seems unrealistic.

For methodology, the dramatic divergence between TWFE and modern estimators demonstrates how traditional methods with staggered adoption can completely reverse conclusions. Given staggered adoption's ubiquity in political science, many published findings may need reconsideration. The literature on institutional effects—from term limits to campaign finance reform to electoral systems—likely contains numerous false positives generated by TWFE's biases.

For policy, results counsel modest expectations: direct democracy serves important functions—expressing citizen preferences on specific issues, providing safety valves for frustrated majorities, enabling policy experimentation—without requiring civic education effects. Rather than expecting ballot measures to transform citizen competence, I should focus on making periodic participation as informed and accessible as possible.

The journey from promise to null provides valuable lessons. As political science emphasizes causal inference, I must ensure methods match inferential goals. The tools exist to handle complex treatment patterns common in political institutions; using them is essential for scientific progress. Future research should apply modern estimators to reexamine other institutional effects and explore whether different contexts might generate genuine engagement. While this study found null effects for morality politics measures, other designs might prove more effective. Only by applying appropriate methods can I build reliable knowledge about how institutions shape political behavior.

\clearpage

\mychapter{Tuition-Free College and Political Participation: Evidence from California's Cal Grant Program}
\vspace{0.15in}
\noindent Published in the \textit{American Economic Journal: Economic Policy} Joint with Daniel Firoozi
\vspace{0.15in}
\label{chap:calgrant}

\providecommand{\sym}[1]{\ifmmode^{#1}\else\(^{#1}\)\fi}

\section{Introduction} \label{sec:introduction}

Identifying the social returns to education expenditures requires estimates of private labor market returns as well as externalities accruing to other parties. Extant research explores the effects of education spending on health, crime, and innovation but overlooks civic externalities despite their historical role in economics.\footnote{For example, recent empirical work shows that education expenditures can have large impacts on criminal activity \citep{bell_why_2022,anders_effect_2023,gray-lobe_long-term_2023,baron_public_2024}, estimates the impact of costly education policies on mortality and long-run health \citep{clark_effect_2013,meghir_education_2018,lundborg_long-term_2022}, and finds that education spending can foster innovation \citep{toivanen_education_2016,andrews_how_2023,babina_cutting_2023}.} Economists like Adam Smith and Milton Friedman justified public spending on primary and secondary schools by asserting that they produce informed and politically active citizens \citep{friedman_capitalism_2020,frame_adam_2022}. Policymakers enacted costly compulsory and universal schooling laws on the same basis, explicitly citing returns to civic participation rather than earnings \citep{mann_republic_1957}. Omitting civic externalities from the social returns to education spending therefore risks ignoring a key external benefit and narrowing the scope of public finance.

A well-documented positive association between educational attainment and political participation exists in most contexts, with educated people more likely to vote, follow the news, and hold public office \citep{gethin_brahmin_2021,chetty_diversifying_2023}. Existing work, however, is mixed on the civic returns to education expenditures, especially later in life.\footnote{Papers on primary schools generally find positive effects \citep{dee_are_2004,sondheimer_using_2010,wantchekon_education_2015} but evidence from secondary schools is inconsistent \citep{milligan_does_2004,tenn_effect_2007,marshall_coarsening_2016,marshall_anti-democrat_2019,cohodes_why_2021,willeck_education_2022, bommel_revisiting_2023}.} In particular, we know little about whether large-scale public education expenditures like college tuition subsidies impact civic engagement, despite 235 billion dollars of financial aid dispersed annually in the US alone \citep{college_board_trends_2023}. Estimates of the civic externalities of higher education spending are limited by the dearth of instruments that allow for a credible identification strategy.\footnote{Most papers match students on observables to estimate the civic externalities of higher educational attainment rather than higher education expenditures, finding conflicting results \citep{kam_reconsidering_2008,henderson_who_2011,mayer_does_2011,willeck_education_2022,scott_does_2022,bell_which_2024}. Three studies use instruments for college enrollment -- specifically distance to college and conscription -- and again arrive at different conclusions \citep{dee_are_2004,berinsky_education_2011,doyle_does_2017}.} A related literature that estimates the returns to tuition subsidies often focuses squarely on graduation rates and earnings, overlooking plausibly large civic returns that accrue to conditional cash transfers in many contexts \citep{manacorda_government_2011}. Identifying the civic externalities of higher education spending remains critical to understanding its social returns, especially given its often modest marginal treatment effects on labor market outcomes.


In this paper, we use data on 16.4 million Free Applications for Federal Student Aid (FAFSAs) and a regression discontinuity (RD) design to identify the micro-level and macro-level civic returns to tuition subsidies. We use strict eligibility rules for the United States' largest tuition-free 4-year college program, the Cal Grant, to estimate impacts on several measures of voter registration and turnout \citep{kane_quasi-experimental_2003,bettinger_long-run_2019,dickler_5_2022}. The setting is ideal because (1) extant evidence on the program's labor market returns are critical for understanding its civic externalities, (2) California is the largest democratic setting where most voters (roughly 3/4ths) self-report party preferences in administrative data, and (3) California has the highest college student retention rate in the United States \citep{hendren_unified_2020,van_dam_states_2022,firoozi_effect_2023}.

We find that each Cal Grant awarded by the California Student Aid Commission (CSAC) raises a student's probability of casting a ballot in 2020 by roughly 10 percentage points relative to a baseline turnout of 56 percent. This implies that the 2.6 million grants awarded over the 2010s led to an additional 259,000 votes (i.e. 1 percentage point of the citizen voting eligible population) being cast in California during the 2020 general election. Effects are similar across demographics and place of origin but are larger among students with the highest GPAs and peak in the most recent in-sample election in 2020. We find some evidence supporting growing treatment effects over time, which we attribute to improved program implementation after a recent policy reform and the rising generosity of California's financial aid programs over the 2010s. Our results are robust to multiple definitions of voter turnout as well as a number of RD implementation choices.

Using location and partisanship records to track Cal Grant recipients and relying on conservative assumptions, we estimate that 1 out of every 66 voters in California cast a ballot in 2020 due to the Cal Grant. The program raises geographic polarization because its largest effects are in politically competitive locations that are home to public research universities. We find that nearly the entire increase in participation occurs among registered independents and Democrats due to the leftward lean of college-educated youth, implying that Cal Grants issued since 2010 raised Democratic margins of victory in California by 0.5 percentage points in 2020.

Evidence from intermediate outcomes and time variation highlight the importance of peer socialization and in-person college attendance among other mechanisms. Cal Grants increase peer socialization by raising the rate at which students enroll in 4-year colleges and live on campus. Effects appear within two years of award receipt, are absent during the year of COVID-19 remote instruction, and raise students' total voter turnout across post-treatment elections between 2010 and 2020. We find suggestive evidence that is relatively more mixed for alternative mechanisms like increased political trust, voter reciprocity, and income effects.

We externally validate our findings with another RD design and 2.5 million students subject to a generosity notch in the Pell Grant, a nationwide financial aid program \citep{denning_propelled:_2019}. Our findings generalize despite the program targeting different types of students and institutions at a different treatment margin. We calculate that the effects of the Pell Grant were likely large enough to change the outcome of the 2020 presidential election. Still, there are two notes of caution in generalizing our results to other settings. First, colleges may differ in the extent to which they facilitate peer socialization, with on-campus housing and non-instructional time spent on campus appearing particularly important. Second, the magnitude of these turnout effects may be amplified or diminished by the extent to which an election cycle mobilizes young voters through youth-oriented issues.

Recent work by \cite{firoozi_effect_2023} identifies the effects of admission to selective colleges on student partisanship, finding that elite institutions do more to influence students' party preferences than their less selective counterparts with an important role for on-campus peer socialization. The common phenomenon across the two papers may explain part of our results to the extent that grant aid shifts students in our sample from community and state colleges toward selective UCs, but we note that changes in enrollment patterns at our policy threshold are smaller than our treatment effects and can explain only a minority of our estimates. Moreover, in this manuscript we tackle the question of how broad-scale \textit{public expenditures} on higher education shape civic engagement, particularly among the less wealthy students and less selective institutions that receive the bulk of college tuition subsidies.


Our evidence illustrates that tuition subsidies can have civic externalities, even when they do not generate meaningful changes in labor market outcomes. The Cal Grant and Pell Grant induce the participation of an informed and active electorate by encouraging low income and college-educated youth with high GPAs to vote. Given our conservative assumptions, we conclude that the civic externalities of education subsidies represent an important marginal external benefit. Our evidence also introduces challenging new questions for the political economy of education finance. Because the benefits of higher education's civic externalities are not symmetric between parties, partisans can have strong private incentives to distort funding levels relative to the social optimum.


\section{California and the Cal Grant in Context} \label{sec:context}
The Cal Grant is the largest tuition-free 4-year college program in the United States, with the California Student Aid Commission (CSAC) offering several awards that have provided over 2.7 million grants to eligible California residents since 2010 \citep{kane_quasi-experimental_2003,bettinger_long-run_2019,dickler_5_2022,scott-clayton_fine_2022}. High school graduates and continuing college students who meet income and academic requirements and enroll at any in-state public university are eligible for four years of tuition-free college under Cal Grant A or three years plus an annual living stipend under Cal Grant B.\footnote{Tuition in this context refers to mandatory systemwide tuition fees. It does not include campus-specific fees for services like student government events, athletics, campus health insurance, etc..} The program also provides a benefit to attend an accredited private institution but cannot be used at community colleges.\footnote{The precise award amount varies based on the type of private institution, but for most of the time period and for the most popular private colleges it was worth roughly 9,000 dollars per year.} These awards are ``first dollar'', meaning that they are provided without consideration of eligibility for most other forms of financial aid, such as student loans or institution-specific grants.

First-time college students who wish to apply for a Cal Grant must file a Free Application for Federal Student Aid (FAFSA) and must separately complete or have their high school submit a Cal Grant-specific form that assists in GPA or test score verification to validate their educational credentials in the December of the calendar year prior to their college enrollment. First-time students typically find out their grant status around the time of their college admission offers in the Spring before they begin college. Renewing students have a streamlined process consisting of filing another FAFSA for the upcoming academic year. For our analysis, we focus on the income eligibility thresholds for the Cal Grant A program between the 2017-2018 and 2019-2020 academic years. These income thresholds vary by family structure as well as applicant characteristics and are adjusted each year based on cost of living increases set by paragraph (1) of subdivision (e) of Section 8 of Article XIII B of the California Constitution. Families whose adjusted assets exceed certain limits, after excluding personal residence and retirement savings, are ineligible. CSAC's switch to ``prior-prior year'' income assessment in 2017-2018 combined with unpredictable cost-of-living adjustments results in several plausibly exogenous discontinuities in eligibility that we use for identification of causal effects.

 Meeting the income and asset requirements is not sufficient for Cal Grant eligibility. Students must also have a 3.0 high school GPA if they are an entering freshman, not hold a bachelor's degree, meet in-state residence requirements, and have a sufficiently low ``expected family contribution''.\footnote{Expected Family Contribution or EFC is a measure of socioeconomic status that depends on family income, family assets, and the cost of attendance at the institution at which a student enrolls.} Hence, we restrict to the subset of students coded as in-state residents without a bachelor's degree who fall below the asset threshold. As we discuss at greater length in Section \ref{sec:method}, we omit students whose family incomes are perfectly divisible by 1,000 dollars throughout our analysis and use the 2017 to 2019 cohorts of FAFSA filers in our preferred specification to address threats to our identification strategy.

The Cal Grant is an ideal policy setting for evaluating the social returns to higher education subsidies, in part, because extant evidence on the pecuniary returns to the program provide important context for assessing civic externalities and social returns \citep{kane_quasi-experimental_2003,bettinger_long-run_2019,hendren_unified_2020}. At the \textit{income threshold}, the Cal Grant has null effects on total college enrollment, but there is a shift in the type of college attended, with an increase in attendance at 4-year private institutions and a reduction in attendance at public 2-year and 4-year colleges and universities.\footnote{We do not use the GPA threshold for identification in our paper. This is because over the 2010-2011 to 2020-2021 timespan in our dataset, the GPA cutoff was fixed at 3.0 and known ex ante, unlike previous work.} The substitution between institutions is not associated with changes in college quality but is associated with higher tuition costs and lower per-student expenditures. The Cal Grant has no impact on labor income for students at the income threshold, contributing to a marginal value of public funds (MVPF) estimated in \cite{hendren_unified_2020} of -0.69, and there is no evidence that it retains individuals in-state prior to 11 years after initially filing for a grant.


The state of California is ideal for estimating the impact of education spending on political participation. Its status as the most populous market of higher education in the United States allows for precise estimates of causal effects. California has the highest post-graduation retention rate of college students among the 50 American states, offering unique advantages for longitudinally tracking students \citep{van_dam_states_2022}. Voter registration records in California are detailed, showing that the political composition of Californian college-educated youth matches that of other American states and making California a setting where roughly three quarters of voters self-identify their political party preferences in administrative data \citep{firoozi_effect_2023}. Notably, party membership is not a prerequisite for participating in any primary elections in California, except presidential primaries, meaning that voters have an incentive to register with the party that best reflects their policy views rather than to strategically register with the party that dominates state politics.

\section{Research Design and Data} \label{sec:method}


    \subsection{Data} \label{sec:ch3_data}

The linked dataset used in this study is a de-identified file comprised of two data sources: the L2 California voter file and financial aid applicant (FAFSA) data provided by the California Student Aid Commission (CSAC) \citep{firoozi_code_2024}. The original dataset contained approximately 16.4 million observations generated by merging the CSAC data and L2 voter data. This matching process was completed by providing CSAC with the L2 voter data, having CSAC match on full name and date of birth, and then receiving back a version of the CSAC data with L2 voter variables and without any personal identifying information.\footnote{We note that there is a match rate of 50 to 60 percent, depending on the particular sample used, between FAFSA files from CSAC and L2's California voter file. Because the merged dataset is de-identified, it is not possible to know the precise registration rate for unique individuals within the dataset. Non-matches between these two datasets are attributable to students moving to another state/country, lack of interest in voting or registering to vote while living in California, or disenfranchisement (for example, due to incarceration or death). Based on the descriptive statistics we observe, we estimate that approximately three quarters of the non-matches between datasets are attributable to Californian residents who simply lack interest in registration or voting, with the remaining quarter attributable to migration out of California.} To comply with the Federal Education Rights and Privacy Act (FERPA) and maintain anonymity, the names and dates of birth of FAFSA filers were never revealed to the authors.

The sample provided by CSAC, which administers the Cal Grant program, includes records on all of the roughly 16.4 million FAFSA filers in California between the academic years 2010-2011 and 2020-2021, excluding 2011-2012.\footnote{Database errors at CSAC prevented retreival of 2011-2012 data. For brevity we refer to the sample excluding this cohort as the full sample in the remainder of the paper.} We estimate that our data cover around three quarters of all Californian first time college applicants, returning college students, and transfers over this timeframe. CSAC's records include detailed information available from the FAFSA form, including family size and structure, adjusted gross income, Cal Grant receipt, housing intent, and ZIP code of origin.

The outcome data used in this study is sourced from L2 Inc., a non-partisan private vendor of political data. Specifically, we use L2 Inc.'s complete California VM2 voter file, which is a retrospective snapshot file that reflects the California voter rolls as of July 2022. This includes identified records on approximately 21 million Californians who are registered to vote, including their political party membership and participation in every election through the 2021 California gubernatorial recall. The file also contains commercial data, which provide additional outcomes and detailed information on the locations where registrants live.

We use several different samples in various parts of this paper. For our main analyses, we focus on the set of students who were likely to be Cal Grant eligible within a 10,000 dollar bandwidth of the income cutoff.\footnote{As mentioned in Section \ref{sec:context}, we restrict to the subset of students coded as in-state residents without a bachelor's degree who fall below the asset threshold. As we discuss at greater length in Section \ref{sec:method}, we omit students whose family incomes are perfectly divisible by 1,000 dollars throughout our analysis to address threats to our identification strategy.} Table \ref{tab:descriptives} shows summary statistics for each sample we use in our paper. We note in Appendix \ref{sec:descriptive} that our main sample is majority non-white, is much more likely to favor the Democratic Party than the Republican Party, has family asset levels (excluding retirement accounts and home equity) that are near zero, and has family income levels roughly around that of the median American family.




    \subsection{Regression Discontinuity Design}\label{sec:RDD}

In this study, we use a fuzzy regression discontinuity design (RDD) to estimate the impact of the Cal Grant program on political participation. The Cal Grant program's main eligibility criterion is family income, which we use as the running variable in our RDD. Specifically, we standardize a student's family income against the income ceilings for Cal Grant A set by the California Student Aid Commission (CSAC). This approach identifies a clear discontinuity in the proportion of students who receive a Cal Grant (as shown in Figure \ref{fig:rd_firststage}).\footnote{Note that there is also a lower income threshold for Cal Grant B eligibility, but it generates little variation in total Cal Grant receipt and leads a subset of students eligible for Cal Grant A to receive Cal Grant B instead. Cal Grant C has the same income eligibility ceiling as Cal Grant A, but there is near zero discontinuity at the threshold as Cal Grant C accounts for only 2 percent of Cal Grants awarded.}

The fundamental assumption of our fuzzy RDD is that the eligibility threshold serves as a clear cutoff for program participation over which students do not have perfect control. In the context of the Cal Grant program, this means that a student whose family income falls just below the eligibility threshold is similar to one whose income falls just above it except for the fact that the latter is ineligible for the program. This assumption implies that the distribution of observed and unobserved characteristics is continuous around the threshold, which is a critical prerequisite for the validity of the design.

We assess the assumptions underlying our fuzzy RDD using several tests. First, we implement a McCrary test to confirm the fuzzy RDD's validity \citep{mccrary_manipulation_2008,cattaneo_manipulation_2018}. This test assesses whether there is a discontinuity in the density of observations at the threshold, which could indicate that individuals are manipulating their reported income to become eligible for the Cal Grant program or selecting into the analysis sample by filing a financial aid application based on program eligibility. As demonstrated in Figure \ref{fig:mccrary_2017} -- and as we confirm through formal tests -- there is no evidence of a density discontinuity at a 90 percent confidence interval ($p$-value=0.50). Second, we conduct balance tests to ensure that observable characteristics trend continuously across the policy threshold. In Figures \ref{fig:rd_covariates1} through \ref{fig:rd_covariates3}, we present evidence that pre-FAFSA covariates including voting patterns are balanced, finding only three rejections out of 18 variables at a 90 percent confidence interval with a 10,000 dollar bandwidth.\footnote{We also show an omnibus test for balance by predicting voter turnout rates using our full set of covariates in Figure \ref{fig:rd_omnibus}. The point estimates and confidence intervals are available for a range of potential bandwidths in Figures \ref{fig:bw_covariates1} through \ref{fig:bw_covariates3}. Our results at a 10,000 dollar bandwidth are roughly consistent with a random rejection rate. However, it is worth noting that these measures likely overstate imbalance because receiving a Cal Grant directly compels the provision of variables like GPA and may make the reporting of other covariate information to CSAC more likely by increasing the rate at which students update their information and refile FAFSAs in future years.} Our balance tests' results are robust to varying the bandwidth around the income ceilings and we find only one rejection of the null hypothesis at a narrower (2,000 dollar) bandwidth, as demonstrated in Figures \ref{fig:bw_covariates1} through \ref{fig:bw_covariates3}. Third, we conduct placebo falsification tests to evaluate the design's validity.\footnote{These tests entail assigning ``placebo'' thresholds and assessing the estimated treatment effect at the true threshold relative to the placebo thresholds. If the placebo test fails, it implies that the design may not be valid but, as we discuss in Section \ref{sec:result}, the outcomes of these tests confirm the validity of our findings.}

We also address potential threats to identification that are specific to the Cal Grant policy by focusing our main results on the 2017-2018 to 2019-2020 cohorts of FAFSA filers before broadening our sample to the 2010-2011 to 2019-2020 filers. The first threat is that some students who barely miss out on the grant may leave the state, which would make them unobservable in the outcome data from L2's California voter file. To address this risk, we use four different approaches that we describe in detail in Appendix \ref{sec:appendixasub}. We start by pointing  to the absence of out-of-state discontinuities in tax filings and college enrollment records. We then use a unique feature of the L2 dataset, in which out-of-state movers have their historical data pruned, and show the absence of a discontinuity in pre-treatment voter turnout. Next, we use a subset of our sample overlapping with the sample in \cite{firoozi_effect_2023} to directly illustrate that the Cal Grant is estimated to have null effects on voter registration outside of California. Finally, we choose to be conservative by beginning our results section with a focus on the 2017-2018 to 2019-2020 sample, who filed a FAFSA less than 5 years prior to our voter file snapshot, minimizing the possibility of selection bias due to out-of-state migration.

The second threat to our identification strategy is that some students may get married and change their legal last name after filing a FAFSA, making it difficult to match their voter registration records to CSAC data. We address this issue by focusing on the 2017-2018 to 2019-2020 sample because few people get married and change their names on the voter roll within 5 years of filing a FAFSA. We also find that women, a group more likely to change legal names, do not drive our estimated effects and note that there is little difference in match rates between female and not female-identifying students who filed a FAFSA for the 2015-2016 academic year or later.

The final threat to our fuzzy regression discontinuity design comes from time-specific concerns. Prior to 2017-2018, families could have seen the income thresholds before the end of the tax year and attempted to manipulate their reported income or selected into FAFSA filing based on unobservable characteristics. Relying on the post-2017 period is helpful, because ``prior-prior year'' income evaluation took effect, making it much harder to anticipate eligibility thresholds ex ante to filing a tax return.\footnote{Our correspondence with CSAC alongside \cite{bettinger_long-run_2019} point out that California disclosed financial aid thresholds ex ante to completion of the tax year in the 2000s and early 2010s, presenting threats to identification for pre-2017 cohorts.} We also exclude students whose family incomes are bunched at perfect multiples of 1,000 dollars, as these families may have greater discretion to manipulate their reported income.\footnote{These students represent roughly 6 percent of the sample within 10,000 dollars of the income ceiling and come from families that are more likely to have salaried workers than hourly workers, are likely to have little to no capital gains income from stocks and bonds, and have fewer total sources of stochastic income relative to fixed income. As a consequence, they may be more capable of manipulating their income or are more likely to be aware of their eligibility because they can easily recall their income and opt-in to filing a FAFSA after finding out the thresholds.} COVID-19 is another potential threat, because Californian colleges were overwhelmingly remote in the 2020-2021 academic year, changing the mechanisms that were discontinuous across the threshold.\footnote{As one example, peer socialization across the threshold would have been greatly curtailed due to extensive remote instruction and the sharp reduction of interaction in on-campus housing.} To address this issue, we use the 2020-2021 cohort solely to analyze mechanisms.

Taking continuity of the conditional expectations function as given, we estimate the following reduced-form RD equation:
\begin{equation} Outcome_{i}=\phi_0 +\phi_1 Above_{i}+ f(Income_{i})+\textbf{X}_{i}'\phi_2+\nu_{i},
\end{equation}
where $Outcome_{i}$ is an outcome for student $i$, $Income_{i}$ is a student's centered income with the Cal Grant A income ceiling normalized to zero, $Above_{i}=\mathbb{I}[Income_{i}> 0]$ is a binary variable for a student being above the income ceiling specific to their cohort and family structure, $f(\cdot)$ is a continuous function, $\textbf{X}_{i}$ is a vector of covariates, and $\nu_{i}$ is an idiosyncratic error term. Assuming the RD assumptions hold, our $-\hat{\phi_1}$ estimate identifies the average effect of being income-eligible for the Cal Grant among students local to the threshold.

We also use a fuzzy RDD approach to estimate the impact of Cal Grant receipt on political participation. Specifically, we treat equation 2 as a first-stage equation:\begin{equation} CalGrant_{i}=\gamma_{0} +\gamma_{1} Above_{i}+ g(Income_{i})+\textbf{X}_{i}'\gamma_{2}+u_{i},
\end{equation}
where $CalGrant_{i}$ is an indicator for individual $i$ having received any Cal Grant award, $Above_{i}=\mathbb{I}[Income_{i}> 0]$ is a dummy variable indicating whether individual $i$'s family income was above the ceiling for Cal Grant eligiblity, $Income_{i}$ is an individual's normalized family income,  $g(\cdot)$ captures the relationship between normalized income and Cal Grant receipt, $\textbf{X}_{i}$ is a vector of pre-FAFSA characteristics, and $u_{i}$ is an idiosyncratic error term. We estimate $\hat\gamma_1$ as the first-stage impact of being above the income ceiling on Cal Grant receipt.

Next, we use equation 3 as an outcome equation to characterize the relationship between political participation and Cal Grant receipt:
\begin{equation} Outcome_{i}=\beta_{0} +\beta_{1} CalGrant_{i}+ h(Income_{i})+\textbf{X}_{i}'\beta_{2}+\varepsilon_{i},
\end{equation}
where $Outcome_{i}$ is an outcome for student $i$,  $h(\cdot)$ reflects the relationship between normalized family income and the outcome of interest, and $\varepsilon_i$ is an idiosyncratic error term. We combine equations (2) and (3) and estimate $\beta_1$ using two-stage least squares with $Above_{i}$ as our excluded instrument. Our estimate $\hat\beta_1$ identifies the average effect of receiving a Cal Grant on each outcome among students local to the threshold. We then test our estimates for robustness to a number of different RDD implementation choices. We vary the order of a polynomial control for the running variable, include an expansive set of pre-FAFSA controls, flexibly change the bandwidth used for inference, and estimate bias-aware confidence intervals to demonstrate robustness \citep{calonico_robust_2014}.\footnote{The controls we use include foster youth status, female self-identification, voter turnout in the general election prior to FAFSA filing, family income, family financial assets excluding personal residence and retirement saving, marital status, first year freshman status, father's college attainment, mother's college attainment, cohort fixed effects, family size, and ZIP code of origin measures of the minority, Black, Hispanic, and Asian population, voter turnout, voter conservatism, as well as mean income and asset levels.}

\section{Results} \label{sec:result}



    \subsection{Main Results}\label{sec:main}

We begin by plotting our outcomes of interest against a student's centered family income in Figure \ref{fig:rd_outcomes}. The first outcome is voter registration in 2022, which we use as a low intensity measure of political participation. The second outcome is voter turnout in the 2020 general election, which is the first and only general election in which all students in our main sample could have participated assuming students enter college at 18 years of age. The third outcome, labeled ``General Election Turnout'', is the voter turnout rate in all general elections between 2010 and 2020 that took place after the academic year in which students filed their FAFSA and received a grant. Finally, we interact the voter turnout rate with an indicator for being registered as a Democrat or independent\footnote{In California, independent refers to voters who are either registered to vote with no stated party preference or with a minor political party like the Green Party, the Peace and Freedom Party, the Libertarian Party, or the American Independent Party.}, which captures the extent to which this turnout increase is attributable to center-left voters. Recent evidence from this setting shows that the interaction term is a strong predictor of support for Democratic candidates because, similar to other American states, roughly 75 percent of Californian, college-educated youth who are registered independents self-report favoring the Democratic Party over the Republican Party \citep{firoozi_effect_2023}.\footnote{\cite{firoozi_effect_2023} finds that registered independents in this setting are as likely to support Democrats as registered Republicans are to support Republicans, are more than twice as ideologically close to registered Democrats as registered Republicans on economic policy and sociocultural issues, and favor Democratic and liberal candidates over Republicans and conservatives by more than a 5 to 1 margin in their political donations. Data from election returns in college campus precincts in both California and other American states illustrate that Republican candidates receive vote totals approximately equal to their share of the registered voters who cast a ballot, corroborating the idea that young, college-educated independents disproportionately vote for Democrats \citep{firoozi_effect_2023}.} On balance, the results show that students who are below the income ceiling and are, therefore, income-eligible for the Cal Grant are more likely to participate in the political process, with essentially the entire increase among registered Democrats and independents.

In Panel A of Table \ref{tab:mainresults}, we formally estimate the effect of receiving a 2017-2019 Cal Grant on students' political outcomes. Each column reflects results for a different outcome of interest and each panel reflects estimated effects for a different policy and sample of students. The same RD specification is used in all estimates: a local linear estimate at a 10,000 dollar bandwidth of the income threshold with a uniform kernel and no pre-treatment covariate controls. The outcome in Column 1 is a binary indicator for being registered to vote in the state of California in 2022. Column 2 shows results for a binary indicator for casting a ballot in the 2020 presidential general election. Column 3 reflects the total voter turnout rate for a student across all post-treatment elections in which they could have participated.\footnote{For these more recent cohorts of students, only the 2020 and/or 2018 election are included as post-treatment elections.} Column 4 shows results for a binary indicator for having ever voted in a federal or state general election. Column 5 tests the results for an interaction term between an indicator for being registered as a Democrat or independent and the voter turnout rate in Column 3. Column 6 tests results for an interaction term between being registered as a Republican and the voter turnout rate in Column 3.

Receiving a Cal Grant sharply increases political participation among these students. Although results for voter registration are noisy and sensitive to RD specification choices in our robustness checks, we find that voter turnout rates rise between 8 and 10 percentage points across the various measures of this outcome in Columns 2 through 4. This indicates that the Cal Grant raises both the intensive and extensive margins of political participation. The entirety of the increase is attributable to voters who are likely to support Democratic Party candidates, with the point estimate on center-left turnout in Column 5 exceeding the point estimate for total turnout in Column 3. There are no detectable effects of receiving a Cal Grant on voter turnout rates among voters registered with the Republican Party, with a negative point estimate in Column 6.

To ensure that the estimates in Panel A are robust, we re-estimate results varying a number of RD implementation choices. We first estimate results across different bandwidths, varying the inclusion of pre-treatment covariate controls and the order of a polynomial control for the runnning variable (See Appendix \ref{sec:cal1719} and Table \ref{tab:results1}). We then estimate results separately by gender identity in Table \ref {tab:gender_results1} to show that our results are not driven by women changing their last names after marriage. We then present findings in Figures \ref{fig:bw_results} through \ref{fig:bw_centerleftvote}, highlighting that the estimated reduced-form discontinuities in our outcome variables are relatively stable across a wide range of bandwidths, the inclusion of pre-FAFSA covariates, and the use of a quadratic control for the running variable.\footnote{The point estimates are negative in this case because we are estimating the intent to treat (ITT) results of being above the ceiling for Cal Grant eligibility (too rich to receive a Cal Grant).} Next, we show in Table \ref{tab:rdrobust1} that our main results are robust to the use of CCT bias-aware confidence intervals. Finally, we conduct falsification tests in Figure \ref{fig:bw_falsif1_results1}, in which we assign a dummy variable for Cal Grant income-eligibility based on placebo income thresholds and then compare the results to those derived from the true income ceiling, finding that our results at the true policy threshold are larger than the 95th percentile of results in these placebo tests for all measures of voter turnout.\footnote{We generate a ``placebo threshold'' at each 500 dollar increment along centered family income, and compare the estimated reduced form impact of these synthetic policies relative to the true policy. Placebo thresholds are bounded between -20,000 and +60,000 dollars relative to the true income ceiling, because this avoids false positives from capturing discontinuities taking place at family incomes of zero at the lower bound and this spans up to the 98th percentile of centered income on the upper bound. A 10,000 dollar bandwidth is used to remain consistent with our preferred specification. We exclude discontinuities within a 10,000 dollar bandwidth of the true cutoff to avoid generating false positives by including the actual policy discontinuity in our placebo estimates.}

\subsection{Time Variation} \label{sec:time}

We extend our analysis to include FAFSA filers from the 2010-2011 to 2020-2021 academic years to examine time variation in treatment effects. While we acknowledge the identification risks that come from earlier and later cohorts (i.e. from potential income manipulation in the earlier time period and from remote instruction in the latter time period), we believe that these results are useful for understanding the causal mechanisms underlying the Cal Grant's impact on political participation.

In Panel B of Table \ref{tab:mainresults}, we show the impact of the Cal Grant on political participation for the 2010-2011 to 2019-2020 sample of likely Cal Grant-eligible FAFSA filers using the same set of outcomes and RD specification as Section \ref{sec:main}. The results are consistent with those of the more recent sample of FAFSA filers, but with smaller magnitudes. Although the impact of Cal Grants on voter registration in Column 1 is again somewhat noisy and sensitive to model specification (See Table \ref{tab:results2}), we find significant increases in voter turnout between 4 and 7 percentage points across the definitions in Columns 2 through 4. In Columns 5 and 6, we again find that over 85 percent of the effect on voter turnout accrues to registered Democrats and independents, consistent with our evidence from more recent cohorts. Each of these results are robust to changes in RD specification choices like the inclusion of pre-treatment controls, varying the order of a polynomial control for the running variable, and changing the bandwidth around the threshold used for inference (See Table \ref{tab:results2}).

In Appendix \ref{sec:timevariation}, we use other dimensions of time variation to understand the Cal Grant's impact on political participation. In Table \ref{tab:results3} we find that the point estimates of these effects are largest within one or two years of receiving a Cal Grant and appear to shrink over longer time windows, but we cannot reject the null hypothesis that they are flat over time. The 2020 COVID-19 student cohort\footnote{These are students whose instruction was almost entirely online prior to election day and had essentially no opportunities for peer socialization.} have null effects in Table \ref{tab:covid1}, in stark contrast to the large effects observed for the 2017-2019 cohorts of students.

Finally, we rerun the specification from Table \ref{tab:mainresults} Panel A Column 3, which includes the most recent four  cohorts, for each respective election in Table \ref{tab:elexovertime}. Our point estimates in Table \ref{tab:elexovertime} are larger in presidential cycles than midterm cycles, suggesting that campaigns in which youth-focused issues are salient may amplify the effects we observe, and peak in the 2020 presidential election for the Cal Grant. We note several reasons why the Cal Grant's effect sizes may be largest in 2020 while other grants' effects are static across presidential elections: changes in the composition of near threshold Cal Grant students over time as the extensive margin of generosity (income thresholds) rises, increases in the intensive margin of the generosity as the Cal Grant's award size grows, random noise given that the confidence intervals of some of these estimates overlap, and unobserved biases in the estimated treatment effects for pre-2017 entering student cohorts for whom the Cal Grant's income thresholds for eligibility were visible ex ante to filing a tax return. We evaluate each of these possibilities in greater depth in Appendix \ref{sec:timevariation}, finding  support for the idea that the pattern is driven by the delay of Cal Grant eligibility announcements paired with growth in Cal Grant generosity.

\subsection{External Validity, Heterogeneity, and Intermediate Outcomes}\label{sec:HTE_int}

We are also able to assess the external validity of our findings for the Cal Grant by estimating the impact of the Pell Grant, the largest nationwide tuition subsidy in the United States. We use an income threshold that generates variation in the generosity of Pell Grants to identify the impact of the intensive margin of financial aid generosity, as opposed to the extensive margin of financial aid receipt we tested with the Cal Grant. We implement an RD design comparing students just above and below an income ceiling that hovers in the 20,000 to 30,000 dollar per year family income range, and provide a detailed analysis and description of our approach in Appendix \ref{sec:pell}. Figure \ref{fig:pell_outcomes} displays each of the same outcomes of interest as Figure \ref{fig:rd_outcomes} from Section \ref{sec:main}, but plots students based on their family income normalized to the threshold for receiving a more generous Pell Grant. Matching the Cal Grant's pattern, the students just below the income limit to receive a more generous Pell Grant are more likely to register to vote, have higher post-treatment voter turnout rates, and are more likely to cast a ballot as either a Democrat or an independent than their counterparts just above the income limit.

In Panel C of Table \ref{tab:mainresults} we estimate the intent to treat effect of this increase in Pell Grant generosity, using the same outcomes and specifications as we did for Cal Grants and noting that the dollar value of this increase in Pell Grant aid is roughly one tenth that of receiving a Cal Grant. Pell Grants have noisy, if any, positive effects on voter registration rates in Column1. Similar to the Cal Grant, the Pell Grant induces clear increases in voter turnout across each definition of the variable in Columns 2 through 4, with over 80 percent of the increase attributable to Democratic and independent voters in Column 5 rather than Republican voters in Column 6. We note that each of these outcomes, with the exception of voter registration, are robust to the use of different RD implementation choices (See Table \ref{tab:pell_results}).

In Appendix \ref{sec:HTE} we evaluate the heterogeneous effects of Cal Grants on political participation across various demographic and socioeconomic dimensions. The analysis looks at treatment effect heterogeneity by race, ethnicity, socioeconomic status, political background, and high school GPA. While we find little variation in the impact of Cal Grants based on racial or ethnic background, socioeconomic factors, or political environment of a student's home location., there are stronger effects on political participation among students with higher high school GPAs, above 3.41 or roughly a B+ on the American grading scale.

After examining heterogeneity we turn to housing choices in Figure \ref{fig:rd_housing1}, using noisy CSAC records on students' reported housing intent for Cal Grant non-recipients and actual housing outcomes for recipients. The RD plots in Figure \ref{fig:rd_housing1} show a clear and significant increase in the share of Cal Grant recipients living on campus, accompanied by sharp declines in the share of students living off campus with family or legal guardians or who did not report their housing intent. The decrease in the latter group is due to the fact that receipt of a Cal Grant requires students to disclose their housing choice \textit{ex post}. We note that, because of the housing choice variable's definition, our estimates will be subject to some amount of bias in the direction of students' underestimated housing expectations.

 The large changes we observe in reported housing status and intent nonetheless point to large impacts from financial aid. We show the impact of the Cal Grant on campus housing choice formally in Table \ref{tab:housing1}. Our results indicate a 17.17 percentage point increase in the share of students living on campus as a result of the Cal Grant under our preferred specification in Column 1, with corresponding decreases of 13.17 percentage points in the share of students with no stated housing intent and 9.20 percentage points in the share of students living off campus with family or guardians. This suggests that the Cal Grant induces some students who did not have a strong preference or who planned to live with family to instead live on campus. We note that this effect could be driven in small part by the fact that Cal Grants reduce enrollment at community colleges, which do not offer on-campus housing, and increase enrollment at 4-year colleges, which house a high proportion of students on campus.

\section{Discussion} \label{sec:discussion}
\subsection{Mechanisms}

We present evidence on four plausible causal mechanisms that could operate independently or collectively: campus-based socialization, electoral reciprocity, civic trust, and income effects. While this is not an exhaustive list of all possible explanations for the civic returns to tuition subsidies, we believe these mechanisms are the most compelling ones in our setting. Moreover, we note that several other sets of mechanisms can effectively be ruled out as meaningful explanations of these externalities, given that they rely on first-stage effects that do not occur or that are small in our policy setting.

Because the 10 percentage point effect of the Cal Grant on voter turnout exceeds the 0.2 percentage point impact on college enrollment and the 6.8 percentage point substitution from 2-year to 4-year colleges, changes in college enrollment are insufficient to fully account for the causal effects of the Cal Grant \citep{bettinger_long-run_2019}. Even assuming that switching from a 2-year college to a 4-year college increases voter turnout by 20 percentage points due to mechanisms like voter registration drives or other differences between campuses, changes in enrollment could account for no more than 15 percent of our total estimated externalities.\footnote{A 20 percentage point estimate is larger than what is suggested by the extant literature and is intended to represent a strict upper bound \citep{bell_which_2024}. } As Appendix Table \ref{tab:conditional_turnout} shows, voter turnout rises even conditional on voter registration, suggesting that plausible mechanisms must raise voter turnout beyond their impact on voter registration status or enrollment choice.

Given that changes in enrollment choices are not sufficiently large for differences between campuses to explain the majority of our observed effects, we turn to our four main plausible mechanisms: on-campus socialization, electoral reciprocity, civic trust, and income effects. First, Cal Grants may facilitate on-campus socialization in 4-year colleges by extending the amount of time students spend at college campuses and increasing the share of students living on-campus independently of their enrollment choice. Second, students may think about the social transfers they received and reciprocate support from Democratic politicians by casting votes for the same politicians. Third, receiving a large conditional cash transfer early in adulthood may increase young people's trust in civic institutions and government. Fourth, financial aid programs may be similar to an increase in total family income, raising voter turnout and shifting students toward the political right.

\subsubsection{Peer Socialization}

We start with the idea that exposure to the college campus environment increases socialization and, in turn, raises voter turnout \citep{persson_education_2015,firoozi_effect_2023}. The results of Table \ref{tab:housing1} demonstrate that Cal Grant receipt increases the rate at which students live on campus by 15 to 20 percentage points and reduces the rate at which they are housed with their family or guardians. Cal Grants cause substitution from 2-year to 4-year colleges in California and generate imprecise but positive increases in the share of students completing 4-year degrees and graduate degrees, extending the amount of time students spend in higher education \citep{bettinger_long-run_2019}. Evidence from time variation also supports the idea that some element of in-person exposure to a 4-year campus environment matters. As Table \ref{tab:results3} illustrates, effects appear within the first two years of receiving a Cal Grant, the same time period over which a student is more likely to live on campus and to remain enrolled in higher education. Null effects for the 2020-2021 COVID-19 year in Table \ref{tab:covid1} may likewise suggest that online instruction does not generate similar effects. In Appendix \ref{sec:externalvalid} we point to evidence that this form of socialization is likely passive, rather than due to active events like on-campus protests, and show that colleges in conservative states are politically similar to those attended by in-sample students.

\subsubsection{Voter Reciprocity}

The second mechanism we test is the idea that students reward government transfer payments like the Cal Grant with their votes, a pattern of reciprocity observed among conditional cash transfers and stimulus spending in other contexts \citep{manacorda_government_2011,pop-eleches_targeted_2012,huet-vaughn_stimulating_2019,firoozi_economic_2024}. We find evidence that is generally mixed on this hypothesis. Contrary to the hypothesis, we show that the largest increase in turnout takes place in the 2020 presidential general election, when California did not vote on state government offices that are critical for tuition policy (see Tables \ref{tab:results1}, \ref{tab:results2}, and \ref{tab:elexovertime}).\footnote{For context, elected officials including the Governor and Superintendent of Public Instruction hold seats on university governing boards, where they can and do directly vote on tuition policy and hold line-item vetoes over Cal Grant funding.} The effects we observe for the intensive margin of financial aid generosity in Section \ref{sec:pell} suggests that the salience of receiving any financial aid is not the primary driver of these effects. Furthermore, null effects on turnout for the 2020-2021 students who received financial aid but did not attend in-person college as a result of COVID-19 suggest that reciprocity alone may not be sufficient to generate our observed results.

While these findings are inconsistent with reciprocity, they do not conclusively rule it out. For example, Cal Grant recipients may have voted for Democratic politicians at the federal level in 2020 because federal Democrats are co-partisans who support national financial aid programs. In Appendix Table \ref{tab:partisanship}, we test whether Cal Grants change students' party registration or vote choice.\footnote{We show the same exercise for the federal Pell Grant in Table \ref{tab:pell_partisanship}.} To the extent that Cal Grants have any impact on registration, the whole of the increase occurs among Democrats and independents with insignificant, negative effects for the Republican Party, shifting partisanship to the left. To the extent that such changes reflect partisan persuasion and differ from an even split between parties, this evidence may be consistent with a voter reciprocity mechanism.

\subsubsection{Civic Trust}

The third mechanism, called civic trust, posits that receiving in-kind transfer payments like financial aid early in adulthood fosters greater trust in civic institutions and government, which then leads to higher voter turnout \citep{alesina_economics_2000,wang_political_2016}. This mechanism is supported by the rapid rise in post-treatment voter turnout and the disproportionate share of newly registered voters affiliating with Democrats, who are generally perceived as favoring more government spending. However, the seemingly short-lived treatment effects in Section \ref{sec:timevariation} are not especially consistent with a durable change in civic trust and support for government. Furthermore, despite predictions by \cite{alesina_economics_2000} that differences in pre-treatment civic trust based on race and income should lead to heterogeneous treatment effects of financial aid by race and socioeconomic status, our data does not support significant racial or socioeconomic heterogeneity in these effects (see Table \ref{tab:hte1} and Section \ref{sec:pell}).

\subsubsection{Income Effects}
The final mechanism we evaluate is that financial aid induces higher voter turnout through income effects if voting is a normal good and reducing the cost of college attendance reduces student labor hours.\footnote{Although the Cal Grant does not generate an increase in long-run earnings at the income threshold we study, it may still function similarly to a short-run increase in family income \citep{bettinger_long-run_2019}.} Under this mechanism we should expect short-run increases in voter turnout that dissipate over time, effect sizes that are roughly in line with the association between income and voting, a partisan shift toward the political right, and larger effects for students who start off with lower family incomes \citep{romer_individual_1975,meltzer_rational_1981}. To the first point, we are able to show in Table \ref{tab:results3} that there is some evidence that effects are largest shortly after receiving the grant.

However, evidence on the latter three points are inconsistent with income effects. We note that the estimated impact of both the Cal Grant and Pell Grant on voter turnout are conservatively in the range of 0.4 to 0.8 percentage points per thousand dollars, which significantly exceed the roughly 0.2 percentage point association between turnout and income. Moreover Tables \ref{tab:partisanship} and \ref{tab:pell_partisanship} show that financial aid programs do not shift registration toward the political right. We also do not find evidence that effect sizes are considerably larger for lower income students either in Table \ref{tab:hte1} or through the comparison of Pell Grant and Cal Grant effect sizes in Section \ref{sec:pell}, which are identified at different points of the income distribution.\footnote{The typical Pell Grant near-threshold student has a family income of roughly 25 thousand dollars per year, whereas the typical Cal Grant near-threshold student has a family income of roughly 73 thousand dollars per year.} The similar per-dollar estimates we find for both the intensive and extensive margin of financial aid across a wide range of family income runs contrary to the diminishing marginal returns we would otherwise expect from an income effects mechanism.\footnote{In Appendix \ref{sec:income} we show that the treatment effects we estimate contradict the diminishing returns to income we find in registration, partisanship, and turnout.}

\subsubsection{Significance}
On balance, our results suggest that the Cal Grant's effects are likely driven by peer socialization, with plausible contributions from voter reciprocity and increased civic trust. The implications of these findings on mechanisms are twofold. First, because our work suggests that in-person socialization within colleges can play an important role in the formation of political behavior, our work contributes to the extant literature on political cleavages by education. By highlighting the role of dormitories and on-campus non-instructional time, our work could help explain the puzzle of why higher education is associated with political engagement in rich and middle-income democracies, but not other contexts \citep{krishna_poor_2008,gallego_understanding_2010,gethin_political_2021}.\footnote{Low-income countries typically have universities and colleges that house a much smaller proportion of students on-campus and in dormitories. Non-democracies do not typically allow for a robust political and activist culture to be established on college campuses.} Second, by raising the possibility that Cal Grants simultaneously change partisanship and voter turnout, this paper poses challenging questions about the incentives elected officials have when setting tuition subsidies.


\subsection{Macro-Level Impact and Policy Implications}\label{sec:policy}
Having tested the generalizability of our findings with the federal Pell Grant and seeing suggestive evidence of large impacts on the political system, a natural question arises about the extent to which the Cal Grant itself has changed California's political landscape. We approach this question in three ways, using the full sample of 16.4 million FAFSAs and 21 million registered California voters. First, we evaluate the aggregate increase in voter turnout induced by the Cal Grant relative to total participation in California to provide a sense of scale. Second, we consider the electoral implications of these results based on the partisan leanings of new voters who are induced to enter the electorate as a result of the policy. Third, we analyze the geographic distribution of Cal Grant recipients after voter registration to identify the types of community that receive the bulk of these effects.

We have made a deliberate effort to err on the side of conservative assumptions throughout this analysis. Our goal is to present an accurate picture of the impact of tuition subsidies on California's political outcomes, but we note that our assumptions will likely understate the full extent of this impact.

To estimate the impact of the Cal Grant on political participation, we assume that the local average treatment effect (LATE) for compliers near the income ceiling is the same as the average treatment effect for all Cal Grant recipients (ATT), including those with family incomes well below the income threshold we study. We view this assumption as reasonable because we find little evidence of heterogeneity by race or socioeconomic status in Table \ref{tab:hte1} and because estimated treatment effects of the Pell Grant in Section \ref{sec:pell} are similar on a per-dollar basis for a pool of students who have roughly 1/3rd as much income.\footnote{The typical Pell Grant near-threshold student has a family income of roughly 25 thousand dollars per year, whereas the typical Cal Grant near-threshold student has a family income of roughly 73 thousand dollars per year.} However, if this assumption is inaccurate, we anticipate that the impact of the Cal Grant on political participation may be plausibly greater for inframarginal Cal Grant recipients, who are lower income on average. Moreover, we have used the presidential election as a benchmark outcome to increase the denominator of our estimates, as it had the highest turnout rate of any election in 2020.

It is worth noting that our analysis is narrow in defining the people we consider treated by the Cal Grant. We assume that Cal Grants issued prior to 2010-2011 and in 2020-2021 have no effect on political outcomes, which excludes more than a decade of treated students. Furthermore, we assume that no other financial aid programs, including the federal Pell Grant, have impacted students' political participation despite our strong evidence to the contrary from Section \ref{sec:pell}. We also assume that there is no persuasion effect despite suggestive evidence in Table \ref{tab:partisanship}, meaning that students who would have voted regardless of Cal Grant receipt are not more likely to cast their vote for Democrats to reward them for the policy. Finally, we assume that there are no spillover effects on the voter turnout rates or party preferences of Cal Grant recipients' parents, families, or any other individuals.

In Appendix \ref{sec:macro} and Table \ref{tab:envelope}, we show our calculations of the estimated impacts of the Cal Grant on the 2020 presidential election results in California. Across three columns, we generate a lower bound estimate of the impact on partisanship and turnout, an upper bound estimate, and our best estimate using preferred specifications and the most accurate data available. We vary three assumptions across the columns to capture the full range of plausible results: the estimated treatment effect on 2020 voter turnout, the partisan composition of new voters who enter the electorate, and the aggregate number of people who received a Cal Grant over the 2010-2011 to 2019-2020 cohorts.

Under our best estimate from our preferred specifications, we find that Cal Grants issued between 2010-2011 and 2019-2020 increased aggregate 2020 general voter turnout in the state of California by approximately 259,000 votes and Joe Biden's margin of victory over Donald Trump by 168,000 votes. These figures correspond to a roughly 1 and 0.5 percentage point increase in the California CVEP turnout rate and the Democratic statewide margin of victory. Stated another way, 1 out of every 66 Californian voters in 2020 voted because of the Cal Grant program along with 1 out of every 55 Californian Biden voters. The magnitude of these effects has important policy implications, especially given that they omit all other forms of financial aid. The increases occur among low-income youth -- one of the most underrepresented groups in the electorate -- and college-goers with high GPAs, who would plausibly contribute to an informed citizenry.

Figure \ref{fig:map} provides insight into the geographic distribution of the Cal Grant's impact on voter turnout. Specifically, we show the number of Cal Grants received per capita in each county based on the 2022 voter registration address of registered Cal Grant recipients. Our findings are noteworthy for two reasons. First, although the sample is disproportionately Hispanic and Asian youth that originate from low income cities and rural areas, we observe that the effects are most prominent in affluent, suburban counties that are dominated by public research universities. The top quartile of California's 58 counties by per capita Cal Grant impact are home to six out of the UC's nine undergraduate campuses (See Figure \ref{fig:campus_map} for campus locations). Second, the new voters' concentration in well-educated suburbs is likely to accelerate the suburbs' leftward shift, which was already happening due to the growing education cleavage in Western democracies \citep{gethin_brahmin_2021}. For example, the largest per capita impacts in Southern California are not concentrated in urban Los Angeles, but Orange County and San Diego, which are each home to one UC and two CSU campuses and have recently supported Republicans running for statewide office.


Given the stark leftward lean of college-educated youth across democratic countries, there are clear partisan incentives in the expansion of higher education and tuition subsidies. For left of center politicians, the results are unambiguous. Tuition subsidies can raise political participation among a group that is disproportionately likely to support their candidates, regardless of its impact on earnings and degree attainment. This poses a risk that left of center candidates could support tuition subsidies to increase their vote share, even when such policies have no benefits for labor market outcomes. Conversely, right of center politicians could face a tradeoff between potentially higher earnings, degree attainment, and political participation on the one hand and a net loss in votes on the other. At worst, this tension may lead to misaligned incentives for policymakers and partisan cycles in subsidies for higher education observed in other contexts \citep{ortega_state_2020}. Binding budget formulas that peg spending to the youth population and price levels or supermajority rules to change tuition subsidies may be welfare enhancing under such circumstances.


\section{Conclusion} \label{sec:conclusion}

We use a regression discontinuity design and 16.4 million financial aid applications to estimate the impact of a tuition-free college program, the Cal Grant, on political participation among college students. The results show that Cal Grant receipt sharply increased voter turnout, almost entirely among left-leaning voters. Under conservative assumptions, our preferred specification suggests that the 2.6 million grants awarded over the 2010s induced an additional 259,000 people to vote in the 2020 general election in California, raising the voter turnout rate and the Democratic margin of victory by 1 and 0.5 percentage points respectively. These effects are amplified by the disproportionate share of the electorate such voters represent in educated suburbs and small college towns.

This paper contributes to a notable shift in the literature on education spending's civic externalities, which increasingly posits socialization rather than non-cognitive skills as a key mechanism \citep{persson_education_2015,mendelberg_college_2017}. Our evidence suggests an important role for in-person time and on-campus residence in higher education, particularly in recent elections like the 2020 US presidential election. Time variation in treatment effects and evidence on mechanisms pose important questions about how our results will apply to  elections in future time periods and other settings. The magnitude and generalizability of these externalities will likely hinge on (1) the extent to which educational institutions facilitate socialization through in-person, non-instructional time, and (2) the salience of youth-oriented, socio-cultural issues and platforms that may be responsible for the growing political-education gradient \citep{krishna_poor_2008,gethin_political_2021}.

A key policy implication of this study is that tuition subsidies benefit the electoral prospects of left-of-center political parties, despite recent evidence suggesting that higher education shifts students toward the economic right \citep{mendelberg_college_2017,scott_does_2022}. This is in large part due to the pre-existing cultural liberalism of college educated youth, but there are plausible contributions from the persuasive effects of socialization as well \citep{apfeld_higher_2022,firoozi_effect_2023}. The partisan externalities of higher education could act as a barrier to the expansion of these programs, as partisan policymakers have a private incentive to distort the level of financial aid provision relative to the social optimum. Such motivations present tough questions for the political economy of education finance because they may encourage elected officials to expand or reduce education subsidies with little regard to beliefs about their impact on earnings or social welfare.

The Cal Grant is a compelling higher education policy from the standpoint of external validity because it has clear analogs to tuition subsidies and tuition-free college programs in other settings. Our findings highlight an important example of how financial aid can increase political participation, even among middle income students who do not earn more money and are no more likely to attend college as a result of the program. Evidence from the Pell Grant helps externally validate these results with an alternative program that targets different students, subsidizes attendance at a different set of institutions, and impacts policy through a different margin of treatment.

Policymakers and economists have asserted for several centuries that public education expenditures enhance civic life by producing an informed citizenry and broader representation in the political process \citep{mann_republic_1957,friedman_capitalism_2020}. Our work shows that large-scale tuition subsidies can do both: inducing middle to low-income and college-educated youth to vote, with effects concentrated among high GPA students. Given that the Cal Grant's effects take place even absent any labor market returns, we conclude that the civic externalities of education expenditures are an important part of their marginal social benefits \citep{finkelstein_welfare_2020,hendren_unified_2020}.

% \section{General Conclusion} \label{sec:general-conclusion}
\chapter*{General Conclusion}
\addcontentsline{toc}{chapter}{General Conclusion}
\markboth{GENERAL CONCLUSION}{}
\label{sec:general-conclusion}

This dissertation has examined three fundamental questions about political institutions, information, and civic engagement. Through careful empirical analysis leveraging quasi-experimental variation, the essays provide new insights into how institutional design shapes political behavior and democratic participation.

\subsection*{Synthesis of Findings}

The three essays reveal both the possibilities and limits of institutional design in promoting democratic values. Chapter~1 finds systematic differences in disaster relief allocation correlated with political representation, but attributes these patterns to compositional differences in disaster vulnerability rather than intentional manipulation---suggesting that while bureaucratic institutions may constrain overt favoritism, underlying geographic and demographic factors produce systematic disparities that warrant policy attention. Chapter~2 highlights the importance of methodological rigor in evaluating institutional effects, demonstrating how traditional estimators can produce misleading conclusions about direct democracy's civic benefits. Chapter~3 shows that education subsidies can generate powerful civic externalities that operate through socialization and network effects.

\subsection*{Future Directions}

Several promising avenues for future research emerge from this work:

\begin{enumerate}
    \item Examining whether similar patterns hold in other institutional contexts and countries
    \item Exploring heterogeneous treatment effects across different populations and settings
    \item Investigating the specific mechanisms through which institutions affect political behavior
    \item Developing new methodological approaches for studying institutional effects
    \item Replicating/testing the findings in this dissertation using larger/different datasets
    \item Continuing to seek causally credible, real-world settings with greater potential for as-if/quasi random identifying variation
  \end{enumerate}

As democracies worldwide grapple with challenges of polarization, disinformation, and declining trust, understanding how institutions shape political behavior becomes ever more critical. The evidence presented here provides both reasons for optimism and important caveats about the limits of institutional reform.


\clearpage
\phantomsection
\addcontentsline{toc}{chapter}{References}

\begin{center}
\textbf{References}
\end{center}

\bibliographystyle{apsr_fs}
\bibliography{dissertation,policy_learning}


\appendix
\renewcommand{\chaptername}{Appendix}

\chapter{Supplementary Materials for Chapter 1: Disaster Relief}



\setcounter{table}{0}
\renewcommand{\thetable}{A\arabic{table}}
\setcounter{figure}{0}
\renewcommand{\thefigure}{A\arabic{figure}}

\section*{A. Data and Sample Description}
\addcontentsline{toc}{section}{A. Data and Sample Description}

This appendix provides detailed information on the data sources, variable definitions, and sample composition for the analysis of partisan favoritism in FEMA disaster relief allocation.

\begin{table}[H]
\centering
\caption{Summary Statistics: Main Analysis Variables}
\label{tab:summary_stats}
\begin{tabular}{lrrrrrrr}
\toprule
Variable & N & Mean & SD & Min & Median & Max \\
\midrule
Total FEMA Obligations (\$) & 3,872 & 19.10M & 132.35M & 90.64 & 2.31M & 3683.11M \\
Log(FEMA Obligations) & 3,872 & 14.518 & 2.251 & 4.518 & 14.654 & 22.027 \\
Running Variable (Dem Share $-$ 0.5) & 3,872 & 0.001 & 0.249 & $-$0.500 & $-$0.026 & 0.500 \\
Copartisan Running Variable & 3,872 & $-$0.005 & 0.249 & $-$0.500 & 0.011 & 0.500 \\
Distance from Threshold & 3,872 & 0.200 & 0.149 & 0.000 & 0.164 & 0.500 \\
\bottomrule
\end{tabular}
\caption*{\footnotesize Notes: Sample includes all CD-years with FEMA Public Assistance obligations, 2001--2020. Running variable is Democratic vote share minus 0.5. Copartisan running variable is positive when district representative shares president's party.}
\end{table}

\begin{table}[H]
\centering
\caption{Sample Composition by Presidential Administration}
\label{tab:sample_president}
\begin{tabular}{lcccccc}
\toprule
President & Years & N (CD-years) & Copartisan & Opposition & Mean (\$M) & Median (\$M) \\
\midrule
Bush & 2001--2008 & 1,638 & 869 & 769 & 11.07 & 1.30 \\
Obama & 2009--2016 & 1,315 & 619 & 696 & 16.23 & 2.46 \\
Trump & 2017--2020 & 919 & 493 & 426 & 37.50 & 7.66 \\
\midrule
Total & 2001--2020 & 3,872 & 1,981 & 1,891 & 19.10 & 2.31 \\
\bottomrule
\end{tabular}
\caption*{\footnotesize Notes: Copartisan districts are those where the representative shares the president's party. Bush years: Republican president, 2001--2008. Obama years: Democratic president, 2009--2016. Trump years: Republican president, 2017--2020.}
\end{table}

\begin{table}[H]
\centering
\caption{Sample Composition by Census Region}
\label{tab:sample_region}
\begin{tabular}{lccccc}
\toprule
Region & N (CD-years) & States & Districts & Mean (\$M) & \% Copartisan \\
\midrule
South & 1,470 & 15 & 155 & 15.86 & 53.0\% \\
Midwest & 869 & 12 & 100 & 8.71 & 55.1\% \\
Northeast & 839 & 9 & 83 & 42.39 & 51.0\% \\
West & 694 & 12 & 101 & 10.79 & 42.5\% \\
\bottomrule
\end{tabular}
\caption*{\footnotesize Notes: Regions follow Census Bureau definitions. Northeast: CT, ME, MA, NH, RI, VT, NJ, NY, PA. Midwest: IL, IN, MI, OH, WI, IA, KS, MN, MO, NE, ND, SD. South: DE, FL, GA, MD, NC, SC, VA, DC, WV, AL, KY, MS, TN, AR, LA, OK, TX. West: AZ, CO, ID, MT, NV, NM, UT, WY, AK, CA, HI, OR, WA.}
\end{table}

\begin{table}[H]
\centering
\caption{Variable Definitions}
\label{tab:variable_definitions}
\begin{tabular}{lp{7.5cm}l}
\toprule
Variable & Definition & Source \\
\midrule
FEMA Obligations & Total FEMA Public Assistance obligations to congressional district in fiscal year, weighted by county-CD population allocation & FEMA EMMIE (FOIA) \\[0.5ex]
Log(FEMA Obligations) & Natural log of FEMA Obligations & Computed \\[0.5ex]
Running Variable (rv1) & Democratic two-party vote share in most recent House election minus 0.5 & MIT Election Data Lab \\[0.5ex]
Copartisan RV & rv1 $\times$ (+1 if Dem president, $-$1 if GOP president) & Computed \\[0.5ex]
Copartisan & District representative shares president's party (rv1\_copartisan $>$ 0) & Computed \\[0.5ex]
Opposition & District representative opposes president's party (rv1\_copartisan $<$ 0) & Computed \\[0.5ex]
\bottomrule
\end{tabular}
\end{table}

% =============================================================================
% APPENDIX B: RD VALIDITY AND DIAGNOSTICS
% =============================================================================

\section*{B. RD Validity and Diagnostics}
\addcontentsline{toc}{section}{B. RD Validity and Diagnostics}

A key assumption of regression discontinuity designs is that units cannot precisely manipulate their position relative to the threshold. In this context, this requires that congressional districts cannot precisely control whether their representative shares the president's party. This assumption is plausible given that House elections are determined by local factors largely independent of presidential party.

Figure \ref{fig:rv_distribution} shows the distribution of the copartisan running variable. The smooth distribution around the threshold, with no apparent bunching at zero, supports the no-manipulation assumption. A formal McCrary density test yields $p = 0.150$ for the Obama subsample, failing to reject the null of no manipulation.

\begin{figure}[H]
\centering
\includegraphics[width=0.85\textwidth]{\basepath/Research/events_extreme_candidate_entry/diss/code/prod/output/plots/descriptives/figure_rv_distribution.pdf}
\caption{Distribution of Copartisan Running Variable}
\label{fig:rv_distribution}
\caption*{\footnotesize Notes: Histogram shows the distribution of the copartisan running variable across all CD-year observations. Positive values indicate districts where the representative shares the president's party (copartisan); negative values indicate opposition districts. The vertical dashed line marks the threshold at zero. The smooth distribution around the threshold supports the no-manipulation assumption required for RD validity.}
\end{figure}

Figure \ref{fig:rv_by_president} shows how the running variable distribution differs across presidential administrations. Under Democratic presidents (Obama), the transformation $\text{rv1\_copartisan} = \text{rv1} \times (+1)$ preserves the original vote share margin. Under Republican presidents (Bush, Trump), the transformation $\text{rv1\_copartisan} = \text{rv1} \times (-1)$ flips the sign, ensuring that positive values consistently indicate copartisan status regardless of which party holds the presidency.

\begin{figure}[H]
\centering
\includegraphics[width=0.85\textwidth]{\basepath/Research/events_extreme_candidate_entry/diss/code/prod/output/plots/descriptives/figure_rv_by_president.pdf}
\caption{Running Variable Distribution by Presidential Administration}
\label{fig:rv_by_president}
\caption*{\footnotesize Notes: Density plots show the distribution of the copartisan running variable separately for each presidential administration. The transformation ensures that positive values indicate copartisan districts and negative values indicate opposition districts under all three presidents.}
\end{figure}

% =============================================================================
% APPENDIX C: DESCRIPTIVE PATTERNS
% =============================================================================

\section*{C. Descriptive Patterns in FEMA Spending}
\addcontentsline{toc}{section}{C. Descriptive Patterns in FEMA Spending}

This appendix presents descriptive patterns in FEMA Public Assistance spending over the 2001--2020 study period.

\begin{figure}[H]
\centering
\includegraphics[width=0.95\textwidth]{\basepath/Research/events_extreme_candidate_entry/diss/code/prod/output/plots/descriptives/figure_time_series_total.pdf}
\caption{Total FEMA Public Assistance Obligations by Year}
\label{fig:time_series_total}
\caption*{\footnotesize Notes: Bar chart shows total FEMA Public Assistance obligations (in billions of dollars) by fiscal year, colored by presidential administration. Notable spikes correspond to major disaster years: 2005 (Hurricane Katrina), 2012 (Hurricane Sandy), and 2017 (Hurricanes Harvey, Irma, and Maria).}
\end{figure}

\begin{figure}[H]
\centering
\includegraphics[width=0.95\textwidth]{\basepath/Research/events_extreme_candidate_entry/diss/code/prod/output/plots/descriptives/figure_time_series_n.pdf}
\caption{Number of Congressional Districts Receiving FEMA Assistance by Year}
\label{fig:time_series_n}
\caption*{\footnotesize Notes: Bar chart shows the number of congressional districts receiving FEMA Public Assistance obligations in each fiscal year. Variation reflects both the geographic scope of disasters and the intensity of FEMA response in different years.}
\end{figure}

\begin{figure}[H]
\centering
\includegraphics[width=0.85\textwidth]{\basepath/Research/events_extreme_candidate_entry/diss/code/prod/output/plots/descriptives/figure_outcome_distribution.pdf}
\caption{Distribution of Log FEMA Obligations}
\label{fig:outcome_distribution}
\caption*{\footnotesize Notes: Histogram shows the distribution of log FEMA Public Assistance obligations across all CD-year observations in the sample. The approximately normal distribution of the log-transformed outcome supports the use of log obligations as the primary dependent variable.}
\end{figure}

\begin{figure}[H]
\centering
\includegraphics[width=0.85\textwidth]{\basepath/Research/events_extreme_candidate_entry/diss/code/prod/output/plots/descriptives/figure_outcome_by_status.pdf}
\caption{FEMA Obligations by Copartisan Status}
\label{fig:outcome_by_status}
\caption*{\footnotesize Notes: Density plots compare the distribution of log FEMA obligations for copartisan districts (where the representative shares the president's party) versus opposition districts. The substantial overlap indicates that the effect operates at the margin rather than representing a wholesale shift in distributions.}
\end{figure}

% =============================================================================
% APPENDIX D: ALTERNATIVE OUTCOMES AND MECHANISM TESTS
% =============================================================================

\section*{D. Alternative Outcomes and Mechanism Tests}
\addcontentsline{toc}{section}{D. Alternative Outcomes and Mechanism Tests}

To investigate potential mechanisms underlying the negative copartisan effect on FEMA obligations, I examine whether copartisan status affects dimensions of disaster relief other than the total amount. If the effect reflects intentional political manipulation, we might expect differences in the \textit{composition} of relief (more visible vs.\ infrastructure spending), the \textit{speed} of processing (faster obligation for favored districts), or the \textit{structure} of awards (more or larger projects). Null results on these alternative outcomes would suggest the effect operates through which districts receive relief rather than how FEMA processes claims differently based on partisanship.

Table \ref{tab:additional_outcomes} presents RD estimates for 17 alternative outcome measures across five categories: (1) within-district funding inequality (Gini coefficient, coefficient of variation, max/min ratio across counties); (2) funding mix (share of highly visible vs.\ infrastructure spending); (3) relief-to-damage ratio (FEMA obligations relative to NOAA-estimated damage); (4) processing speed (days from declaration to obligation); and (5) project structure (number, mean size, and concentration of individual awards).

\begin{table}[H]
\centering
\caption{Alternative Outcome RD Estimates: Mechanism Tests}
\label{tab:additional_outcomes}
\small
\begin{tabular}{lcccc}
\toprule
Outcome & Estimate & SE & p-value & N$_{\text{eff}}$ \\
\midrule
\multicolumn{5}{l}{\textit{Panel A: Within-District Inequality}} \\
\quad Gini Coefficient & $-$0.026 & (0.026) & 0.322 & 1,228 \\
\quad Coefficient of Variation & $-$0.043 & (0.074) & 0.561 & 1,061 \\
\quad Log(Max/Min Ratio) & $-$31.8\% & (24.8\%) & 0.198 & 937 \\[0.5ex]
\multicolumn{5}{l}{\textit{Panel B: Funding Mix}} \\
\quad Share: Highly Visible (A+B) & 0.011 & (0.012) & 0.359 & 10,841 \\
\quad Share: Infrastructure (C--G) & $-$0.010 & (0.008) & 0.230 & 11,457 \\
\quad Log(Infrastructure/Visible) & $-$19.0\% & (36.5\%) & 0.604 & 9,118 \\[0.5ex]
\multicolumn{5}{l}{\textit{Panel C: Relief-to-Damage Ratio}} \\
\quad Relief/Damage (level) & $-$5.539 & (4.11) & 0.178 & 896 \\
\quad Log(Relief/Damage) & $-$37.3\% & (30.1\%) & 0.217 & 1,013 \\[0.5ex]
\multicolumn{5}{l}{\textit{Panel D: Processing Speed}} \\
\quad Mean Days to Obligation & 10.694 & (40.8) & 0.793 & 1,392 \\
\quad Median Days to Obligation & 17.341 & (32.1) & 0.590 & 1,426 \\
\quad Share Obligated $\leq$ 30 Days & $-$0.003 & (0.004) & 0.439 & 1,091 \\
\quad Share Obligated $\leq$ 60 Days & 0.000 & (0.006) & 0.971 & 1,591 \\
\quad Share Obligated $\leq$ 90 Days & $-$0.018 & (0.035) & 0.602 & 1,528 \\[0.5ex]
\multicolumn{5}{l}{\textit{Panel E: Project Structure}} \\
\quad Log(Number of Projects) & $-$22.9\% & (15.9\%) & 0.150 & 836 \\
\quad Log(Mean Project Size) & 18.7\% & (33.3\%) & 0.574 & 1,335 \\
\quad Log(Median Project Size) & 15.3\% & (20.9\%) & 0.464 & 1,156 \\
\quad Share from Large Projects & 0.001 & (0.004) & 0.804 & 1,897 \\
\bottomrule
\end{tabular}
\caption*{\footnotesize Notes: Regression discontinuity estimates using copartisan running variable. Local linear estimation with triangular kernel and MSE-optimal bandwidth selection. Robust standard errors in parentheses. Percent effects shown for log outcomes; coefficients shown for level outcomes. N$_{\text{eff}}$ is the effective sample size within the RD bandwidth. Categories A--B are Debris Removal and Emergency Protective Measures (highly visible); Categories C--G are Roads/Bridges, Water Control, Public Buildings, Utilities, and Recreation (infrastructure). $^{*}p<0.10$, $^{**}p<0.05$, $^{***}p<0.01$.}
\end{table}

All 17 alternative outcome estimates are statistically insignificant at conventional levels. The null results on processing speed are particularly informative: copartisan districts do not receive faster FEMA service, ruling out administrative favoritism as a mechanism. The null results on funding composition suggest that FEMA does not strategically allocate more visible spending (which might generate political credit) to favored districts.

These findings are consistent with a compositional interpretation of the main negative effect. The $-45\%$ difference in total FEMA obligations between copartisan and opposition districts does not appear to reflect how FEMA treats claims differently based on partisanship. Instead, the effect likely reflects systematic differences in which types of districts---with different disaster exposure and relief application patterns---happen to be represented by copartisan versus opposition members at any given time.

% =============================================================================
% APPENDIX E: RUNNING VARIABLE CONSTRUCTION (formerly A)
% =============================================================================

\section*{E. Running Variable Construction}
\addcontentsline{toc}{section}{E. Running Variable Construction}

A critical methodological issue in this analysis is constructing the running variable to properly identify copartisan effects when pooling data across Democratic and Republican administrations. This issue is subtle but consequential: using the wrong construction can produce spurious null results by conflating opposite-signed effects.

The naive approach uses the Democratic vote margin directly:
\begin{equation}
\text{rv1} = \text{dem\_share} - 0.5
\end{equation}

This approach is problematic because $\text{rv1} > 0$ means different things under different presidents. Under Democratic presidents, $\text{rv1} > 0$ indicates a Democratic representative who is copartisan with the president. Under Republican presidents, $\text{rv1} > 0$ indicates a Democratic representative who is in opposition to the president. When pooling across administrations, this conflation produces opposite-signed effects that partially cancel, potentially masking substantial partisan effects.

The magnitude of this problem is illustrated by comparing estimates across subsamples:

\begin{center}
\begin{tabular}{lccc}
\toprule
Period & Effect Using rv1 & Interpretation & Effect Using rv1\_copartisan \\
\midrule
Obama & $-61.8\%$ & Copartisan effect & $-61.8\%$ \\
Bush/Trump & $+68.0\%$ & Opposition effect & $-40.5\%$ \\
Pooled & $-2.0\%$ & Cancels out & $-50.7\%$ \\
\bottomrule
\end{tabular}
\end{center}

The solution is to transform the running variable so that positive values always indicate copartisan representation:
\begin{equation}
\text{rv1\_copartisan} = \text{rv1} \times \text{prez\_dem\_indicator}
\end{equation}
where $\text{prez\_dem\_indicator} = +1$ under Democratic presidents and $-1$ under Republican presidents.

This transformation ensures consistent interpretation: $\text{rv1\_copartisan} > 0$ always means the district's representative shares the president's party, regardless of which party that is. All estimates in this paper use this transformed running variable.

\section*{F. Robust Inference}
\addcontentsline{toc}{section}{F. Robust Inference}

All regression discontinuity estimates use the robust bias-corrected inference procedure from \cite{calonico2014robust}. This methodological choice is important because standard RD estimation faces a bias-variance tradeoff: larger bandwidths reduce variance but introduce bias from curvature in the conditional expectation function. The \cite{calonico2014robust} procedure addresses this by using bias correction while appropriately adjusting standard errors to account for the additional uncertainty introduced by estimating the bias.

The rdrobust package produces three rows of output for each estimate: conventional (which ignores bias), bias-corrected (which corrects the point estimate but uses conventional standard errors), and robust (which both corrects the point estimate and inflates standard errors to account for bias estimation). I report the robust inference (third row) throughout, as this provides valid coverage under the maintained assumptions. Specifically, for each estimate I report standard errors from \texttt{rd\$se[3]}, p-values from \texttt{rd\$pv[3]}, and confidence intervals from \texttt{rd\$ci[3,]}.

This choice follows best practices in the RD literature. Using conventional inference (the first row) would understate uncertainty; using bias-corrected point estimates with conventional standard errors (the second row) would provide invalid coverage. The robust procedure ensures that reported confidence intervals have appropriate coverage properties under local linear estimation with bias correction.

\section*{G. Balance Tests}
\addcontentsline{toc}{section}{G. Balance Tests}

The validity of the regression discontinuity design rests on the assumption that potential outcomes vary smoothly through the threshold, which implies that predetermined covariates should be balanced across the cutoff. Table \ref{tab:balance_full} presents balance tests for all available predetermined covariates, including demographic characteristics from the Census, prior disaster history, and geographic variables.

For each covariate, I estimate the same local linear specification used in the main analysis but with the covariate as the outcome. A significant discontinuity would suggest that districts just above and below the threshold differ systematically in ways that could confound the main estimates. The table reports the estimated discontinuity, robust standard error, and p-value for each covariate, as well as a joint F-test for balance across all covariates.

The results support the validity of the design: no individual covariate shows a statistically significant discontinuity at conventional levels, and the joint test fails to reject balance. This provides reassurance that the comparison of copartisan and opposition districts near the threshold approximates a randomized experiment with respect to observable characteristics. Of course, balance on observables does not guarantee balance on unobservables, but it is a necessary condition for the RD identifying assumptions and increases confidence in the causal interpretation.

\input{\basepath/Research/events_extreme_candidate_entry/diss/code/prod/output/tables/appendix/table_balance_tests.tex}
\label{tab:balance_full}

\section*{H. Placebo Threshold Tests}
\addcontentsline{toc}{section}{H. Placebo Threshold Tests}

A key identifying assumption in regression discontinuity designs is that potential outcomes vary smoothly through the threshold. If this assumption holds, we should observe discontinuities only at the true threshold (where treatment status changes) and not at placebo thresholds elsewhere in the running variable distribution. Significant effects at placebo thresholds would suggest either misspecification of the functional form or the presence of other discontinuities unrelated to treatment.

Table \ref{tab:placebo_thresholds} tests for discontinuities at placebo thresholds located at various points away from zero. I estimate the standard RD specification but centered at these artificial thresholds rather than the true threshold. The placebo thresholds are chosen to span the range of the running variable where data density is sufficient for estimation.

The results provide qualified support for the identification strategy. While three of the four placebo tests are statistically insignificant (p $>$ 0.07), the test at $c = -0.10$ shows a significant effect ($+94.8\%$, p = 0.019). However, this result does not undermine the main finding for two reasons. First, the direction differs markedly: the true threshold shows a negative effect ($-49.2\%$) while the significant placebo shows a positive effect, indicating these discontinuities capture different phenomena. Second, a single significant result among multiple tests is consistent with expected false positive rates. The pattern overall suggests that the negative copartisan effect is specific to the true threshold rather than an artifact of functional form misspecification.

\input{\basepath/Research/events_extreme_candidate_entry/diss/code/prod/output/tables/appendix/table_placebo_thresholds.tex}
\label{tab:placebo_thresholds}

\section*{I. Electoral Cycle Analysis}
\addcontentsline{toc}{section}{I. Electoral Cycle Analysis}

Political incentives may vary systematically across the electoral cycle. If presidents strategically allocate disaster relief to influence elections, we would expect stronger effects in election years when the electoral stakes are highest. Conversely, if the pattern reflects compositional differences rather than strategic behavior, effects should be relatively stable across the electoral cycle.

Table \ref{tab:electoral_cycle_full} presents comprehensive estimates by year within presidential terms. Panel A shows estimates for each year of the four-year presidential term separately. Panel B distinguishes midterm election years from other years. Panel C focuses specifically on presidential election years. Panel D examines position in the presidential term (first term vs. last year vs. middle years).

The results provide no support for the electoral targeting hypothesis. Effects are not systematically larger in election years; if anything, the pattern is reversed, with the strongest effects appearing in non-election years. This finding is important for interpretation: it undermines the hypothesis that presidents strategically direct disaster relief to opposition districts for electoral gain, and is more consistent with the compositional explanation developed in the main text.

\input{\basepath/Research/events_extreme_candidate_entry/diss/code/prod/output/tables/appendix/table7_electoral_cycle_robustness.tex}
% Label provided by input file

\section*{J. Bandwidth Robustness}
\addcontentsline{toc}{section}{J. Bandwidth Robustness}

The choice of bandwidth in regression discontinuity estimation involves a bias-variance tradeoff: larger bandwidths include more observations (reducing variance) but may introduce bias if the conditional expectation function is not well-approximated by a linear function over the wider range. While the main analysis uses the MSE-optimal bandwidth from \cite{calonico2014robust}, it is important to verify that results are not sensitive to this choice.

Table \ref{tab:bandwidth_robustness} presents estimates across a range of bandwidths, from half the optimal bandwidth to double the optimal bandwidth. Narrower bandwidths focus on observations closest to the threshold where the local linear approximation is most credible but have fewer observations and thus wider confidence intervals. Wider bandwidths include more observations but may introduce curvature bias.

The results demonstrate that the negative copartisan effect is robust to bandwidth choice. Point estimates remain consistently negative across all bandwidths, ranging from approximately $-52\%$ at the narrowest bandwidth to $-38\%$ at the widest. The consistency of sign and approximate magnitude across this range provides confidence that the finding is not an artifact of a particular bandwidth selection. Statistical significance varies with bandwidth due to the bias-variance tradeoff, but the qualitative conclusion---that copartisan districts receive less disaster relief---is stable.

\input{\basepath/Research/events_extreme_candidate_entry/diss/code/prod/output/tables/appendix/table6_robustness_FINAL.tex}
\label{tab:bandwidth_robustness}

\section*{K. Detailed Demographic Heterogeneity}
\addcontentsline{toc}{section}{K. Detailed Demographic Heterogeneity}

The main text documents substantial heterogeneity in the copartisan effect across district demographic characteristics, with the strongest negative effects concentrated in predominantly white, affluent, highly-educated districts. Table \ref{tab:demographic_heterogeneity} in the main text provides detail on these patterns.

The table presents RD estimates separately for districts above and below the median on each demographic dimension, as well as estimates by quintile for key variables. This granular analysis reveals whether heterogeneity represents a sharp threshold (with effects appearing only above or below some cutoff) or a more gradual gradient (with effects varying smoothly across the distribution).

For racial diversity, the gradient is particularly striking and approximately monotonic: effects are most negative in the whitest quintile and attenuate progressively through the distribution, eventually reversing sign (though imprecisely estimated) in the most diverse quintile. This gradient is difficult to explain through strategic political targeting but is consistent with compositional differences---if disaster-prone areas with particular infrastructure characteristics differ in their racial composition, and if this composition correlates with political representation, we would observe exactly this pattern.

The demographic patterns are also informative for understanding potential mechanisms. The concentration of effects in high-income, low-poverty, high-education districts is consistent with a story in which these areas have different disaster experiences or relief needs. These characteristics may proxy for housing quality, infrastructure investment, or other factors that affect disaster vulnerability and relief requirements.

% Table already appears in main text as Table \ref{tab:demographic_heterogeneity}
% Removed duplicate inclusion to avoid multiply-defined label warning

\section*{L. Causal Forest Details}
\addcontentsline{toc}{section}{L. Causal Forest Details}

Causal forests \citep{wager2018estimation, athey2016recursive} provide a data-driven approach to discovering treatment effect heterogeneity. Unlike the hypothesis-driven heterogeneity analyses in the main text (which examine specific dimensions suggested by theory), causal forests allow the data to reveal which covariates, if any, predict variation in treatment effects. This provides a complementary test of whether strategic manipulation (which should produce heterogeneity by electoral context) or compositional differences (which should produce homogeneity across electoral contexts but potential heterogeneity by compositional factors) better explains the observed pattern.

The causal forest is estimated using the grf package with specifications chosen to balance bias and variance while maintaining honesty (using separate subsamples for tree construction and estimation). The sample is restricted to observations within bandwidth 0.10 of the threshold (N = 1,044) to focus on the local comparison that identifies causal effects. Treatment is defined as an indicator for the copartisan running variable being non-negative. Covariates include year, state, the running variable itself, year in presidential term, and presidential party. Overlap weighting addresses extreme propensity scores that could otherwise distort estimates. The forest uses 2,000 trees with honest splitting to ensure valid inference.

The calibration test confirms the forest is well-specified: the mean forest prediction coefficient is statistically significant (p = 0.006), indicating that the forest captures meaningful variation in outcomes. However---and this is the key finding---the distribution of conditional average treatment effects (CATEs) is homogeneous. The standard deviation of CATEs is only 4.6 percentage points, with all estimates falling in the narrow range from $-89\%$ to $-77\%$. When examining CATEs by subgroup (administration, election year, etc.), the estimates are virtually identical.

This homogeneity is inconsistent with strategic manipulation, which should produce heterogeneous effects varying with electoral incentives. The finding instead supports the compositional interpretation: whatever produces the negative effect operates uniformly across observable contexts.

\input{\basepath/Research/events_extreme_candidate_entry/diss/code/prod/output/causal_forest/tables/table_causal_forest.tex}
\chapter{Supplementary Materials for Chapter 2: Direct Democracy}



\setcounter{table}{0}
\renewcommand{\thetable}{B\arabic{table}}
\setcounter{figure}{0}
\renewcommand{\thefigure}{B\arabic{figure}}

\subsection*{A. Complete Model Specifications}
\addcontentsline{toc}{section}{A. Complete Model Specifications}

Table \ref{tab:full_specs} presents the full regression output for all main specifications comparing traditional two-way fixed effects with modern difference-in-differences estimators. This table includes coefficients for the morality politics treatment indicator (marijuana, gambling, abortion, and marriage measures), individual-level controls (age, age squared, gender, education, party affiliation), state and year fixed effects, and clustered standard errors. The divergence between TWFE and modern estimators illustrates how heterogeneous treatment timing can severely bias traditional estimates when treatment effects vary across adoption cohorts.

\begin{table}[H]
\centering
\caption{Full Regression Specifications}
\label{tab:full_specs}
\begin{singlespace}
{\footnotesize \textit{Notes:} * p $<$ 0.1, ** p $<$ 0.05, *** p $<$ 0.01. Standard errors clustered by state in parentheses.}
\end{singlespace}
\vspace{0.5em}
\small
\begin{tabular}{lccccc}
\toprule
& Simple DiD & State FE & Two-way FE & With Controls & State Trends \\ \midrule
(Intercept) & 3.277*** &  &  &  &  \\
& (0.002) &  &  &  &  \\
Ballot Measure $\times$ Post & $-$0.084*** & 0.102*** & 0.029* & 0.013 & 0.026* \\
& (0.014) & (0.033) & (0.015) & (0.011) & (0.015) \\
Treatment Group & 0.169*** &  &  &  &  \\
& (0.005) &  &  &  &  \\
Post-Election & 0.049*** &  &  &  &  \\
& (0.003) &  &  &  &  \\
age &  &  &  & 0.024*** &  \\
&  &  &  & (0.001) &  \\
age\_squared &  &  &  & $-$0.000*** &  \\
&  &  &  & (0.000) &  \\
female &  &  &  & $-$0.322*** &  \\
&  &  &  & (0.007) &  \\
college &  &  &  & 0.310*** &  \\
&  &  &  & (0.008) &  \\
democrat &  &  &  & 0.394*** &  \\
&  &  &  & (0.009) &  \\
republican &  &  &  & 0.330*** &  \\
&  &  &  & (0.010) &  \\
independent &  &  &  & 0.414*** &  \\
&  &  &  & (0.009) &  \\ \midrule
Num.Obs. & 477,880 & 477,880 & 477,880 & 477,717 & 477,880 \\
R2 & 0.003 & 0.003 & 0.021 & 0.206 & 0.021 \\
R2 Adj. & 0.003 & 0.003 & 0.021 & 0.206 & 0.021 \\
R2 Within &  & 0.000 & 0.000 & 0.189 & 0.000 \\
R2 Within Adj. &  & 0.000 & 0.000 & 0.189 & 0.000 \\
F & 509.439 &  &  &  &  \\
Std.Errors &  & by: state & by: state & by: state & by: state \\
FE: state &  & X & X & X & X \\
FE: year &  &  & X & X & X \\
\bottomrule
\end{tabular}
\end{table}

\subsection*{B. Sun-Abraham Event Study Decomposition}
\addcontentsline{toc}{section}{B. Sun-Abraham Event Study Decomposition}

Table \ref{tab:sun_abraham} reports the complete Sun-Abraham (2021) cohort-specific average treatment effects for each relative time period. These estimates properly account for heterogeneous treatment effects across adoption cohorts by avoiding the problematic comparisons inherent in TWFE. The near-zero aggregate ATT emerges from averaging heterogeneous effects across early and late adopters, demonstrating why traditional methods that pool these different treatment responses produce misleading results. The table shows effects for periods from three years before to three years after initial treatment.

\input{\basepath/Research/policy_learning/output/tables/appendix/table_a2_sun_abraham.tex}
\label{tab:sun_abraham}

\subsection*{C. Parallel Trends Diagnostic Tests}
\addcontentsline{toc}{section}{C. Parallel Trends Diagnostic Tests}

To formally assess the parallel trends assumption, I estimate a model interacting treatment group status with a linear time trend in the pre-treatment period. This test examines whether treatment and control states followed similar trajectories before morality politics measures appeared. The interaction coefficient is $-0.031$ (SE = 0.003, p $<$ 0.001), indicating a statistically significant violation of parallel trends: treatment states were already on different trajectories before ballot measures appeared. This violation of the key identifying assumption helps explain discrepancies between traditional TWFE and modern estimator results, and reinforces concerns about the credibility of positive TWFE estimates.

\subsection*{D. Heterogeneous Treatment Effects Analysis}
\addcontentsline{toc}{section}{D. Heterogeneous Treatment Effects Analysis}

Table \ref{tab:heterogeneity} provides comprehensive heterogeneity analysis examining how treatment effects vary across demographic subgroups and political characteristics. The table reports separate treatment effects and interaction terms for education levels (college vs. non-college), party affiliation (Democrat, Republican, Independent), age cohorts (young, middle-aged, elderly), and baseline political awareness terciles. While TWFE suggests differential effects across groups, these patterns likely inherit the same biases as the main estimates and should be interpreted cautiously given the findings from modern estimators.

\input{\basepath/Research/policy_learning/output/tables/appendix/table_a5_heterogeneity.tex}
\label{tab:heterogeneity}

\subsection*{E. Leave-One-State-Out Sensitivity Analysis}
\addcontentsline{toc}{section}{E. Leave-One-State-Out Sensitivity Analysis}

Table \ref{tab:leave_one_out} documents how the estimated marijuana effect—the only marginally significant finding in traditional specifications—changes when individual states are systematically excluded from the analysis. The table reports the coefficient and standard error when each state is dropped, revealing sensitivity to specific states, particularly early-adopting states with established direct democracy traditions. This sensitivity analysis, combined with the loss of significance under alternative specifications, raises questions about the robustness of even this limited finding.

\input{\basepath/Research/policy_learning/output/tables/appendix/table_a6_leave_one_out.tex}
\label{tab:leave_one_out}

\subsection*{F. Effects by Specific Morality Politics Topics}
\addcontentsline{toc}{section}{F. Effects by Specific Morality Politics Topics}

Table \ref{tab:topics} disaggregates the treatment effect by type of morality politics measure, presenting separate estimates for marijuana legalization, gambling expansion, abortion restrictions, and marriage equality ballot measures. The table reveals striking heterogeneity across issue areas, with some measures showing much larger effects than others despite similar theoretical predictions. This heterogeneity suggests that pooling different morality politics measures may mask important variation and that outlier effects from specific issues may drive aggregate estimates.

\input{\basepath/Research/policy_learning/output/tables/appendix/table_a7_topics.tex}
\label{tab:topics}

\subsection*{G. Evolution of Sample Composition and Balance}
\addcontentsline{toc}{section}{G. Evolution of Sample Composition and Balance}

Table \ref{tab:summary_stats} tracks summary statistics by treatment status across the study period, documenting changes in sample composition, covariate balance, and the distribution of morality politics measures over time. The table shows how treated and control units differ on observable characteristics including education, party affiliation, and baseline political interest. It also reveals temporal patterns in the types of morality politics measures appearing on ballots, with certain issues concentrated in specific time periods, which contributes to the heterogeneous timing effects identified by modern estimators.

\input{\basepath/Research/policy_learning/output/tables/appendix/table_a8_summary_stats.tex}

\subsection*{H. Sample Construction and Attrition}
\addcontentsline{toc}{section}{H. Sample Construction and Attrition}

Table \ref{tab:sample_selection} details the sample selection process from the initial Cooperative Election Study cumulative file through the final analysis dataset. The table documents attrition at each stage of data cleaning: dropping observations with missing state identifiers, excluding non-continental states, removing cases with missing outcome data, and restricting to the study period. Understanding sample construction is crucial for assessing external validity and potential selection biases. The final sample represents the vast majority of eligible CES respondents during the study period.

\input{\basepath/Research/policy_learning/output/tables/appendix/table_a9_sample_selection.tex}
\label{tab:sample_selection}

\subsection*{I. Variable Definitions and Measurement}
\addcontentsline{toc}{section}{I. Variable Definitions and Measurement}

Table \ref{tab:variables} provides precise definitions for all variables used in the analysis, including construction details for composite measures. Key variables include the news interest score (measured on a four-point scale), treatment indicators (both binary and continuous measures), timing variables for staggered adoption, and the political awareness index combining multiple indicators of political sophistication. The table also documents the classification criteria for identifying morality politics measures from ballot descriptions and the coding decisions for control variables.

\input{\basepath/Research/policy_learning/output/tables/appendix/table_a10_variables.tex}
\label{tab:variables}

\subsection*{J. Wild Cluster Bootstrap Inference}
\addcontentsline{toc}{section}{J. Wild Cluster Bootstrap Inference}

Table \ref{tab:bootstrap} addresses concerns about inference with a limited number of clusters by reporting wild cluster bootstrap p-values for all main specifications. The bootstrap procedure uses Webb weights to improve small-sample performance. These alternative inference methods help validate which findings are robust to different assumptions about the error structure and confirm that most effects lose statistical significance under more conservative inference approaches, reinforcing concerns about the fragility of traditional TWFE results.

\input{\basepath/Research/policy_learning/output/tables/appendix/table_a11_bootstrap.tex}
\label{tab:bootstrap}

\subsection*{K. Treatment Adoption Timing Across States}
\addcontentsline{toc}{section}{K. Treatment Adoption Timing Across States}

Treatment timing based on first morality politics ballot measure (marijuana, gambling, abortion, or same-sex marriage) appearing on state ballot during the 2006--2020 sample period reveals distinct waves of adoption. \textbf{Early adopters (2006, N=17)}: AL, AZ, CA, CO, FL, LA, MA, MO, NE, NV, OH, OK, OR, RI, SD, VA, WI. \textbf{Mid-period adopters (2007--2012, N=10)}: ME (2007); AR, MD, MI, ND (2008); WA (2010); MS, NJ (2011); MT, NC (2012). \textbf{Late adopter (2013, N=1)}: NY. \textbf{Never-treated states (N=19)}: CT, DE, GA, HI, ID, IN, IA, KS, KY, MN, NH, PA, SC, TN, TX, UT, VT, WV, WY. Never-treated states provide the clean control group essential for modern difference-in-differences estimators. This staggered adoption pattern, combined with effect heterogeneity across cohorts, explains why TWFE produces biased estimates and why modern estimators are necessary for credible identification.

\chapter{Supplementary Materials for Chapter 3: Cal Grant}


\setcounter{table}{0}
\renewcommand{\thetable}{C\arabic{table}}
\setcounter{figure}{0}
\renewcommand{\thefigure}{C\arabic{figure}}

\subsection*{A. Tables}\label{sec:tab}
    \addcontentsline{toc}{section}{A. Tables}



\begin{table}[h!]\centering
        \def\sym#1{\ifmmode^{#1}\else\(^{#1}\)\fi}
        \caption{Estimated Impact of College Tuition Subsidies on Political Outcomes\label{tab:mainresults}}
        \begin{tabular}{l*{6}{c}}
            \hline\hline
            &\multicolumn{1}{c}{(1)}&\multicolumn{1}{c}{(2)}&\multicolumn{1}{c}{(3)}&\multicolumn{1}{c}{(4)}&\multicolumn{1}{c}{(5)}&\multicolumn{1}{c}{(6)}\\
            Outcome&\multicolumn{1}{c}{Registered}&\multicolumn{1}{c}{2020 Vote}&\multicolumn{1}{c}{Turnout}&\multicolumn{1}{c}{Ever Voted}&\multicolumn{1}{c}{ (D $\cup$ I) x Vote}&\multicolumn{1}{c}{R x Vote}\\
            \hline
            &&&&&&\\
            \multicolumn{7}{l}{\textit{A. Estimated Impact of Receiving a 2017-2019 Cal Grant}} \\ [0.5em]
            \input{\basepath/Research/college_politics/Tables/row_cal1719.tex}
            Mean&[0.6914]&[0.5661]&[0.4884]&[0.5819]&[0.4152]&[0.0731]\\
            N      &   258,329     &    258,329      &      258,329      &   258,329    &   258,329      &   258,329      \\ [0.5em]

            \multicolumn{7}{l}{\textit{B. Estimated Impact of Receiving a 2010-2019 Cal Grant}} \\ [0.5em]

            \input{\basepath/Research/college_politics/Tables/row_cal1019.tex}
            Mean&[0.6292]&[0.5111]&[0.3863]&[0.5515]&[0.3294]&[0.0569]\\
            N& 738,046     &    738,046      &      738,046      &   738,046 &738,046      &   738,046\\ [0.5em]


            \multicolumn{7}{l}{\textit{C. Estimated Impact of Increasing 2010-2019 Pell Grant Generosity}} \\ [0.5em]

            \input{\basepath/Research/college_politics/Tables/row_pell1019.tex}
            Mean&[0.5894]&[0.4372]&[0.3256]&[0.4802]&[0.2906]&[0.0349]\\
            N&2,561,537    &    2,561,537      &      2,561,537     &   2,561,537 & 2,561,537     &   2,561,537 \\


            &&&&&&\\

            \hline
            Bandwidth & \$10,000       &   \$10,000        &   \$10,000        &   \$10,000       &   \$10,000        &   \$10,000        \\
            Kernel&   Uniform      &   Uniform        &   Uniform        &   Uniform       &   Uniform       &   Uniform       \\
            Polynomial&   1       &   1        &   1        &   1       &   1       &   1       \\
            Controls&   No       &   No        &   No        &   No        &   No        &   No        \\
            \hline
        \end{tabular}
        \caption*{\footnotesize Note: \sym{+} \(p<0.1\), \sym{*} \(p<0.05\), \sym{**} \(p<0.01\). Heteroskedasticity robust standard errors in parentheses. Control group mean values of each outcome variable are shown in brackets below each estimate. ``N'' refers to the sample size used in each regression. ``Registered'' refers to a binary indicator for being registered to vote in California in 2022. ``Turnout'' refers to the share of all post-treatment federal general elections between 2010 and 2020 in which a student voted after the academic year in which they filed a FAFSA, which includes up to two elections in 2018 and 2020 for students in the 2017 to 2019 cohorts. ``Ever Voted'' refers to a binary indicator variable for ever having cast a ballot in a federal or state general election. ``(D $\cup$ I ) x Vote'' refers to the interaction between voter turnout in Column 3 and an indicator for whether the student was a registered Democrat or independent, following \cite{firoozi_effect_2023}.  ``R x Vote'' refers to the interaction between voter turnout in Column 3 and an indicator for being registered with the Republican Party.}
    \end{table}


    \clearpage







    \begin{table}[h!]\centering
    \def\sym#1{\ifmmode^{#1}\else\(^{#1}\)\fi}
    \caption{Estimated Impacts of 2017-2019 Cal Grant Receipt on Student Housing Choice\label{tab:housing1}}
    \begin{tabular}{l*{6}{c}}
        \hline\hline
        Outcome&\multicolumn{1}{c}{(1)}&\multicolumn{1}{c}{(2)}&\multicolumn{1}{c}{(3)}&\multicolumn{1}{c}{(4)}&\multicolumn{1}{c}{(5)}&\multicolumn{1}{c}{(6)}\\
        \hline
        &&&&&&\\

        \input{\basepath/Research/college_politics/Tables/oncampus.tex} [0.5em]

        \input{\basepath/Research/college_politics/Tables/offcampus.tex} [0.5em]

        \input{\basepath/Research/college_politics/Tables/withparents.tex} [0.5em]

        \input{\basepath/Research/college_politics/Tables/house_mis.tex}

        &&&&&&\\

        \hline
        Bandwidth & \$10,000       &   \$10,000        &   \$10,000        &   \$10,000       &   \$5,000        &   \$5,000        \\
        Kernel&   Uniform      &   Uniform        &   Uniform        &   Uniform       &   Uniform       &   Uniform       \\
        Polynomial&   1       &   1        &   2        &   2       &   1       &   1       \\
        Controls&   No       &   Yes        &   No        &   Yes        &   No        &   Yes        \\
        Sample Size      &   258,329     &    258,329      &      258,329      &   258,329    &   123,774       &   123,774        \\
        \hline
    \end{tabular}
    \caption*{\footnotesize Note: \sym{+} \(p<0.1\), \sym{*} \(p<0.05\), \sym{**} \(p<0.01\). Heteroskedasticity robust standard errors in parentheses. ``On Campus'' refers to students who received a Cal Grant and lived on campus or who did not receive a Cal Grant and stated an intent to live on campus.``Off Campus'' refers to students who received a Cal Grant and lived off campus at a residence without family or guardians or who did not receive a Cal Grant and stated an intent to live off campus at a residence without family or guardians. ``With Guardians'' refers to students who received a Cal Grant and lived off campus at a residence with family or guardians or who did not receive a Cal Grant and stated an intent to live off campus at a residence with family or guardians.  ``No Housing Info'' refers to students who did not receive a Cal Grant and did not express a preference for their housing intent.}
\end{table}

\clearpage











\subsection*{B. Figures}\label{sec:fig}
    \addcontentsline{toc}{section}{B. Figures}







\begin{figure}[h!]
    \centering{\includegraphics[width=\textwidth]
        {\basepath/Research/college_politics/Figures/rd_outcomes1.png}}\caption{\label{fig:rd_outcomes}Regression Discontinuity Plots of Main Outcomes at Cal Grant Income Ceiling}
    \caption*{\footnotesize Note: Centered income values are normalized to the income ceiling for Cal Grant A for a given individual. Outcomes correspond directly to those in Tablse \ref{tab:mainresults} and \ref{tab:results1}. The sample includes Cal Grant students in the 2017 through 2019 entering student cohorts.}
\end{figure}


\clearpage

\begin{figure}[h!]
    \centering{\includegraphics[width=\textwidth]
        {\basepath/Research/college_politics/Figures/pell_outcomes1.png}}\caption{\label{fig:pell_outcomes}Regression Discontinuity Plots of Main Outcomes at Pell Grant Generosity Notch}
    \caption*{\footnotesize Note: Centered income values are normalized to the Pell Grant's automatic zero expected family contribution threshold. Outcomes correspond directly to those in Tables \ref{tab:mainresults} and \ref{tab:pell_results}. The sample includes Pell Grant students in the 2010 to 2019 entering student cohorts.}
\end{figure}



\clearpage




\begin{figure}[h!]
    \centering{\includegraphics[width=\textwidth]
        {\basepath/Research/college_politics/Figures/rd_housing1.png}}\caption{\label{fig:rd_housing1}Regression Discontinuity Plots of Housing Choice and Intent}
    \caption*{\footnotesize Note: Centered income values are normalized to the income ceiling for Cal Grant A for a given individual. Outcomes correspond directly to those in Table \ref{tab:housing1}. The sample includes Cal Grant students in the 2017 through 2019 entering student cohorts.}
\end{figure}


\clearpage









    \clearpage

    \appendix
    \textbf{\huge{Online Appendices}}

    \subsection*{C. Summary and Descriptive Statistics}\label{sec:appendixc}
    \addcontentsline{toc}{section}{C. Summary and Descriptive Statistics}

    \subsection{Descriptive Statistics}\label{sec:descriptive}
    Table \ref{tab:descriptives} shows summary statistics for the full sample of all FAFSA filers from 2010-2011 to 2020-2021 in Column 1, likely Cal Grant eligible FAFSA filers from 2010-2011 to 2020-2021 in Column 2, likely Cal Grant eligible FAFSA filers from 2017-2018 to 2019-2020 in Column 3, near-threshold and likely eligible FAFSA filers from 2017-2018 to 2019-2020 in Column 4, and near-threshold and likely eligible FAFSA filers for the 2010-2011 to 2019-2020 cohorts in Column 5.

    Beginning with Column 1, we note that the sample of FAFSA filers is likely to be majority Hispanic and Asian, has family income and asset levels of 53,000 dollars and 39,000 dollars, and overwhelmingly favors the Democratic Party relative to the Republican Party.\footnote{Conditional on observing race, a majority of students are Hispanic or Asian American. We expect that consistent with general patterns in the American electorate, non-registered voters are more likely to be Hispanic and Asian American than registered voters. This suggests that our racial and ethnic composition descriptive statistics conditional on registration understate the share of FAFSA filers who are racial or ethnic minorities.} This is similar to the subsample of likely Cal Grant eligible FAFSA filers in Column 2, except for higher political participation and much lower family asset levels due to the Cal Grant's asset eligibility ceiling. In Column 3, we show that the characteristics of likely Cal Grant eligible FAFSA filers in the analysis sample does not change much in the time period after prior-prior year income assessment took effect, suggesting that restricting the sample to those less able to manipulate their income does not change the observable demographics of the sample.

    Finally, in Columns 4 and 5, we note that our main and expanded samples of near threshold FAFSA filers who were likely to be Cal Grant eligible are more likely to be white, have higher incomes and lower family assets, and are more likely to be politically active at baseline than their peers in the full FAFSA filer sample. The intuition behind these differences is that we are focusing on students local to the maximum allowable income to remain eligible for financial aid, while excluding students who are ineligible based on other characteristics. We note that our main sample is a policy-relevant group because debates over the expansion of financial aid often center on extensions of eligibility to higher quantiles of the family income distribution.

    \clearpage

        \begin{table}[h!]\centering
        \def\sym#1{\ifmmode^{#1}\else\(^{#1}\)\fi}
        \caption{Descriptive Statistics for In-Sample FAFSA Filers\label{tab:descriptives}}
        \begin{tabular}{l*{1}{c c c c c}}
            \hline\hline
            Sample  &  Full  & Eligible  & Analysis & Main & Expanded \\
            \hline
            &&&&&\\
            Sample Size & 16,393,526 & 8,563,732 & 3,318,159 & 258,329 & 738,046\\ [0.5em]
            \multicolumn{6}{l}{\textit{A. Race and Ethnicity (Missing if Not Registered to Vote)}} \\ [0.5em]
            Share Latino & 0.225 & 0.280 & 0.307 & 0.277 & 0.243\\
            Share White  & 0.161  & 0.166  & 0.174  & 0.240  & 0.224 \\
            Share Asian & 0.068 & 0.067 & 0.065 & 0.060 & 0.060\\
            Share Black  & 0.019  & 0.020  & 0.020  & 0.020  & 0.019 \\
            Share Other  & 0.028  & 0.029  & 0.031  & 0.031  & 0.029 \\[0.5em]
            \multicolumn{6}{l}{\textit{B. FAFSA Household Characteristics}} \\ [0.5em]
            Family Income & 53,494 & 47,607 & 51,493 & 72,787 & 67,237\\
            Family Assets  & 39,109  & 1,533  & 1,772  & 2,833  & 2,451 \\
            Family Size & 3.184 & 3.442 & 3.444 & 2.839 & 2.904\\
            FAFSA Year  & 2015.68  & 2015.31  & 2018.01  & 2018.00  & 2015.05 \\
            Share Married & 0.085 & 0.086 & 0.083 & 0.064 & 0.067\\[0.5em]
            \multicolumn{6}{l}{\textit{C. Political Characteristics}} \\ [0.5em]
            Democratic  Party & 0.294  & 0.327  & 0.348  & 0.355  & 0.326 \\
            Republican  Party & 0.069 & 0.077 & 0.081 & 0.102 & 0.093\\
            Voted in 2020  & 0.413  & 0.460  & 0.494  & 0.555  & 0.502 \\
            &&&&&\\
            \hline
            Eligiblity Limits  &  No  & Yes  & Yes& Yes& Yes \\
            Earliest Year  &  2010  & 2010  & 2017& 2017& 2010 \\
            Latest Year  &  2020  & 2020  & 2019& 2019& 2019 \\
            Bandwidth Limit  &  None  & None  & None& \$10,000& \$10,000 \\
            \hline
        \end{tabular}
        \caption*{\footnotesize Note: Each column shows the characteristics of a different sample of our FAFSA filer data. Labels for each of these samples are displayed in the first row of the table. Race and ethnicity are only available for registered voters and are missing for all people who are not registered to vote, because race and ethnicity data are only recorded through the voter file. Registered voters may choose to register as a Democrat, a Republican, a member of another political party, or as no party preference.}
    \end{table}


    \clearpage

        \begin{figure}[h!]
        \centering{\includegraphics[width=\textwidth]
            {\basepath/Research/college_politics/Figures/rd_firststage.png}}\caption{\label{fig:rd_firststage}Regression Discontinuity Plot of Cal Grant Receipt}
        \caption*{\footnotesize Note: Centered income values are normalized to the income ceiling for Cal Grant A for a given individual. ``Received Any Cal Grant Award'' refers to an indicator for having received any Cal Grant award in the academic year following FAFSA filing. Observations includes all main sample 2017-2019 Cal Grant eligible students within a 10,000 dollar bandwidth of a Cal Grant income threshold.}
    \end{figure}


    \clearpage

    \subsection{Calculating Macro-level Outcomes}\label{sec:macro}


    In Step 1 of Table \ref{tab:envelope}, we explicitly state the assumptions for each column. The lower bound assumes that the smallest point estimate for voter turnout in the 2020 general election is true, despite the fact that this point estimate relies on data that is  susceptible to potential bias from manipulation of the running variable, students exiting the state, and legal name changes after marriage. The lower bound also relies on the lowest ratio of new center-left to new center-right voters across our specifications and assumes that 80 percent of center-left voters support Biden and 0 percent of registered Republicans support Biden.\footnote{Based on actual election returns and surveys of California college students that matched party registration records to policy preferences, this is likely to be an underestimate. At both private and public 4-year college campus precincts in California and other states, between 80 and 95 percent of voters cast a ballot for Democratic candidates. Likewise, a survey of Californians who applied to college between 2007 and 2011 that was fielded in 2022 found that approximately 75 percent of registered third party and non-partisan voters, 94 percent of registered Democrats, and 23 percent of registered Republicans reported that they favored the Democratic Party over the Republican Party \citep{firoozi_effect_2023}. Evidence from political contribution data suggest even higher rates of support for Democratic candidates.} Lastly, the lower bound estimate assumes that the number of Cal Grant recipients in 2011-2012 was the same as 2010-2011, despite the recipient totals monotonically increasing each year.

    Our best estimate uses the results of our preferred specification for 2020 turnout and the partisan composition of new voters, making the same assumptions on vote by party as the lower bound column. It also assumes that the number of Cal Grant recipients in 2011-2012 was the average of the 2010-2011 and 2012-2013 cohorts. Finally, the upper bound uses the largest point estimate on the treatment effect for 2020 turnout, assumes the full net increase accrues to Democrats\footnote{This is mathematically feasible because some specifications find a net decrease in GOP turnout, with more than the whole of the net increase accruing to registered independents and Democrats.}, and that the 2011-2012 Cal Grant awards were equal to the 2012-2013 total.

    In Step 2, we multiply out the assumptions from Step 1 and compare them to the actual results of the 2020 presidential election in California across each column. In Step 3, we deduct the estimated impacts of 2010-2019 Cal Grants from the actual results and show the projected vote totals if the grants had not been awarded. In Step 4, we show the net impact of the 2010-2019 Cal Grants on the Biden-Trump margin of victory in the state of California and the aggregate voter turnout rate relative to the total citizen voting-eligible population (CVEP).

    \clearpage


    \begin{table}[h!]\centering
        \centering
        \caption{Estimated Impact of 2010-2019 Cal Grants on the 2020 Election in California}
        \label{tab:envelope}
        \begin{tabular}{l*{3}{r}}
            \hline\hline
            & Lower & Best & Upper  \\
            & Bound &  Estimate &  Bound \\
            \hline
            \multicolumn{4}{l}{\textit{Step 1: Assumptions}} \\ [0.5em]
            1.a) Average Treatment Effect& +3.86 pp & +9.85 pp & +11.93 pp \\[0.5em]
            1.b) Partisan Assignment& 70\% D & 80\% D & 100\% D \\
            & 25\% R & 15\% R& 0\% R \\
            & 5\% O & 5\% O & 0\% O \\[0.5em]
            1.c) Total Grants Awarded& 2,612,744& 2,626,306& 2,639,867  \\[0.5em]
            \hline
            \multicolumn{4}{l}{\textit{Step 2: Actual 2020 Outcomes and Calculated Effects}} \\ [0.5em]
            2.a) Actual 2020 Biden Votes & & 11,110,250 & \\
            & &  (63.48 pp)& \\
            Impact of Cal Grants & +70,596 & +206,953 & +314,936 \\[0.5em]
            2.b) Actual 2020 Trump Votes&  & 6,006,429 &  \\
            & &  (34.32 pp)& \\
            Impact of Cal Grants & +25,213 & +38,803 & +0 \\[0.5em]
            2.c) Actual 2020 Other Votes&  & 384,202 &  \\
            & &  (2.20 pp)& \\
            Impact of Cal Grants & +5,043 & +12,935 & +0 \\[0.5em]
            2.d) Actual 2020 Total Votes&  & 17,500,881 &  \\
            & &  (100.00 pp)& \\
            Impact of Cal Grants & +100,852 & +258,691 & +314,936 \\[0.5em]
            \hline
            \multicolumn{4}{l}{\textit{Step 3: Estimated Outcomes without 2010-2019 Cal Grants}} \\ [0.5em]
            3.a) Projected 2020 Biden Votes & 11,039,654& 10,903,297 &10,795,314 \\
            & (63.45 pp)&  (63.24 pp)&(62.81 pp) \\[0.5em]
            3.b) Projected 2020 Trump Votes & 5,981,216& 5,967,626&6,006,429\\
            & (34.37 pp)&  (34.61 pp)&  (34.95 pp)\\[0.5em]
            3.c) Projected 2020 Other Votes & 379,159& 371,267&384,202 \\
            & (2.18 pp)&  (2.15 pp)&(2.24 pp) \\[0.5em]
            3.d) Projected 2020 Total Votes & 17,400,029& 17,242,190 &17,185,945  \\
            & (100.00 pp)&  (100.00 pp)&(100.00 pp) \\[0.5em]
            \hline
            \multicolumn{4}{l}{\textit{Step 4: Estimated Impacts of 2010-2019 Cal Grants}} \\ [0.5em]
            4.a) Impact on Biden-Trump Margin& +45,383& +168,150&+314,936\\
            &+0.08 pp &  +0.53 pp& +1.30 pp \\[0.5em]
            4.b)  Impact on Voter Turnout Rate & +0.40 pp& +1.03 pp &+1.26 pp\\
            \hline
        \end{tabular}
        \caption*{\footnotesize Note: The citizen voting eligible population (CVEP) was 25,090,517 in California for the 2020 presidential general election.}
    \end{table}

    \clearpage

    \begin{figure}[h!]
        \centering{\includegraphics[width=\textwidth]
            {\basepath/Research/college_politics/Figures/map_recipients_by_county.png}}\caption{\label{fig:map}Map of 2010-2020 Cal Grant Recipients in 2022}
        \caption*{\footnotesize Note: The map above shows the number of Cal Grant recipients who were registered to vote per capita in each of California's counties. Locations are assessed based on the 2022 voter registration address of Cal Grant recipients. See Figure \ref{fig:campus_map} for a map of public university campuses.}
    \end{figure}


    \clearpage

    \begin{figure}[h!]
        \centering{\includegraphics[width=\textwidth]
            {\basepath/Research/college_politics/Figures/Campus_Map.png}}\caption{\label{fig:campus_map}Map of Public 4-Year Colleges in California}
        \caption*{\footnotesize Note: California's public universities are plotted with symbols and county lines in the map above. Blue triangles represent the University of California's (UC) undergraduate campuses, which are highly selective research universities. Of the 9 UC campuses, 8 are categorized as R1 research universities. Red circles represent the California State University's (CSU) campuses, which are selective local comprehensive universities that primarily focus on teaching. See Figure \ref{fig:map} for a heat map of the Cal Grant's political impact by county.}
    \end{figure}


    \clearpage




\subsection*{D. RD Validation}\label{sec:appendixa}
    \addcontentsline{toc}{section}{D. RD Validation}



        \begin{figure}[h!]
        \centering{\includegraphics[width=\textwidth]
            {\basepath/Research/college_politics/Figures/McCrary_main.png}}\caption{\label{fig:mccrary_2017}McCrary Density Test}
        \caption*{\footnotesize Note: Centered income values are normalized to the income ceiling for Cal Grant A for a given individual. Observations include all main sample 2017-2019 Cal Grant eligible students within a 10,000 dollar bandwidth of a Cal Grant income threshold.}
    \end{figure}

    \clearpage


\begin{figure}[h!]
    \centering{\includegraphics[width=\textwidth]
        {\basepath/Research/college_politics/Figures/rd_covariates1.png}}\caption{\label{fig:rd_covariates1}Regression Discontinuity Plot for Covariates}
    \caption*{\footnotesize Note: Centered income values are normalized to the income ceiling for Cal Grant A for a given individual. Observations include all main sample 2017-2019 Cal Grant eligible students within a 10,000 dollar bandwidth of a Cal Grant income threshold.}
\end{figure}


\clearpage

\begin{figure}[h!]
    \centering{\includegraphics[width=\textwidth]
        {\basepath/Research/college_politics/Figures/rd_covariates2.png}}\caption{\label{fig:rd_covariates2}Regression Discontinuity Plot for Covariates}
    \caption*{\footnotesize Note: Centered income values are normalized to the income ceiling for Cal Grant A for a given individual. Observations include all main sample 2017-2019 Cal Grant eligible students within a 10,000 dollar bandwidth of a Cal Grant income threshold.}
\end{figure}


\clearpage

\begin{figure}[h!]
    \centering{\includegraphics[width=\textwidth]
        {\basepath/Research/college_politics/Figures/rd_covariates3.png}}\caption{\label{fig:rd_covariates3}Regression Discontinuity Plot for Covariates}
    \caption*{\footnotesize Note: Centered income values are normalized to the income ceiling for Cal Grant A for a given individual. Observations include all main sample 2017-2019 Cal Grant eligible students within a 10,000 dollar bandwidth of a Cal Grant income threshold.}
\end{figure}


\clearpage


\begin{figure}[h!]
    \centering{\includegraphics[width=\textwidth]
        {\basepath/Research/college_politics/Figures/omnibustest.png}}\caption{\label{fig:rd_omnibus}Omnibus Covariate Balance Test}
    \caption*{\footnotesize Note: Centered income values are normalized to the income ceiling for Cal Grant A for a given individual. Observations include all main sample 2017-2019 Cal Grant eligible students within a 10,000 dollar bandwidth of a Cal Grant income threshold.}
\end{figure}


\clearpage


\begin{figure}[h!]
    \centering{\includegraphics[width=\textwidth]
        {\basepath/Research/college_politics/Figures/bw_covariates1.png}}\caption{\label{fig:bw_covariates1}Covariate Discontinuities by Bandwidth}
    \caption*{\footnotesize Note: Each panel shows the reduced form disconitinuity in a coviarate across a range of potential bandwidths using a local linear specification and a uniform kernel without covariates. The graphs reflect the point estimate, 95 percent confidence interval, and 90 percent confidence interval. Observations include all main sample 2017-2019 Cal Grant eligible students within a given bandwidth of a Cal Grant income threshold.}
\end{figure}


\clearpage

\begin{figure}[h!]
    \centering{\includegraphics[width=\textwidth]
        {\basepath/Research/college_politics/Figures/bw_covariates2.png}}\caption{\label{fig:bw_covariates2}Covariate Discontinuities by Bandwidth}
    \caption*{\footnotesize Note: Each panel shows the reduced form disconitinuity in a coviarate across a range of potential bandwidths using a local linear specification and a uniform kernel without covariates. The graphs reflect the point estimate, 95 percent confidence interval, and 90 percent confidence interval. Observations include all main sample 2017-2019 Cal Grant eligible students within a given bandwidth of a Cal Grant income threshold.}
\end{figure}


\clearpage

\begin{figure}[h!]
    \centering{\includegraphics[width=\textwidth]
        {\basepath/Research/college_politics/Figures/bw_covariates3.png}}\caption{\label{fig:bw_covariates3}Covariate Discontinuities by Bandwidth}
    \caption*{\footnotesize Note: Each panel shows the reduced form disconitinuity in a coviarate across a range of potential bandwidths using a local linear specification and a uniform kernel without covariates. The graphs reflect the point estimate, 95 percent confidence interval, and 90 percent confidence interval. Observations include all main sample 2017-2019 Cal Grant eligible students within a given bandwidth of a Cal Grant income threshold.}
\end{figure}


\clearpage





\subsection{Interstate Migration} \label{sec:appendixasub}
We note that there is a match rate of 50 to 60 percent, depending on the particular sample used, between FAFSA files from CSAC and L2's California voter file. Because the merged dataset is de-identified, it is not possible to know the precise registration rate for unique individuals within the dataset. Non-matches between these two datasets are attributable to students moving to another state/country, lack of interest in voting or registering to vote while living in California, or disenfranchisement (for example, due to incarceration or death). Based on the descriptive statistics we observe for the tables generated in this appendix, we estimate that approximately three quarters of the non-matches between datasets are attributable to Californian residents who simply lack interest in registration or voting, with the remaining quarter attributable to migration out of California.


We use four approaches to address the risk that out-of-state migration may bias our estimates.

First, we note that National Student Clearinghouse data and IRS tax records have been used to track movers at the Cal Grant income ceiling, confirming the validity of our approach. Receiving a Cal Grant has insignificant effects on the share of students enrolling at in-state colleges in the short-run and null effects on IRS tax filings in California until 11 years after FAFSA filing. \cite{bettinger_long-run_2019} formally estimate an insignificant increase in \textit{out-of-state} enrollment for each Cal Grant received. Because moving out of California makes a student's political participation records unobservable in our dataset, a literal interpretation of the authors' estimates would mean that our short run estimates are biased \textit{toward zero} by roughly 2 percentage points (assuming a 50 percent turnout rate). The paper also uses IRS data to show null effects of the Cal Grant on California residence ($<$2 pp) followed by a significant increase in retention within the state of California of approximately 2 percentage points after 11 years have elapsed since FAFSA filing. Because the falling share of students who reside in California over time should partially offset the selective retention induced by the Cal Grant after 11 years, we expect relocation to be a minor source of upward bias of approximately 1 percentage point in estimates after 11 years have elapsed (again assuming a 50 percent turnout rate).

Second, we use the subsample of our dataset in which we can observe out-of-state migration to confirm that our results match those of earlier work on the Cal Grant. Specifically, we use the subset of students that overlap between our main dataset and the main dataset used by \cite{firoozi_effect_2023}, which would be UC applicants filing a financial aid application in anticipation of the 2010 and 2011 academic years. We show that receiving a Cal Grant has no measured impact on out-of-state voter registration in Table \ref{tab:extension} for this sample and that, under most specifications, the point estimate is positive. Each specification reflects a result that is both small in absolute magnitude and in economic terms. The observation that most specifications generate a positive point estimate suggests that out-of-state migration may very slightly bias our estimates toward zero, which is consistent with previous work on the Cal Grant's impact on migration.

Third, we use a unique feature of our L2 dataset to illustrate that attrition has not made the sample unbalanced. In the L2 dataset, people who move out of the state of California will have no California voter record in either pre-treatment or post-treatment elections. This is because L2, which observes all state voter files, prunes historical voting records of interstate migrants from the dataset of the origin state. Therefore, if ambitious students with a high propensity to vote who are barely too rich to receive the Cal Grant attrit to other states, we should observe a discontinuity in pre-treatment voter turnout mirroring the results for post-treatment voter turnout. As we show in Table \ref{tab:placebo_prevturnout}, there are null effects of receiving a Cal Grant on pre-treatment turnout.

Fourth, we choose to be conservative in our approach to identifying causal effects. We begin our results section with a focus on the 2017-2018 to 2019-2020 sample, who filed a FAFSA less than 5 years prior to our voter file snapshot, minimizing the possibility of selection bias due to out-of-state migration. We also externally validate our results with a discontinuity in the generosity of the federal Pell Grant, showing that our findings generalize to policies that subsidize tuition at colleges outside of California. For our results to be explained by out-of-state migration, the Cal Grant's effect on short-run out-of-state migration would need to have changed from a negative value in the early 2000s to a present day effect size of roughly 20 percentage points per grant awarded and the Pell Grant would need to induce selection into California at a similar per dollar rate.

\clearpage

    \begin{table}[h!]\centering
    \def\sym#1{\ifmmode^{#1}\else\(^{#1}\)\fi}
    \caption{Estimated Impacts of 2010-2011 Cal Grants on Out-of-State Registration\label{tab:extension}}
    \begin{tabular}{l*{6}{c}}
        \hline\hline
        Outcome&\multicolumn{1}{c}{(1)}&\multicolumn{1}{c}{(2)}&\multicolumn{1}{c}{(3)}&\multicolumn{1}{c}{(4)}&\multicolumn{1}{c}{(5)}&\multicolumn{1}{c}{(6)}\\
        \hline
        &&&&&&\\

        \input{\basepath/Research/college_politics/Tables/extension_oosreg.tex}


        &&&&&&\\

        \hline
        Bandwidth & \$10,000       &   \$10,000        &   \$10,000        &   \$10,000       &   \$5,000        &   \$5,000        \\
        Kernel&   Uniform      &   Uniform        &   Uniform        &   Uniform       &   Uniform       &   Uniform       \\
        Polynomial&   1       &   1        &   2        &   2       &   1       &   1       \\
        Controls&   No       &   Yes        &   No        &   Yes        &   No        &   Yes        \\
        Sample Size      &      6,589     &      6,589      &   6,589    &   6,589       &   3,318     &    3,318      \\
        \hline
    \end{tabular}
    \caption*{\footnotesize Note: \sym{+} \(p<0.1\), \sym{*} \(p<0.05\), \sym{**} \(p<0.01\). Heteroskedasticity robust standard errors in parentheses.    ``Non-CA Registered'' refers to the proportion of students who were observed as registered to vote in a state other than California. Data is from a sample of UC applicants between the 2010-2011 and 2011-2012 academic years who filed a FAFSA \citep{firoozi_effect_2023}. }
\end{table}

\clearpage

    \begin{table}[h!]\centering
    \def\sym#1{\ifmmode^{#1}\else\(^{#1}\)\fi}
    \caption{Placebo Test of 2017-2019 Cal Grant Receipt on Pre-Treatment Voter Turnout\label{tab:placebo_prevturnout}}
    \begin{tabular}{l*{6}{c}}
        \hline\hline
        Outcome&\multicolumn{1}{c}{(1)}&\multicolumn{1}{c}{(2)}&\multicolumn{1}{c}{(3)}&\multicolumn{1}{c}{(4)}&\multicolumn{1}{c}{(5)}&\multicolumn{1}{c}{(6)}\\
        \hline
        &&&&&&\\

        \input{\basepath/Research/college_politics/Tables/placebo_prev_turnout.tex} [0.5em]


        &&&&&&\\

        \hline
        Bandwidth & \$10,000       &   \$10,000        &   \$10,000        &   \$10,000       &   \$5,000        &   \$5,000        \\
        Kernel&   Uniform      &   Uniform        &   Uniform        &   Uniform       &   Uniform       &   Uniform       \\
        Polynomial&   1       &   1        &   2        &   2       &   1       &   1       \\
        Controls&   No       &   Yes        &   No        &   Yes        &   No        &   Yes        \\
        Sample Size      &   258,329     &    258,329      &      258,329      &   258,329    &   123,774       &   123,774        \\
        \hline
    \end{tabular}
    \caption*{\footnotesize Note: \sym{+} \(p<0.1\), \sym{*} \(p<0.05\), \sym{**} \(p<0.01\). Heteroskedasticity robust standard errors in parentheses. ``Voted Pre-FAFSA'' refers to having voted in the general election immediately prior to filing a FAFSA and potentially having received a Cal Grant.}
\end{table}

\clearpage

\subsection*{E. Robustness Checks}\label{sec:appendixb}
    \addcontentsline{toc}{section}{E. Robustness Checks}




\subsection{Effects of 2017-2019 Cal Grants}\label{sec:cal1719}


We test our main outcomes formally in Table \ref{tab:results1}. Each row of the table reflects the IV estimate of the effect of receiving a Cal Grant on the aforementioned outcomes and each column represents a different specification. Column 1 begins with our preferred specification, which uses local linear estimation with a uniform kernel at a 10,000 dollar bandwidth without covariates. Column 2 adds a set of pre-FAFSA covariate controls. Columns 3 and 4 increase the order of the polynomial control for the running variable to a quadratic functional form, with and without covariates. Columns 5 and 6 use local linear estimation at a much narrower 5,000 dollar bandwidth, again varying the inclusion of pre-FAFSA covariates.

Beginning with voter registration in Panel A, we find that the Cal Grant generates noisy, if any, increases in aggregate voter registration. Our preferred specification in Column 1 yields a point estimate of 6.34 percentage points per grant awarded. However, the estimates in Columns 2 through 6 are smaller and not significant, which suggests that the estimated impact of Cal Grant receipt on voter registration are sensitive to model specification and could be null or positive and small.

In Panel B, we present evidence that the rate at which students actually cast a ballot rises sharply as a result of California's tuition-free college program. For each grant awarded by the state, a student's odds of casting a ballot in the 2020 general election rose by 9.85 percentage points, which is significant at a 99 percent confidence interval.  For all general elections held after the academic year in which students filed their FAFSA, we estimate a similar increase of 8.55 percentage points per grant awarded. These findings are consistent across each successive column and remain significant at a 90 percent confidence interval, suggesting that the effect of Cal Grants is robust to different specifications and definitions of turnout.

Panel C concludes by interacting voter turnout with measures of student partisanship. Row 4 begins with an indicator for Democratic or independent registration status interacted with a student's general election turnout rate. We find that, for each grant awarded, the rate at which students turnout for elections as a Democrat or independent increases by 8.91 percentage points, which is significant at a 99 percent confidence interval. This result is much larger than the effect on Republican turnout in Row 5, which is estimated to be -0.36 percentage points per grant awarded. Taken together, these results suggest that the Cal Grant program substantially raises political participation, exclusively among left-leaning voters.

\clearpage


\begin{table}[h!]\centering
    \def\sym#1{\ifmmode^{#1}\else\(^{#1}\)\fi}
    \caption{Estimated Impacts of 2017-2019 Cal Grant Receipt on Political Participation\label{tab:results1}}
    \begin{tabular}{l*{6}{c}}
        \hline\hline
        Outcome&\multicolumn{1}{c}{(1)}&\multicolumn{1}{c}{(2)}&\multicolumn{1}{c}{(3)}&\multicolumn{1}{c}{(4)}&\multicolumn{1}{c}{(5)}&\multicolumn{1}{c}{(6)}\\
        \hline
        &&&&&&\\
        \multicolumn{7}{l}{\textit{A. Voter Registration}} \\ [0.5em]
        \input{\basepath/Research/college_politics/Tables/registered.tex} [0.5em]

        \multicolumn{7}{l}{\textit{B. Voter Turnout}} \\ [0.5em]

        \input{\basepath/Research/college_politics/Tables/g2020.tex} [0.5em]

        \input{\basepath/Research/college_politics/Tables/turnoutrate.tex} [0.5em]

        \multicolumn{7}{l}{\textit{C. Voter Turnout by Partisanship}} \\ [0.5em]

        \input{\basepath/Research/college_politics/Tables/centerleftvote.tex} [0.5em]

        \input{\basepath/Research/college_politics/Tables/centerrightvote.tex}

        &&&&&&\\

        \hline
        Bandwidth & \$10,000       &   \$10,000        &   \$10,000        &   \$10,000       &   \$5,000        &   \$5,000        \\
        Kernel&   Uniform      &   Uniform        &   Uniform        &   Uniform       &   Uniform       &   Uniform       \\
        Polynomial&   1       &   1        &   2        &   2       &   1       &   1       \\
        Controls&   No       &   Yes        &   No        &   Yes        &   No        &   Yes        \\
        Sample Size      &   258,329     &    258,329      &      258,329      &   258,329    &   123,774       &   123,774        \\
        \hline
    \end{tabular}
    \caption*{\footnotesize Note: \sym{+} \(p<0.1\), \sym{*} \(p<0.05\), \sym{**} \(p<0.01\). Heteroskedasticity robust standard errors in parentheses. ``Voter Turnout'' refers to the share of all federal general elections between 2010 and 2020 in which a student voted after the academic year in which they filed a FAFSA, which includes up to two elections in 2018 and 2020 for these student cohorts. ``Center-left Turnout'' refers to the interaction between voter turnout and an indicator for whether the student was a registered Democrat or independent, following \cite{firoozi_effect_2023}.  ``Center-right Turnout'' refers to the interaction between voter turnout and an indicator for being registered with the Republican Party.}
\end{table}

\clearpage




    \begin{table}[h!]\centering
    \def\sym#1{\ifmmode^{#1}\else\(^{#1}\)\fi}
    \caption{Impacts of 2017-2019 Cal Grants on Political Participation by Gender\label{tab:gender_results1}}
    \begin{tabular}{l*{6}{c}}
        \hline\hline
        Outcome&\multicolumn{1}{c}{(1)}&\multicolumn{1}{c}{(2)}&\multicolumn{1}{c}{(3)}&\multicolumn{1}{c}{(4)}&\multicolumn{1}{c}{(5)}&\multicolumn{1}{c}{(6)}\\
        \hline
        &&&&&&\\
        \multicolumn{7}{l}{\textit{A. Non-Female}} \\ [0.5em]
        \input{\basepath/Research/college_politics/Tables/men_registered.tex} [0.5em]

        \input{\basepath/Research/college_politics/Tables/men_g2020.tex} [0.5em]

        \input{\basepath/Research/college_politics/Tables/men_turnoutrate.tex} [0.5em]

            Sample Size      &    112,527    &     112,527   &      112,527    &     112,527   &   53,478    &   53,478        \\[0.5em]

        \multicolumn{7}{l}{\textit{B. Female}} \\ [0.5em]

        \input{\basepath/Research/college_politics/Tables/women_registered.tex} [0.5em]

        \input{\basepath/Research/college_politics/Tables/women_g2020.tex} [0.5em]

        \input{\basepath/Research/college_politics/Tables/women_turnoutrate.tex} [0.5em]

            Sample Size      &  145,802    &   145,802     &     145,802     &   145,802    &   70,296    &   70,296        \\

        &&&&&&\\

        \hline
        Bandwidth & \$10,000       &   \$10,000        &   \$10,000        &   \$10,000       &   \$5,000        &   \$5,000        \\
        Kernel&   Uniform      &   Uniform        &   Uniform        &   Uniform       &   Uniform       &   Uniform       \\
        Polynomial&   1       &   1        &   2        &   2       &   1       &   1       \\
        Controls&   No       &   Yes        &   No        &   Yes        &   No        &   Yes        \\

        \hline
    \end{tabular}
    \caption*{\footnotesize Note: \sym{+} \(p<0.1\), \sym{*} \(p<0.05\), \sym{**} \(p<0.01\). Heteroskedasticity robust standard errors in parentheses. ``Voter Turnout'' refers to the share of all federal general elections between 2010 and 2020 in which a student voted after the academic year in which they filed a FAFSA.}
\end{table}

\clearpage

\begin{figure}[h!]
    \centering{\includegraphics[width=\textwidth]
        {\basepath/Research/college_politics/Figures/bw_results1.png}}\caption{\label{fig:bw_results}Bandwidth Robustness Tests for Main Outcomes}
    \caption*{\footnotesize Note: Each panel shows the reduced form discontinuity in an outcome variable across a range of potential bandwidths using a local linear specification and a uniform kernel without covariates. The graphs reflect the point estimate, 95 percent confidence interval, and 90 percent confidence interval of the reduced form discontinuity in the outcome of interest. Outcomes correspond directly to those in Table \ref{tab:results1}.}
\end{figure}


\clearpage

\begin{figure}[h!]
    \centering{\includegraphics[width=\textwidth]
        {\basepath/Research/college_politics/Figures/bw_registered.png}}\caption{\label{fig:bw_registered}Voter Registration in 2022 Bandwidth Robustness Tests}
    \caption*{\footnotesize Note: The graphs reflect the point estimate, 95 percent confidence interval, and 90 percent confidence interval of the reduced form discontinuity in the outcome of interest.  Each panel represents a different specification. Observations include all main sample 2017-2019 Cal Grant eligible students within a 14,000 dollar bandwidth of a Cal Grant income threshold.}
\end{figure}


\clearpage

\begin{figure}[h!]
    \centering{\includegraphics[width=\textwidth]
        {\basepath/Research/college_politics/Figures/bw_g2020.png}}\caption{\label{fig:bw_g2020}2020 General Election Turnout Bandwidth Robustness Tests}
    \caption*{\footnotesize Note: The graphs reflect the point estimate, 95 percent confidence interval, and 90 percent confidence interval of the reduced form discontinuity in the outcome of interest.  Each panel represents a different specification. Observations include all main sample 2017-2019 Cal Grant eligible students within a 14,000 dollar bandwidth of a Cal Grant income threshold.}
\end{figure}


\clearpage

\begin{figure}[h!]
    \centering{\includegraphics[width=\textwidth]
        {\basepath/Research/college_politics/Figures/bw_turnoutrate.png}}\caption{\label{fig:bw_turnoutrate}General Election Turnout Bandwidth Robustness Tests}
    \caption*{\footnotesize Note: The graphs reflect the point estimate, 95 percent confidence interval, and 90 percent confidence interval of the reduced form discontinuity in the outcome of interest.  Each panel represents a different specification. Observations include all main sample 2017-2019 Cal Grant eligible students within a 14,000 dollar bandwidth of a Cal Grant income threshold.}
\end{figure}


\clearpage

\begin{figure}[h!]
    \centering{\includegraphics[width=\textwidth]
        {\basepath/Research/college_politics/Figures/bw_centerleftvote.png}}\caption{\label{fig:bw_centerleftvote}Democratic or Independent Turnout Bandwidth Robustness Tests}
    \caption*{\footnotesize Note: The graphs reflect the point estimate, 95 percent confidence interval, and 90 percent confidence interval of the reduced form discontinuity in the outcome of interest.  Each panel represents a different specification. Observations include all main sample 2017-2019 Cal Grant eligible students within a 14,000 dollar bandwidth of a Cal Grant income threshold.}
\end{figure}


\clearpage

\begin{table}[h!]\centering
    \def\sym#1{\ifmmode^{#1}\else\(^{#1}\)\fi}
    \caption{Main Results for 2017-2019 with CCT Bias-Aware Confidence Intervals\label{tab:rdrobust1}}
    \begin{tabular}{{l} {c} {c}}
        \hline\hline
        Outcome &\multicolumn{1}{c}{(1)}&\multicolumn{1}{c}{(2)}\\
        \hline
        \multicolumn{3}{l}{\textit{Voter Registration}} \\
        \input{\basepath/Research/college_politics/Tables/rdrobust_registered.tex}[0.5em]
        \multicolumn{3}{l}{\textit{Voted in 2020}} \\
        \input{\basepath/Research/college_politics/Tables/rdrobust_g2020.tex}[0.5em]
        \multicolumn{3}{l}{\textit{Voter Turnout}} \\
        \input{\basepath/Research/college_politics/Tables/rdrobust_turnoutrate.tex}[0.5em]
        \multicolumn{3}{l}{\textit{Center-Left Turnout}} \\
        \input{\basepath/Research/college_politics/Tables/rdrobust_centerleftvote.tex}[0.5em]
        \multicolumn{3}{l}{\textit{Center-Right Turnout}} \\
        \input{\basepath/Research/college_politics/Tables/rdrobust_centerrightvote.tex}[0.5em]
        \hline
        Polynomial&   1       &   1        \\
        Kernel &   Triangular       &   Triangular        \\
        Bandwidth&   \$10,000 & \$5,000              \\
        \hline
    \end{tabular}
    \caption*{\footnotesize Note: \sym{+} \(p<0.1\), \sym{*} \(p<0.05\), \sym{**} \(p<0.01\). Each row titled ``RD Estimate'' shows the conventional point estimate and standard errors in parentheses for a given outcome variable. These are calculated using a triangular kernel for a local linear specification without covariates. The rows ``Robust 95\% CI'' and ``Robust p-value'' show the bias-corrected confidence interval and the bias-corrected p-value for the same outcome variable \citep{calonico_robust_2014}. These outcomes correspond to those in Table \ref{tab:results1}.}
\end{table}

\clearpage




\begin{figure}[h!]
    \centering{\includegraphics[width=\textwidth]
        {\basepath/Research/college_politics/Figures/falsif1_results1.png}}\caption{\label{fig:bw_falsif1_results1}Placebo Falsification Tests for Main Outcomes}
    \caption*{\footnotesize Note: Each panel shows the cumulative distribution of t-statistics estimated at placebo thresholds for an outcome of interest with a red dashed line indicating the estimated t-statistic at the true policy threshold. We generate a ``placebo threshold'' at each 500 dollar increment along centered family income, and compare the estimated reduced form impact of these synthetic policys relative to the true policy. Placebo thresholds are bounded between -20,000 and +60,000 dollars relative to the true income ceiling, because this avoids false positives from capturing discontinuities taking place at family incomes of zero at the lower bound and this spans up to the 98th percentile of centered income on the upper bound. A 10,000 dollar bandwidth is used to remain consistent with our preferred specification. We exclude discontinuities within a 10,000 dollar bandwidth of the true cutoff to avoid generating false positives by including the actual policy discontinuity in our placebo estimates. Observations include all main sample 2017-2019 Cal Grant eligible students.}
\end{figure}


\clearpage



\subsection{External Validity and the Pell Grant}\label{sec:pell}

We externally validate our findings from the Cal Grant using a notch in the generosity of the United States' largest tuition subsidy, the Pell Grant. The Pell Grant is a federal need-based tuition subsidy awarded to between 6 and 10 million students per year with identical eligibility and selection criteria in all American states.\footnote{The Pell Grant does not have a strict 3.0 GPA cutoff like the Cal Grant. Roughly 1 out of 6 Pell Grant recipients are located in the state of California in any given year.} The Pell Grant is less generous than the Cal Grant, covering at most half the cost of tuition and fees at University of California campuses as opposed to full tuition and fees. Students submit a FAFSA to apply for a Pell Grant and eligibility is determined by a mix of characteristics including family income, household size, family assets, and a college's cost of attendance. The program targets a much poorer subset of students than the Cal Grant, with few families with incomes over 60,000 dollars per year receiving a Pell Grant award as opposed to the upper bound of roughly 120,000 dollars per year for the richest Cal Grant students. The Pell Grant also subsidizes a much wider subset of post-secondary institutions with close to half of recipients attending 2-year colleges or vocational programs that are essentially excluded from tuition subsidies under the Cal Grant. We view the Pell Grant as an effective context to assess the generalizability of our findings because, relative to the Cal Grant, it is less generous, it subsidizes a wider and much less selective set of post-secondary institutions, and its eligibility criteria are both broader than state aid and consistent throughout the United States.

To identify the impact of the Pell Grant, we exploit a notch in the intensive margin of subsidy generosity rather than the extensive margin of subsidy receipt. When a students' family income falls below a certain level, they are automatically assigned an expected family contribution (EFC) level of zero dollars toward their cost of college attendance. This means that the generosity of the federal Pell Grant as measured through the size of the financial aid grant awarded to the student rises sharply below the zero EFC income threshold \citep{denning_propelled:_2019}. Our estimate of the change in political participation at the Pell Grant income threshold will compare students who receive marginally more financial aid dollars to people who receive marginally fewer dollars from the same tuition subsidy.\footnote{It is worth bearing in mind that RD designs estimate the average treatment effects for compliers near some eligibility threshold. The notch at the zero EFC threshold we study takes place for family incomes in roughly the 20,000 to 30,000 dollar per year range whereas our Cal Grant thresholds are usually in the 40,000 to 120,000 dollar per year range for most students.}

We test the external validity of our findings formally with data on 2.5 million FAFSA filers from our eligible sample (see Column 2 of Table \ref{tab:descriptives}) who are local to the income notch used to determine Pell Grant generosity. Following the strategy in  \cite{denning_propelled:_2019}, we center students' adjusted gross family income (AGI) relative to the level that automatically qualifies a student for zero EFC and increases Pell Grant generosity by 700 dollars.\footnote{A small part of this increase in generosity is due to crowding-in of state-level financial aid. While we find that the zero EFC threshold crowds in Cal Grant awards at a 0.3 percentage point rate, we know from Section \ref{sec:result} of our paper that this crowd-in is unable to explain more than a 0.03 percentage point increase in voter turnout at the notch for Pell Grant generosity.} Because we do not directly observe total financial aid awarded, we assume a discontinuity of 1,000 dollars to be conservative.\footnote{The size of this discontinuity was generally decreasing over time, so this is likely to be an overestimate \citep{denning_propelled:_2019}. Assuming a larger first-stage discontinuity deflates the estimated effect per dollar awarded. For perspective, a typical Cal Grant recipient receives around ten times as many dollars over our sample timeframe. Hence, a quick point of comparison to our results for the Cal Grant in Table \ref{tab:results2} entails multiplying estimates in this section by 10, which we discuss in greater detail in Section \ref{sec:policy}.}

We begin by plotting our main outcomes of interest against a student's centered AGI in Figure \ref{fig:pell_outcomes}. Mirroring the results in Figure \ref{fig:rd_outcomes}, the four panels of Figure \ref{fig:pell_outcomes} display the share of students who were registered to vote in 2022, the share of students who voted in the 2020 general election, the total voter turnout rate across all post-treatment general elections betwen 2010 and 2020, and the interaction between the voter turnout rate across post-treatment general elections and an indicator for being registered as a member of the Democratic Party or an independent. There is clear evidence of a discontinuity in each outcome and the results are identical to those of the Cal Grant. Receiving more generous tuition subsidies from America's largest financial aid program, the federal Pell Grant, increases voter registration and turnout, largely among left-leaning students. We take this figure to provide strong evidence that the civic externalities of the Cal Grant and Pell Grant generalize to tuition subsidies that cover different types of post-secondary institutions, target different populations of students, and have different levels and margins of generosity.

In Table \ref{tab:pell_results}, we formally estimate the intent to treat (ITT) effects of this notch in Pell Grant generosity. We replicate our outcomes and specifications from Table \ref{tab:results2}'s estimates of the Cal Grant and find similar results. Our preferred specification suggests that raising Pell Grant generosity by 1,000 dollars at this notch increases voter turnout in the 2020 election by 0.52 percentage points, with around four fifths of the effect occurring among left-of-center voters. These estimates are robust to alternative definitions of political participation and a number of RD implementation choices like changing the bandwidth used for inference, including pre-treatment covariates, and changing the order of a polynomial control for the running variable. Each point estimate is just under one tenth of the corresponding effect size in Table \ref{tab:results2}, meaning that our results across the Pell Grant and Cal Grant imply similar per-dollar estimates of the civic externalities of tuition subsidies. We find the similarity especially notable given the programs target different populations, subsidize different postsecondary institutions, and represent differences in the extensive versus intensive margins of treatment with tuition subsidies.

These estimates suggest that Pell Grants \textit{issued during the 2010s alone}, which disbursed 349.8 billion dollars, increased voter turnout in the 2020 American elections by 1,819,000 votes and were responsible for the turnout of 1 out of every 87 American voters.\footnote{This is calculated by dividing 349.8 billion by 1,000 dollars and then multiplying by the effect size of 0.0052 votes per 1,000 dollars awarded.} Using conservative assumptions about the partisan composition of treated students\footnote{See the middle column of Table \ref{tab:envelope} and Section \ref{sec:policy} for a detailed explanation behind assumptions on partisanship. We note that we find similar results when using ANES 2020 survey data for the types of voters who are likely to receive financial aid and are similar to recipients of the Pell and Cal Grants. Specifically, we limit the sample to young voters aged 18 to 28 in 2020, who have family incomes below 100,000 dollars per year, and report tertiary education or postgraduate education. Approximately 77 percent of these voters nationwide report supporting Biden, compared to 17 percent who report supporting Trump and the remainder supporting minor party candidates. For voters in states outside of California, the comparable figures are 76 percent for Biden and 18 percent for Trump. In each case, the results are similar to our imputation of 80 percent of our sample favoring Biden and 15 percent favoring Trump, which we estimated using survey data from Firoozi (2023) and the partisan registration of in-sample students.}, the Pell Grant increased the Democratic Party's lead in the national popular vote by 1,182,000 votes (0.74 percentage points of total 2020 ballots) in the 2020 presidential election, with large enough effects to change the 2020 winner in Arizona, Georgia, Wisconsin, and the overall Electoral College.
\clearpage



\begin{table}[h!]\centering
    \def\sym#1{\ifmmode^{#1}\else\(^{#1}\)\fi}
    \caption{Estimated Impacts of 2010-2019 Pell Grant Generosity on Political Participation\label{tab:pell_results}}
    \begin{tabular}{l*{6}{c}}
        \hline\hline
        Outcome&\multicolumn{1}{c}{(1)}&\multicolumn{1}{c}{(2)}&\multicolumn{1}{c}{(3)}&\multicolumn{1}{c}{(4)}&\multicolumn{1}{c}{(5)}&\multicolumn{1}{c}{(6)}\\
        \hline
        &&&&&&\\
        \multicolumn{7}{l}{\textit{A. Voter Registration}} \\ [0.5em]
        \input{\basepath/Research/college_politics/Tables/pell_registered.tex} [0.5em]

        \multicolumn{7}{l}{\textit{B. Voter Turnout}} \\ [0.5em]

        \input{\basepath/Research/college_politics/Tables/pell_g2020.tex} [0.5em]

        \input{\basepath/Research/college_politics/Tables/pell_evervoted.tex} [0.5em]

        \input{\basepath/Research/college_politics/Tables/pell_turnoutrate.tex} [0.5em]

        \multicolumn{7}{l}{\textit{C. Voter Turnout by Partisanship}} \\ [0.5em]

        \input{\basepath/Research/college_politics/Tables/pell_centerleftvote.tex} [0.5em]

        \input{\basepath/Research/college_politics/Tables/pell_centerrightvote.tex}

        &&&&&&\\

        \hline
        Bandwidth & \$10,000       &   \$10,000        &   \$10,000        &   \$10,000       &   \$5,000        &   \$5,000        \\
        Kernel&   Uniform      &   Uniform        &   Uniform        &   Uniform       &   Uniform       &   Uniform       \\
        Polynomial&   1       &   1        &   2        &   2       &   1       &   1       \\
        Controls&   No       &   Yes        &   No        &   Yes        &   No        &   Yes        \\
        Sample Size      &   2,561,537    &    2,561,537      &      2,561,537     &   2,561,537   &    1,279,637       &    1,279,637        \\
        \hline
    \end{tabular}
    \caption*{\footnotesize Note: \sym{+} \(p<0.1\), \sym{*} \(p<0.05\), \sym{**} \(p<0.01\). Heteroskedasticity robust standard errors in parentheses. Sample excludes 2011-2012 FAFSA filers, because the data for that cohort was not available. ``Ever Voted'' refers to the extensive margin of ever having participated in a general election in the academic year after a student filed a FAFSA. ``Voter Turnout'' refers to the share of all federal general elections between 2010 and 2020 in which a student voted after the academic year in which they filed a FAFSA. ``Center-left Turnout'' refers to the interaction between voter turnout and an indicator for whether the student was a registered Democrat or independent, following \cite{firoozi_effect_2023}.  ``Center-right Turnout'' refers to the interaction between voter turnout and an indicator for being registered with the Republican Party. Outcomes correspond to those in Table \ref{tab:results2}.}
\end{table}




\clearpage



\subsection*{F. Heterogeneity and Time Variation}\label{sec:ch3_heterogeneity}
    \addcontentsline{toc}{section}{F. Heterogeneity and Time Variation}

\subsection{Time Variation}\label{sec:timevariation}

In Table \ref{tab:results3}, we examine the longitudinal impact of the Cal Grant on voter turnout, standardizing election outcomes relative to the year in which a student would have received a Cal Grant. In Row 1, we show the impact on voter turnout in the first general election that takes place after a student would have started receiving a Cal Grant. Row 2 shows the same outcome for the next general election, taking place after another 2 years. Row 3 shows results for the extensive margin of casting a ballot in any federal general election taking place 4 or more years after the first general election after Cal Grant receipt.

Our preferred specification in Column 1 shows that receiving a Cal Grant increases voter turnout by 4.96 percentage points in the first general election after Cal Grant receipt. This result is between 2 and 6 percentage points and significant at a 95 percent confidence interval across most RD implementation choices. While there are positive point estimates of meaningfully large magnitudes for elections taking place later on, we lack the precision to detect them and can reject neither the null hypothesis that there is no impact, nor the null hypothesis that the effects are the same as in the first election after Cal Grant receipt.

To make further use of time variation, we display results for the 2020-2021 cohort, whose instruction was fully online due to the COVID-19 pandemic in the year of their Cal Grant receipt, and compare them to earlier cohorts in Table \ref{tab:covid1}. Panel A shows results for the earlier 2017-2018 to 2019-2020 cohorts, while Panel B shows results for the 2020-2021 COVID-19 cohort. We find that there are not obvious differences in the impact on voter registration in 2022, but there are large gaps in the impact on voter turnout in the 2020 general election. Students whose schooling took place in person in the year in which they received their Cal Grant have significant positive effects on 2020 election turnout between 7 and 12 percentage points, whereas the COVID-19 cohort has null effects between -4 and 3 percentage points. Although we lack the statistical power under most specifications to reject the null hypothesis that the results are the same, we nonetheless view this as suggestive evidence consistent with on-campus residence and other campus-based socialization activities as a key mechanism in increasing student political participation.

Finally, in Table \ref{tab:elexovertime} we estimate the impact of the Cal Grant and Pell Grant across the four different elections in which we have sufficiently large sample sizes to match that of our main specification from Table \ref{tab:mainresults} Panel A Column 2, using all students within the four most recent entering cohorts.  Table \ref{tab:elexovertime} Row 1 Column 4 uses the exact same specification and, therefore, has the exact same sample size. Table \ref{tab:elexovertime} Row 1 Columns 1 through 3 do the same for the elections in 2014, 2016, and 2018, again restricting the sample to students who have had the same amount of time elapsed (within the 4 most recent years) since treatment: 2011 to 2014 students, 2013 to 2016 students, and 2015 to 2018 students, respectively. This leads to different sample sizes across years because cohorts have different numbers of students. Table D.4 Row 2 Columns 1 through 4 repeat the same exercise with the Pell Grant, including only students treated within the most recent 4 years of any given election shown in the respective columns.

We find larger point estimates on the effect of subisidies in the 2020 election than previous elections, growing point estimates over time for the California-based Cal Grant, and a pattern of larger effects in presidential cycles for the federal Pell Grant. However, it is worth cautioning against reading too much into heterogeneous point estimates over different election cycles because our confidence intervals remain too wide in many cases to reject the null hypothesis that these effects are consistent over time. We formally test for heterogeneity in treatment effects over time by comparing student cohorts whose first presidential elections were 2016 and 2020 to remain consistent with the samples used in Columns 2 and 4 of Table \ref{tab:elexovertime}.

We test the heterogeneity of the Pell Grant's effects on voter turnout over time by estimating the causal effect of being *below* the income ceiling for the Zero EFC Pell Bonus (being just poor enough to qualify for bonus grant funding) and test the significance on the interaction term between the excluded instrument (being below the threshold) and being in one of the more recent cohorts whose outcome is turnout in the 2020 election. We also fully interact the indicator for the 2020 election cohorts with all other parts of the RD specification and vary the bandwidth used for inference, the order of a polynomial control for the running variable, and the inclusion of pretreatment covariate controls to mirror Table \ref{tab:results2}. Results are shown in Table \ref{tab:pell_ttest} for the Pell Grant. A positive coefficient on ``Below Pell Limit'' in this case reflects positive impacts of grant funding on voter turnout. A positive coefficient on the interaction term likewise indicates stronger effects on the 2020 election.

In Table \ref{tab:first_ttest}, we show that the ITT approach for testing the time-varying effects of the Cal Grant is valid because the first stage results are consistent over time, falling between 12.5 and 13.8 percentage points.\footnote{Although the point estimates are significant in two out of six specifcation the differences are small (less than a percentage point) in economic terms.} Unfortunately we cannot repeat this for the Pell Grant because we do not observe the first stage of a Zero EFC rating for the Pell Grant over time.

Next, we test the heterogeneity of the Cal Grant's effects on voter turnout over time by estimating the causal effect of being *above* the income ceiling for Cal Grant eligibility (being too rich to qualify for grants) and test the significance on the interaction term between the excluded instrument (being above the threshold) and being in one of the more recent cohorts whose outcome is turnout in the 2020 election. We also fully interact the indicator for the 2020 election cohorts with all other parts of the RD specification and vary the bandwidth used for inference, the order of a polynomial control for the running variable, and the inclusion of pretreatment covariate controls to mirror Table \ref{tab:results1}. Results are shown in Table \ref{tab:calgrant_ttest} for the Cal Grant. A negative coefficient on ``Above CG limi'' in this case reflects positive impacts of grants on voter turnout. A negative coefficient on the interaction term indicates stronger positive effects on the 2020 election.

As the results across each table illustrate there is no consistent evidence that the Pell Grant has stronger treatment effects in the more recent 2020 election. In three of the six specifications there is evidence that the Cal Grant has stronger effects in more recent elections and in the other three the coefficient on the interaction term is not significant at the 10 percent level. On balance the results suggest that the Pell Grants effects are consistent over time with mixed evidence on the suggestion that the Cal Grant’s impact has grown.

To the extent that estimated effect sizes are growing over time, for the Cal Grant in particular, we offer the following hypotheses in order of our assessment of their plausibility:
\begin{itemize}
    \item[A)] Before the 2017 cohort, the California Student Aid Commission advertised income thresholds ex ante.

    As we discuss in response to Point 4, CSAC previously disclosed and advertised the Cal Grants' thresholds for eligibility. This may have caused unobservable selection into sample in the earlier years, drawing in students into the treatment group who are less likely to engage in pro-social behavior with small personal returns like voting. This is consistent with our results for the Cal Grant being biased toward zero in the pre-2017 cohorts and is also consistent with the fact that we find no time heterogeneity in the Pell Grant estimates.

    \item[B)] The intensive margin of aid generosity has risen for the Cal Grant, but less so for the Pell Grant.

The maximum award sizes for the Cal Grant have risen as the UC and CSU systems' tuition hike moratoria expired in the mid-2010s and the sticker price of tuition rose. The maximum award size for the Pell Grant, which is offered to students below the Zero EFC threshold, last experienced a significant increase in 2009, which is at the very beginning of our sample time period.

Given that we find consistent per-dollar estimates of treatment effects across the programs, it is plausible that some of the rise in the point estimates of the Cal Grant's treatment effect over time is attributable to rising generosity. The fact that the Pell Grant's maximum award size has not changed whereas the Cal Grants has is also consistent with the pattern of results we show in response to Point 2 and in Table D.4. Still, we point to one cause for skepticism in this explanation; The generosity of the Cal Grant has risen by a smaller magnitude than the point estimate in the Cal Grant’s treatment effect in Table D.4. That said, we cannot observe net tuition at the threshold and it is possible that net tuition at the threshold has risen more quickly, matching the point estimates.

    \item[C)] The extensive margin of income thresholds for Cal Grant eligibility have risen over time, but the Pell Grant has not.

A family of four had a Cal Grant A income threshold of 80,200 dollars in 2010 and 102,500 dollars in 2019. The Zero EFC threshold for the Pell Grant has gone from 30,000 dollars in 2010 to 26,000 dollars in 2019. These changes mean that the sample of near threshold students has changed as the policies bind to different percentiles within the income distribution over time.

Since Cal Grant students are increasingly upper income over time, part of the growth in the point estimate of the Cal Grant's treatment effects may be explained if richer students have larger treatment effects. The change in the point estimate of the Cal Grant’s effects and the Pell Grant's constant effects across presidential and midterm elections may be  consistent with this explanation. However, we repeatedly test for and do not find heterogeneity by various measures of socioeconomic status in Appendix D, contradicting this explanation.

\item[D)] Random noise will cause some estimates to be larger than others.

Testing for heterogeneity by election involves splitting a larger sample into subsamples, which will have some above average and some below average effect sizes by construction. There is some risk of false positives when testing the null hypothesis of equality across groups. The risk of a false positive is higher when the groups that we are testing heterogeneity for are not pre-specified or we do not have ex ante reasons to believe there is heterogeneity. There is some risk that this is going on with the Cal Grant and that the effect sizes are actually constant over time, given that we can only reject the null hypothesis of equality at a 10 percent level in 3 out of 6 specifications.

\item[E)] The types of students who are compliers at the threshold may have changed over time.

It is possible that the types of students who choose to apply for and accept a Cal Grant have changed over time in such a way that the students who accept Cal Grants in later time periods have larger treatment effects. The reason why we are skeptical of this explanation is that we found little evidence of heterogeneity based on individual-level characteristics in Appendix D and there is not an obvious explanation for why the types of students who choose to apply for and accept the Cal Grant would change in such a way that would not be replicated at the Pell Grant Zero EFC notch.
\end{itemize}





\clearpage


\begin{table}[h!]\centering
    \def\sym#1{\ifmmode^{#1}\else\(^{#1}\)\fi}
    \caption{Estimated Impacts of 2010-2019 Cal Grant Receipt on Political Participation\label{tab:results2}}
    \begin{tabular}{l*{6}{c}}
        \hline\hline
        Outcome&\multicolumn{1}{c}{(1)}&\multicolumn{1}{c}{(2)}&\multicolumn{1}{c}{(3)}&\multicolumn{1}{c}{(4)}&\multicolumn{1}{c}{(5)}&\multicolumn{1}{c}{(6)}\\
        \hline
        &&&&&&\\
        \multicolumn{7}{l}{\textit{A. Voter Registration}} \\ [0.5em]
        \input{\basepath/Research/college_politics/Tables/fulltime_registered.tex} [0.5em]

        \multicolumn{7}{l}{\textit{B. Voter Turnout}} \\ [0.5em]

        \input{\basepath/Research/college_politics/Tables/fulltime_g2020.tex} [0.5em]

        \input{\basepath/Research/college_politics/Tables/fulltime_evervoted.tex} [0.5em]

        \input{\basepath/Research/college_politics/Tables/fulltime_turnoutrate.tex} [0.5em]

        \multicolumn{7}{l}{\textit{C. Voter Turnout by Partisanship}} \\ [0.5em]

        \input{\basepath/Research/college_politics/Tables/fulltime_centerleftvote.tex} [0.5em]

        \input{\basepath/Research/college_politics/Tables/fulltime_centerrightvote.tex}

        &&&&&&\\

        \hline
        Bandwidth & \$10,000       &   \$10,000        &   \$10,000        &   \$10,000       &   \$5,000        &   \$5,000        \\
        Kernel&   Uniform      &   Uniform        &   Uniform        &   Uniform       &   Uniform       &   Uniform       \\
        Polynomial&   1       &   1        &   2        &   2       &   1       &   1       \\
        Controls&   No       &   Yes        &   No        &   Yes        &   No        &   Yes        \\
        Sample Size      &   738,046     &    738,046      &      738,046      &   738,046    &   354,091       &   354,091        \\
        \hline
    \end{tabular}
    \caption*{\footnotesize Note: \sym{+} \(p<0.1\), \sym{*} \(p<0.05\), \sym{**} \(p<0.01\). Heteroskedasticity robust standard errors in parentheses. Sample excludes 2011-2012 FAFSA filers, because the data for that cohort was not available. ``Ever Voted'' refers to the extensive margin of ever having participated in a general election in the academic year after a student filed a FAFSA. ``Voter Turnout'' refers to the share of all federal general elections between 2010 and 2020 in which a student voted after the academic year in which they filed a FAFSA. ``Center-left Turnout'' refers to the interaction between voter turnout and an indicator for whether the student was a registered Democrat or independent, following \cite{firoozi_effect_2023}.  ``Center-right Turnout'' refers to the interaction between voter turnout and an indicator for being registered with the Republican Party.}
\end{table}

\clearpage


\begin{table}[h!]\centering
    \def\sym#1{\ifmmode^{#1}\else\(^{#1}\)\fi}
    \caption{Estimated Impacts of 2010-2019 Cal Grant Receipt on Voter Turnout\label{tab:results3}}
    \begin{tabular}{l*{6}{c}}
        \hline\hline
        Outcome&\multicolumn{1}{c}{(1)}&\multicolumn{1}{c}{(2)}&\multicolumn{1}{c}{(3)}&\multicolumn{1}{c}{(4)}&\multicolumn{1}{c}{(5)}&\multicolumn{1}{c}{(6)}\\
        \hline
        &&&&&&\\
        \input{\basepath/Research/college_politics/Tables/vote1.tex}
        &   [738,046]     &    [738,046]      &      [738,046]    &   [738,046]    &   [354,091]       &   [354,091]        \\[0.5em]
        \input{\basepath/Research/college_politics/Tables/vote2.tex}
        &   [651,212]     &    [651,212]      &      [651,212]    &   [651,212]    &   [312,477]       &   [312,477]        \\[0.5em]
        \input{\basepath/Research/college_politics/Tables/vote3plus.tex}
        &   [479,717]     &    [479,717]      &      [479,717]    &   [479,717]    &   [230,317]       &   [230,317]        \\

        &&&&&&\\

        \hline
        Bandwidth & \$10,000       &   \$10,000        &   \$10,000        &   \$10,000       &   \$5,000        &   \$5,000        \\
        Kernel&   Uniform      &   Uniform        &   Uniform        &   Uniform       &   Uniform       &   Uniform       \\
        Polynomial&   1       &   1        &   2        &   2       &   1       &   1       \\
        Controls&   No       &   Yes        &   No        &   Yes        &   No        &   Yes        \\
        \hline
    \end{tabular}
    \caption*{\footnotesize Note: \sym{+} \(p<0.1\), \sym{*} \(p<0.05\), \sym{**} \(p<0.01\). Heteroskedasticity robust standard errors in parentheses. Sample size in brackets. Estimates exclude 2011-2012 FAFSA filers, because the data for that cohort was not available. ``Voted 1st Chance'' refers to an indicator for participating in the first general election taking place after a student filed their FAFSA, between one and two years later. ``Voted 2nd Chance'' refers to an indicator for participating in the second general election taking place after a student filed their FAFSA, between three and four years later. ``Voted 3rd+ Chance'' refers to an indicator for ever participating in any general election taking place subsequent to the second general election after a student filed their FAFSA.}
\end{table}

\clearpage

\begin{table}[h!]\centering
    \def\sym#1{\ifmmode^{#1}\else\(^{#1}\)\fi}
    \caption{Estimated Impacts of Recent Cal Grant Receipt on Political Participation\label{tab:covid1}}
    \begin{tabular}{l*{6}{c}}
        \hline\hline
        Outcome&\multicolumn{1}{c}{(1)}&\multicolumn{1}{c}{(2)}&\multicolumn{1}{c}{(3)}&\multicolumn{1}{c}{(4)}&\multicolumn{1}{c}{(5)}&\multicolumn{1}{c}{(6)}\\
        \hline
        &&&&&&\\
        \multicolumn{7}{l}{\textit{A. Sample: Academic Years 2017-2018 to 2019-2020}} \\ [0.5em]
        \input{\basepath/Research/college_politics/Tables/registered.tex}
        &   [258,329]     &   [258,329]     &     [258,329]     &   [258,329]    &   [123,744]       &   [123,744] \\[0.5em]

        \input{\basepath/Research/college_politics/Tables/g2020.tex}
        &   [258,329]     &   [258,329]     &     [258,329]     &   [258,329]    &   [123,744]       &   [123,744] \\[0.5em]

        \multicolumn{7}{l}{\textit{B. Sample: Academic Year 2020-2021 (COVID-19 Cohort)}} \\ [0.5em]

        \input{\basepath/Research/college_politics/Tables/2020_registered.tex}
        &   [83,664]     &   [83,664]      &     [83,664]    &   [83,664]   &   [40,290]       &   [40,290] \\[0.5em]

        \input{\basepath/Research/college_politics/Tables/2020_g2020.tex}
        &   [83,664]     &   [83,664]      &     [83,664]    &   [83,664]   &   [40,290]       &   [40,290] \\

        &&&&&&\\

        \hline
        Bandwidth & \$10,000       &   \$10,000        &   \$10,000        &   \$10,000       &   \$5,000        &   \$5,000        \\
        Kernel&   Uniform      &   Uniform        &   Uniform        &   Uniform       &   Uniform       &   Uniform       \\
        Polynomial&   1       &   1        &   2        &   2       &   1       &   1       \\
        Controls&   No       &   Yes        &   No        &   Yes        &   No        &   Yes        \\
        \hline
    \end{tabular}
    \caption*{\footnotesize Note: \sym{+} \(p<0.1\), \sym{*} \(p<0.05\), \sym{**} \(p<0.01\). Heteroskedasticity robust standard errors in parentheses. Sample size in brackets.}
\end{table}

\clearpage


\begin{table}[h!]\centering
    \def\sym#1{\ifmmode^{#1}\else\(^{#1}\)\fi}
    \caption{Impacts of Cal Grants and Pell Grants by Election\label{tab:elexovertime}}
    \begin{tabular}{l*{4}{c}}
        \hline\hline
        &\multicolumn{1}{c}{(1)}&\multicolumn{1}{c}{(2)}&\multicolumn{1}{c}{(3)}&\multicolumn{1}{c}{(4)}\\
        Election Year&\multicolumn{1}{c}{2014}&\multicolumn{1}{c}{2016}&\multicolumn{1}{c}{2018}&\multicolumn{1}{c}{2020}\\
        \hline
        &&&&\\

        \input{\basepath/Research/college_politics/Tables/fulltime_elexcal.tex}
        Sample Size      &       253,196     &      330,899      &   337,420   &   258,329      \\[0.75em]


        \input{\basepath/Research/college_politics/Tables/fulltime_elexpell.tex}
        Sample Size      &     828,132     &     1,181,103     &   1,261,717    &   1,003,200        \\


        &&&&\\

        \hline
        Bandwidth & \$10,000       &   \$10,000        &   \$10,000        &   \$10,000            \\
        Kernel&   Uniform      &   Uniform        &   Uniform        &   Uniform            \\
        Polynomial&   1       &   1        &  1        &   1            \\
        Controls&   No       &   No        &   No        &   No               \\

        \hline
    \end{tabular}
    \caption*{\footnotesize Note: \sym{+} \(p<0.1\), \sym{*} \(p<0.05\), \sym{**} \(p<0.01\). Heteroskedasticity robust standard errors in parentheses. Each column reflects the results for students entering in the respective calendar year or any of the three preceding years, excluding missing data and sample restrictions. Each of these specifications is meant to mimic that of Table \ref{tab:mainresults} Panel A Column 2, which includes cohorts within four most recent years of treatment to any given election. This is intended to make results across elections more directly comparable without capturing differences between short-run and longer-run treatment effects.}
\end{table}


\clearpage

\begin{table}[h!]\centering
    \def\sym#1{\ifmmode^{#1}\else\(^{#1}\)\fi}
    \caption{Pell Grant Time Heterogeneity in Treatment Effects Test\label{tab:pell_ttest}}
    \begin{tabular}{l*{6}{c}}
        \hline\hline
        Outcome&\multicolumn{1}{c}{(1)}&\multicolumn{1}{c}{(2)}&\multicolumn{1}{c}{(3)}&\multicolumn{1}{c}{(4)}&\multicolumn{1}{c}{(5)}&\multicolumn{1}{c}{(6)}\\
        \hline
        &&&&&&\\
        \multicolumn{7}{l}{\textit{A. Presidential Voter Turnout (2016 if  pre-2017, 2020 if  post-2017)}} \\ [0.5em]
        \input{\basepath/Research/college_politics/Tables/pell_ttest.tex} [0.5em]


        &&&&&&\\

        \hline
        Bandwidth & \$10,000       &   \$10,000        &   \$10,000        &   \$10,000       &   \$5,000        &   \$5,000        \\
        Kernel&   Uniform      &   Uniform        &   Uniform        &   Uniform       &   Uniform       &   Uniform       \\
        Polynomial&   1       &   1        &   2        &   2       &   1       &   1       \\
        Controls&   No       &   Yes        &   No        &   Yes        &   No        &   Yes        \\
        Sample Size      &   2,184,303     &     2,184,303      &      2,184,303     &   2,184,303   &    2,184,303      &   2,184,303        \\
        \hline
    \end{tabular}
    \caption*{\footnotesize Note: \sym{+} \(p<0.1\), \sym{*} \(p<0.05\), \sym{**} \(p<0.01\). Heteroskedasticity robust standard errors in parentheses. ``Presidentiall Voter Turnout'' refers to a student's participation in the first post-treatment presidential election, which is the 2016 election for the 2013 to 2016 cohorts and 2020 election for the 2017 and later cohorts.}
\end{table}



\clearpage



\begin{table}[h!]\centering
    \def\sym#1{\ifmmode^{#1}\else\(^{#1}\)\fi}
    \caption{Cal Grant Time Heterogeneity in First Stage Test\label{tab:first_ttest}}
    \begin{tabular}{l*{6}{c}}
        \hline\hline
        Outcome&\multicolumn{1}{c}{(1)}&\multicolumn{1}{c}{(2)}&\multicolumn{1}{c}{(3)}&\multicolumn{1}{c}{(4)}&\multicolumn{1}{c}{(5)}&\multicolumn{1}{c}{(6)}\\
        \hline
        &&&&&&\\
        \multicolumn{7}{l}{\textit{A. First Stage: Share of Students who Received a Cal Grant}} \\ [0.5em]
        \input{\basepath/Research/college_politics/Tables/first_ttest.tex} [0.5em]


        &&&&&&\\

        \hline
        Bandwidth & \$10,000       &   \$10,000        &   \$10,000        &   \$10,000       &   \$5,000        &   \$5,000        \\
        Kernel&   Uniform      &   Uniform        &   Uniform        &   Uniform       &   Uniform       &   Uniform       \\
        Polynomial&   1       &   1        &   2        &   2       &   1       &   1       \\
        Controls&   No       &   Yes        &   No        &   Yes        &   No        &   Yes        \\
        Sample Size      &   589,228     &    589,228      &      589,228       &   589,228    &   589,228      &  589,228       \\
        \hline
    \end{tabular}
    \caption*{\footnotesize Note: \sym{+} \(p<0.1\), \sym{*} \(p<0.05\), \sym{**} \(p<0.01\). Heteroskedasticity robust standard errors in parentheses.}
\end{table}



\clearpage

\begin{table}[h!]\centering
    \def\sym#1{\ifmmode^{#1}\else\(^{#1}\)\fi}
    \caption{Cal Grant Time Heterogeneity in Treatment Effects Test\label{tab:calgrant_ttest}}
    \begin{tabular}{l*{6}{c}}
        \hline\hline
        Outcome&\multicolumn{1}{c}{(1)}&\multicolumn{1}{c}{(2)}&\multicolumn{1}{c}{(3)}&\multicolumn{1}{c}{(4)}&\multicolumn{1}{c}{(5)}&\multicolumn{1}{c}{(6)}\\
        \hline
        &&&&&&\\
        \multicolumn{7}{l}{\textit{A. Presidential Voter Turnout (2016 if  pre-2017, 2020 if  post-2017)}} \\ [0.5em]
        \input{\basepath/Research/college_politics/Tables/calgrant_ttest.tex} [0.5em]


        &&&&&&\\

        \hline
        Bandwidth & \$10,000       &   \$10,000        &   \$10,000        &   \$10,000       &   \$5,000        &   \$5,000        \\
        Kernel&   Uniform      &   Uniform        &   Uniform        &   Uniform       &   Uniform       &   Uniform       \\
        Polynomial&   1       &   1        &   2        &   2       &   1       &   1       \\
        Controls&   No       &   Yes        &   No        &   Yes        &   No        &   Yes        \\
        Sample Size      &  589,228      &    589,228      &     589,228       &   589,228     &   589,228       &   589,228         \\
        \hline
    \end{tabular}
    \caption*{\footnotesize Note: \sym{+} \(p<0.1\), \sym{*} \(p<0.05\), \sym{**} \(p<0.01\). Heteroskedasticity robust standard errors in parentheses. ``Presidentiall Voter Turnout'' refers to a student's participation in the first post-treatment presidential election, which is the 2016 election for the 2013 to 2016 cohorts and 2020 election for the 2017 and later cohorts.}
\end{table}



\clearpage


\subsection{Heterogeneity}\label{sec:HTE}

Having demonstrated the Cal Grant's strong overall impact on political participation, we pivot to heterogeneous treatment effects and intermediate outcomes. The estimated results of our heterogeneity analysis are displayed in Table \ref{tab:hte1}, which includes four panels that show heterogeneity by the racial or ethnic composition of a student's home ZIP code in Panel A, the socioeconomic composition of their ZIP code in Panel B, the political composition of their ZIP code in Panel C, and the high school GPA of GPA-eligible students for whom this data is available in Panel D. Each column of the table represents the results for one quartile of the distribution of the respective variable for which we are assessing heterogeneity. We use total general election voter turnout in all post-treatment elections between 2010 and 2020 as our outcome of interest and repeat our preferred specification from Row 3, Column 1 of Table \ref{tab:results1} to provide a common point of comparison.

We start with heterogeneity by race and ethnicity in Panel A, using the racial and ethnic composition of a student's home region to proxy for these characteristics. Specifically, we use L2 voter file data to calculate the racial and ethnic shares of registered California voters who filed a FAFSA between 2010-2011 to 2020-2021 and collapse this data on ZIP code of origin. Rows 1 through 4 show the estimated impact of the Cal Grant on 2017-2018 to 2019-2020 FAFSA filers by the probability of being of a non-European ethnicity, Hispanic, Asian, and Black, respectively. We see little heterogeneity in the treatment effects of Cal Grants across these dimensions and are unable to detect significant differences by quartile of racial or ethnic composition. This suggests that the impact of higher education spending on political participation is unlikely to be driven by a single racial or ethnic group.

In terms of heterogeneity by socioeconomic composition in Panel B, we do not identify any significant differences in the treatment effects of Cal Grants on political participation. We repeat our methods from Panel A and find that there are no obvious patterns across quartiles of mean ZIP code income or mean ZIP code asset levels. We acknowledge that two limits to this analysis are the reality that Cal Grant recipients must themselves have low assets and that we are examining impacts for students local to an income eligibility ceiling. Our interpretation is that the absence of a clear pattern by neighborhood SES nonetheless provides suggestive evidence that there is unlikely to be large SES heterogeneity at the individual level.

Panel C repeats these methods again to examine whether effects are concentrated among students who originate from politically different households or neighborhoods, using the political composition of a FAFSA filer's ZIP code of origin as a proxy. Specifically, we collapse the 2020 voter turnout rate and Republican to Democratic ratio at the ZIP code of origin level, defining these as ZIP Code Voter Turnout and ZIP Code Conservatism, respectively. We find no detectable heterogeneity along either of these dimensions, suggesting that the political climate of one's upbringing is relatively unimportant in determining the magnitude of Cal Grants' impact on future political participation.

Finally, we examine heterogeneity in the Cal Grant's impact on voter turnout by a student's high school GPA, restricting to the subset of students for whom this data is available. We note that Cal Grant receipt increases the share of students reporting a GPA across the policy threshold, as GPA verification is required to determine the eligibility of high school students who are first-time college applicants. Because of potential issues with selection into reporting GPA, mixed listing of high school and community college GPAs, and the minimum 3.0 GPA eligibility limit for Cal Grants, we restrict our sample in this row to students with GPAs above 3.0 for whom we can identify a high school of origin. Our findings suggest that Cal Grants have a stronger impact on political participation among students with the highest high school GPAs, with the largest effects observed among students with GPAs above 3.41 (roughly a B+ average).

\clearpage

\begin{table}[h!]\centering
    \def\sym#1{\ifmmode^{#1}\else\(^{#1}\)\fi}
    \caption{Estimated Impacts of 2017-2019 Cal Grant Receipt on General Election Turnout\label{tab:hte1}}
    \begin{tabular}{l*{4}{c}}
        \hline\hline
        Dimension of Heterogeneity&\multicolumn{1}{c}{(1)}&\multicolumn{1}{c}{(2)}&\multicolumn{1}{c}{(3)}&\multicolumn{1}{c}{(4)}\\
        \hline
        &&&&\\
        \multicolumn{5}{l}{\textit{A. Heterogeneity by Racial and Ethnic Composition of Home ZIP Code}} \\ [0.5em]
        \input{\basepath/Research/college_politics/Tables/shareminority.tex} [0.5em]
        \input{\basepath/Research/college_politics/Tables/sharehispanic.tex} [0.5em]
        \input{\basepath/Research/college_politics/Tables/shareasian.tex} [0.5em]
        \input{\basepath/Research/college_politics/Tables/shareblack.tex} [0.5em]
        \multicolumn{5}{l}{\textit{B. Heterogeneity by Socioeconomic Composition of Home ZIP Code}} \\ [0.5em]
        \input{\basepath/Research/college_politics/Tables/regionincome.tex} [0.5em]
        \input{\basepath/Research/college_politics/Tables/regionasset.tex} [0.5em]
        \multicolumn{5}{l}{\textit{C. Heterogeneity by Political Composition of Home ZIP Code}} \\ [0.5em]
        \input{\basepath/Research/college_politics/Tables/regionvote.tex} [0.5em]
        \input{\basepath/Research/college_politics/Tables/regioncon.tex} [0.5em]
        \multicolumn{5}{l}{\textit{D. Heterogeneity by High School GPA (Subset of Full Sample)}} \\ [0.5em]
        \input{\basepath/Research/college_politics/Tables/gpa_mod.tex}
        &&&&\\

        \hline
        Bandwidth & \$10,000       &   \$10,000        &   \$10,000        &   \$10,000             \\
        Kernel&   Uniform      &   Uniform        &   Uniform        &      Uniform       \\
        Polynomial&   1       &   1        &   1        &   1         \\
        Controls&   No       &   No        &   No        &   No              \\
        Quartile      &   1     &    2      &      3      &   4           \\        \hline
    \end{tabular}
    \caption*{\footnotesize Note: \sym{+} \(p<0.1\), \sym{*} \(p<0.05\), \sym{**} \(p<0.01\). Heteroskedasticity robust standard errors in parentheses. The outcome for all regressions in this table is ``General Election Turnout'', which refers to the share of all federal general elections between 2010 and 2020 in which a student voted after the academic year in which they filed a FAFSA. ``ZIP Code Voter Turnout'' refers to the 2020 voter turnout rate in a student's home ZIP code among FAFSA filers who originated from their home ZIP code. ``ZIP Code Conservatism'' refers to the share of FAFSA filers originating from a student's home ZIP code that registered with a major political party and joined the Republican Party.}
\end{table}

\clearpage






\subsection*{G. Mechanisms}\label{sec:appendixd}
    \addcontentsline{toc}{section}{G. Mechanisms}



    \begin{table}[h!]\centering
    \def\sym#1{\ifmmode^{#1}\else\(^{#1}\)\fi}
    \caption{Turnout Effects of 2017-2019 Cal Grants Conditional on Voter Registration\label{tab:conditional_turnout}}
    \begin{tabular}{l*{6}{c}}
        \hline\hline
        Outcome&\multicolumn{1}{c}{(1)}&\multicolumn{1}{c}{(2)}&\multicolumn{1}{c}{(3)}&\multicolumn{1}{c}{(4)}&\multicolumn{1}{c}{(5)}&\multicolumn{1}{c}{(6)}\\
        \hline
        &&&&&&\\


        \input{\basepath/Research/college_politics/Tables/conditional_turnoutrate.tex}

        &&&&&&\\

        \hline
        Bandwidth & \$10,000       &   \$10,000        &   \$10,000        &   \$10,000       &   \$5,000        &   \$5,000        \\
        Kernel&   Uniform      &   Uniform        &   Uniform        &   Uniform       &   Uniform       &   Uniform       \\
        Polynomial&   1       &   1        &   2        &   2       &   1       &   1       \\
        Controls&   No       &   Yes        &   No        &   Yes        &   No        &   Yes        \\
        Sample Size      &     177,356    &     177,356    &      177,356     &   177,356  &  85,147     &   85,147    \\
        \hline
    \end{tabular}
    \caption*{\footnotesize Note: \sym{+} \(p<0.1\), \sym{*} \(p<0.05\), \sym{**} \(p<0.01\). Heteroskedasticity robust standard errors in parentheses. ``Voter Turnout'' refers to the share of all federal general elections in which a student voted after the academic year in which they filed a FAFSA. Outcomes directly correspond to those in Table \ref{tab:results1}.}
\end{table}

\clearpage

\subsection{External Validity Across States and Peer Socialization}\label{sec:externalvalid}
In our view, the peer socialization mechanism has more to do with time spent living with other students and discussing politics on campus in casual settings rather than short-lived but intense direct action like protests. In essence, financial aid increases time spent on campus with peers who come from relatively more educated families than a students' counterfactual time spent in the low-wage labor force, living off-campus, or with their parents/guardians. We think that this type of socialization is likely to generalize to other states (albeit with different magnitudes) based on the descriptive fact that voters with higher levels of educational attainment are more likely to vote than their less educated counterparts in every American state. We also view this as the most probable form of socialization, because protests tend to be concentrated at elite private colleges with wealthy students, which does not match the low income students and primarily less-selective institutions attended by the types of students in our sample.\footnote{See ``Are Gaza Protests Happening Mostly at Elite Colleges?'' in \textit{The Washington Monthly}, which employs the same dataset we use in this Appendix.}

That said, we implement an exercise using publicly available data from Harvard's Crowd Counting Consortium, which maintains records on public protests taking place in the United States.\footnote{The data are available at https://github.com/nonviolent-action-lab/crowd-counting-consortium.} We limit the sample to protests taking place since October of 2023 at locations including the words “college” or “university” and restrict to protests that are flagged as focusing on foreign affairs issues related to Israel, the IDF, Gaza, Palestine, Hamas, etc.. There are a total of 4,824 protest events in the sample.

We find that 14.2 percent of college and university campuses with such protests (88 out of 620) were located in the state of California, which is roughly in line with the 12 to 18 percent of American college students who study at Californian colleges in any given year. A higher fraction of total arrests, 20.5 percent (730 out of 3,567), and total protest events, 22.7 percent (1,094 out of 4,024), took place in California but in both cases results are driven by a single elite college. UCLA alone accounts for 253 (35 percent of) arrests and Stanford alone accounts for 292 (27 percent of) protest events. These elite colleges sum to less than 2 percent of statewide college enrollment and are far more selective, with admissions rates of 11 percent or less, than the typical 4-year university or 2-year community college that near threshold in-sample students attend.

Californian colleges may be especially effective at socialization, but we caution against over-interpreting the results of this exercise. The political culture of a college that the typical in-sample students attend is distantly removed from the elite colleges that give California its reputation for activism. Typical Pell and Cal Grant students attend post-secondary institutions that are indistinguishable from similarly selective institutions in Republican-leaning states, which is why we expect our results to generalize (see Table \ref{tab:externalvalid}).\footnote{Precinct level data for college election results are retrieved manually from https://www.nytimes.com/interactive/2021/upshot/2020-election-map.html. College admission rates are retrieved manually from https://www.nytimes.com/interactive/2023/07/03/opinion/for-most-college-students-affirmative-action-was-not-enough.html.} Two thirds of California 4-year public university enrollment, for example, is in the California State University system at campuses like CSU Fullerton, SDSU, and CSU Northridge, which closely match the partisanship and degree of activism of the modal college campuses in other states. Close to half of Pell Grant students attend open-access 2-year community college programs in California, which have even less protest activity than CSU campuses.


\clearpage


\begin{table}[h!]\centering
    \def\sym#1{\ifmmode^{#1}\else\(^{#1}\)\fi}
    \caption{Characteristics of Colleges in California, Texas, Florida, and North Carolina\label{tab:externalvalid}}
    \begin{tabular}{l*{1}{c c c c c c}}
        \hline\hline
        College  &  Admit \%  & State  & Biden 2020 & Trump 2020 &  Protests & Arrests \\
        \hline
        &&&&&&\\
        Stanford&4\%&CA&94\%&4\%&292&13\\
        Duke&6\%&    NC&    90\%&    8\%    &9&    0\\
        UCLA&11\%&    CA&    93\%&6\%    &59&    253\\
        UC Berkeley&14\%&    CA&    92\%&5\%    &57&    13\\
        UNC&20\%&    NC&    86\%&12\%    &26&    6\\
        UT Austin&29\%&    TX&    83\%&14\%    &17&    136\\
        Univ. Florida&30\%&    FL&    74\%&24\%    &16&    9\\
        UCF&36\%&    FL&    74\%&24\%    &21&    0\\
        SDSU&38\%&    CA&    75\%&23\%    &7&    0\\
        NCSU Raleigh&48\%&    NC&    79\%&19\%    &3&    0\\
        CSUF&59\%&    CA&    65\%&32\%    &6&    0\\
        Texas A\&M&64\%&    TX&    66\%&31\%    &9&    0\\
        FIU&64\%&    FL&    69\%&29\%    &14&    0\\
        UT Dallas&87\%&    TX&    63\%&35\%    &15&    22\\
        CSUN&88\%&    CA&    71\%&26\%    &4&    0\\
        &&&&&&\\
        \hline
    \end{tabular}
    \caption*{\footnotesize Note: Each column shows different characteristics of colleges. Admission rates and voting totals from precinct results are pulled manually from the New York Times. Gaza protests and arrest totals are recorded manually and totaled from Harvard’s Crowd Counting Consortium.}
\end{table}


\clearpage


\begin{table}[h!]\centering
    \def\sym#1{\ifmmode^{#1}\else\(^{#1}\)\fi}
    \caption{Effects of 2017-2019 Cal Grants on Partisanship and Vote Choice\label{tab:partisanship}}
    \begin{tabular}{l*{6}{c}}
        \hline\hline
        Outcome&\multicolumn{1}{c}{(1)}&\multicolumn{1}{c}{(2)}&\multicolumn{1}{c}{(3)}&\multicolumn{1}{c}{(4)}&\multicolumn{1}{c}{(5)}&\multicolumn{1}{c}{(6)}\\
        \hline
        &&&&&&\\
            \multicolumn{7}{l}{\textit{A. Total Registration}} \\ [0.5em]

        \input{\basepath/Research/college_politics/Tables/registered.tex}  [0.5em]

    \multicolumn{7}{l}{\textit{B. Registration by Partisanship in 2022}} \\ [0.5em]

        \input{\basepath/Research/college_politics/Tables/gop.tex}

        \% of New Registrants    &[-9\%]&[-18\%]&[-17\%]&[-37\%]&[-33\%]&[-73\%]\\ [0.5em]

        \input{\basepath/Research/college_politics/Tables/notgop.tex}

        \% of New Registrants    &[109\%]&[118\%]&[117\%]&[137\%]&[133\%]&[173\%]\\[0.5em]


            \multicolumn{7}{l}{\textit{C. Imputed 2020 Vote Probabilities}} \\ [0.5em]

                    \input{\basepath/Research/college_politics/Tables/probdemvote.tex} [0.5em]

                    \input{\basepath/Research/college_politics/Tables/probgopvote.tex} [0.5em]

                    \input{\basepath/Research/college_politics/Tables/didntvote.tex}

        &&&&&&\\

        \hline
        Bandwidth & \$10,000       &   \$10,000        &   \$10,000        &   \$10,000       &   \$5,000        &   \$5,000        \\
        Kernel&   Uniform      &   Uniform        &   Uniform        &   Uniform       &   Uniform       &   Uniform       \\
        Polynomial&   1       &   1        &   2        &   2       &   1       &   1       \\
        Controls&   No       &   Yes        &   No        &   Yes        &   No        &   Yes        \\
        Sample Size      &   258,329     &    258,329      &      258,329      &   258,329    &   123,744       &   123,744        \\
        \hline
    \end{tabular}
    \caption*{\footnotesize Note: \sym{+} \(p<0.1\), \sym{*} \(p<0.05\), \sym{**} \(p<0.01\). Heteroskedasticity robust standard errors in parentheses. Outcomes directly correspond to those in Table \ref{tab:results1}. ``Biden Vote 2020'' and ``Trump Vote 2020'' refer to the probability that a person cast a ballot for the respective candidate in the 2020 general election using actual data on whether or not they voted and Cal Grant voters' probabilities of favoring Democrats and Republicans match those of a sample of young, college-educated independents in recent survey data of California college applicants that were linked to voter registration records \citep{firoozi_effect_2023}.}
\end{table}

\clearpage


\begin{table}[h!]\centering
    \def\sym#1{\ifmmode^{#1}\else\(^{#1}\)\fi}
    \caption{Effects of 2010-2019 Pell Grant Generosity on Partisanship and Vote Choice\label{tab:pell_partisanship}}
    \begin{tabular}{l*{6}{c}}
        \hline\hline
        Outcome&\multicolumn{1}{c}{(1)}&\multicolumn{1}{c}{(2)}&\multicolumn{1}{c}{(3)}&\multicolumn{1}{c}{(4)}&\multicolumn{1}{c}{(5)}&\multicolumn{1}{c}{(6)}\\
        \hline
        &&&&&&\\
        \multicolumn{7}{l}{\textit{A. Total Registration}} \\ [0.5em]

        \input{\basepath/Research/college_politics/Tables/pell_registered.tex}  [0.5em]

        \multicolumn{7}{l}{\textit{B. Registration by Partisanship in 2022}} \\ [0.5em]

        \input{\basepath/Research/college_politics/Tables/pell_gop.tex}

        \% of New Registrants    &[-6\%]&[-4\%]&[11\%]&[23\%]&[3\%]&[12\%]\\ [0.5em]

        \input{\basepath/Research/college_politics/Tables/pell_notgop.tex}

        \% of New Registrants    &[106\%]&[104\%]&[89\%]&[77\%]&[97\%]&[88\%]\\[0.5em]


        \multicolumn{7}{l}{\textit{C. Imputed 2020 Vote Probabilities}} \\ [0.5em]

        \input{\basepath/Research/college_politics/Tables/pell_probdemvote.tex} [0.5em]

        \input{\basepath/Research/college_politics/Tables/pell_probgopvote.tex} [0.5em]

        \input{\basepath/Research/college_politics/Tables/pell_didntvote.tex}

        &&&&&&\\

        \hline
        Bandwidth & \$10,000       &   \$10,000        &   \$10,000        &   \$10,000       &   \$5,000        &   \$5,000        \\
        Kernel&   Uniform      &   Uniform        &   Uniform        &   Uniform       &   Uniform       &   Uniform       \\
        Polynomial&   1       &   1        &   2        &   2       &   1       &   1       \\
        Controls&   No       &   Yes        &   No        &   Yes        &   No        &   Yes        \\
        Sample Size      &   258,329     &    258,329      &      258,329      &   258,329    &   123,744       &   123,744        \\
        \hline
    \end{tabular}
    \caption*{\footnotesize Note: \sym{+} \(p<0.1\), \sym{*} \(p<0.05\), \sym{**} \(p<0.01\). Heteroskedasticity robust standard errors in parentheses. Outcomes directly correspond to those in Table \ref{tab:results1}. ``Biden Vote 2020'' and ``Trump Vote 2020'' refer to the probability that a person cast a ballot for the respective candidate in the 2020 general election using actual data on whether or not they voted and Pell Grant voters' probabilities of favoring Democrats and Republicans match those of a sample of young, college-educated independents in recent survey data of California college applicants that were linked to voter registration records \citep{firoozi_effect_2023}.}
\end{table}

\clearpage


\subsection{Income Effects and Per-Dollar Estimates}\label{sec:income}

The terminology we use with income effects is somewhat imprecise. In a technical sense, higher effective incomes afforded by increasing the intensive or extensive margin of financial aid receipt may be expended in the form of greater consumption of on-campus housing. Financial aid is also relatively easy from an administrative standpoint to apply toward on-campus housing when students receive aid from multiple sources that already partly cover tuition and fees.

Although this may be thought of as “income effects” because it reflects that students have a high marginal propensity to consume financial aid income through housing, we refer to this in our manuscript as peer socialization because the impact of the aid is coming through the channel of consumption of a good/service (on-campus housing) that is inducing peer socialization.

There are two major reasons why we think that similar per-dollar effects between the Cal Grant and Pell Grant are inconsistent with a direct ``income effect'' mechanism. The first is that there should be a diminishing marginal return of income to political participation if income effects are a core mechanism. The second is that income effects should make students more Republican whereas socialization effects should make students more Democratic.

We show in Figures \ref{fig:partisanship_byincome} and \ref{fig:incomegradient} that income is positively associated with voter registration and turnout, but the gradient is largest toward the bottom of the income distribution and becomes flat around the middle of the distribution. While changes to family income could plausibly impact students’ voter turnout and registration, we do not believe that our results are explained by that mechanism. This is because we find similar per-dollar effects at the middle of the income distribution with a change to the extensive margin of financial aid (with our marginal Cal Grant students) as we do at the bottom of the income distribution with a change to the intensive margin of the generosity of financial aid (with our marginal Pell Grant students).

Similarly, our results on partisanship align with peer socialization rather than income effects. Students from higher income families are more likely to be Republicans later in life and theoretical work suggests income effects should shift students rightward \citep{romer_individual_1975,meltzer_rational_1981}. However, to the extent that we find any changes in partisanship, the policies we study shift students toward the left. This is consistent with both Pell and Cal Grant marginal students consuming this financial aid through greater leisure time and in-person time with other left-leaning peers rather than counterfactual time spent with their (on average) more conservative guardians.

\clearpage

\begin{figure}[h!]
    \centering{\includegraphics[width=\textwidth]
        {\basepath/Research/college_politics/Figures/partisanship_byincome.png}}\caption{\label{fig:partisanship_byincome}Party Registration by Family Income}
    \caption*{\footnotesize Note: This figure shows the share of students at each level of family income with different party registration statuses.}
\end{figure}


\clearpage


\begin{figure}[h!]
    \centering{\includegraphics[width=\textwidth]
        {\basepath/Research/college_politics/Figures/incomegradient.png}}\caption{\label{fig:incomegradient}Voter Turnout Rate by Family Income}
    \caption*{\footnotesize Note: This figure shows the average voter turnout rate in post-treatment general elections at each level of family income.}
\end{figure}


\clearpage





\end{document}
